\documentclass[11pt,a4paper]{article}

\usepackage[T1]{fontenc}
\usepackage[utf8]{inputenc}
\usepackage[british]{babel}
\usepackage{lmodern}
\usepackage{microtype}
\usepackage[a4paper,margin=1in]{geometry}
\usepackage{amsmath,amssymb,amsthm,mathtools}
\usepackage{enumitem}
\usepackage{booktabs,tabularx,array}
\usepackage{xcolor}
\usepackage{hyperref}
\hypersetup{
  colorlinks=true,
  linkcolor=blue!60!black,
  citecolor=blue!60!black,
  urlcolor=blue!60!black,
  hypertexnames=false,
  pdftitle={A Preference for Traceable Agency: An Axiomatization of Expected Utility plus Mutual Information},
  pdfauthor={Daniele Condorelli},
  pdfsubject={A behavioural characterization of expected utility plus mutual information}
}
\setlength{\emergencystretch}{2em}
\interfootnotelinepenalty=10000

\newtheorem{theorem}{Theorem}
\newtheorem{proposition}[theorem]{Proposition}
\newtheorem{lemma}[theorem]{Lemma}
\newtheorem{corollary}[theorem]{Corollary}
\theoremstyle{definition}
\newtheorem{definition}[theorem]{Definition}
\newtheorem{remark}[theorem]{Remark}

\newcommand{\D}{\Delta}
\newcommand{\E}{\mathbb E}
\newcommand{\I}{I}
\newcommand{\Sh}{H}
\newcommand{\supp}{\operatorname{supp}}
\newcommand{\Id}{\operatorname{Id}}
\newcommand{\Rnum}{\mathbb R}
\newcommand{\R}{\mathbb R}
\newcommand{\KL}{D_{\mathrm{KL}}}
\newcommand{\1}{\mathbf 1}

\title{A Preference for Traceable Agency\\[0.35em]
\large An Axiomatization of Expected Utility plus Mutual Information\thanks{Email:
\href{mailto:d.condorelli@gmail.com}{d.condorelli@gmail.com}.  All parts of this manuscript have been written with assistance
from a multi-model LLM workflow; the author personally contributed very little to the writing of the proofs
beyond encouragement.  A Lean~4 verification of the main theorem is available at
\url{https://github.com/dcond/TraceableAgency}; it uses the equivalent finite-branch form of the recordwise
sure-thing principle.  All claims, errors, and omissions remain the sole responsibility of the author.}}
\author{Daniele Condorelli\\University of Warwick}
\date{ \today}

\begin{document}
\maketitle

\begin{abstract}
\noindent
An action can matter twice: through the payoff it produces and the visible trace it leaves.  A channel maps
actions into consequences; preferences rank action lotteries in each channel.  Data-processing conditions and a
recordwise sure-thing principle, with standard axioms, characterize expected utility plus $\lambda$ times the
mutual information between action and consequence.  The theorem identifies $\lambda$, the price of trace in
payoff units; expected utility is the case $\lambda=0$.  The criterion reads the coupling of act and
consequence, not consequences alone; randomisation can be strictly optimal.  When each action leaves a distinct
trace, the optimal lottery has multinomial-logit probabilities; when traces overlap, substitution depends on the
similarity of their consequence distributions and independence of irrelevant alternatives can fail.
\end{abstract}

%\tableofcontents
\clearpage

% ===============================================================
\section{Introduction}
% ===============================================================


\begin{quote}
\emph{Shower upon him every earthly blessing \ldots\ He would even risk his cakes and would deliberately desire the
most fatal rubbish, the most uneconomical absurdity \ldots\ simply in order to prove to himself---as though that were
so necessary---that men still are men and not the keys of a piano.\ \ldots\ If you say that all this, too, can be
calculated and tabulated \ldots\ then man would purposely go mad in order to be rid of reason and gain his point!}

\hfill --- Fyodor Dostoevsky, \emph{Notes from the Underground}, Part I, Chapter VIII \cite{Dostoevsky1864}
\end{quote}



Dostoevsky's Underground Man faces a choice problem.  The cakes are material payoff, and he will risk them so
that his conduct can attest to being his own, not the mechanical implication of a table.  Yet the proof he seeks
is self-defeating.  An act fixed in advance cannot provide it, and deliberate rebellion, too, can be
“calculated and tabulated.”  So he would “purposely go mad in order to be rid of reason”---not because
madness is what he values, but because it is his last imagined way of preventing the table from closing before he
acts.  This paper takes the evidentiary idea literally.  In Dostoevsky's
example, conduct is observed directly: the act supplies its own evidence.  We consider the general case, in which
an act is observed only through a possibly noisy payoff-bearing outcome and an explicit visible record.

For visible consequences to carry information about an act, two conditions are needed.  There must be uncertainty
beforehand about which act will occur, and the consequence must discriminate among the
possible acts.  Either condition can fail on its own.  A lottery over acts leaves no trace when every act generates
the same distribution of visible consequences; a consequence that perfectly reveals the act conveys
nothing new when the act was already known.  We call the intrinsic preference for the two together a preference for \emph{traceable agency}.  The paper axiomatizes it jointly with ordinary
preferences over payoff risk.

Let $A$ be a finite set of actions, $O$ a finite set of payoff-bearing
outcomes, $R$ a finite set of explicit visible records, and $K:A\to\Delta(O\times R)$ the channel relating actions to
visible consequences.  An action lottery $q\in\Delta(A)$ gives the distribution of the realised action.  The primitive is a family
$\{\succeq_K\}_K$, with one binary relation $\succeq_K$ on $\Delta(A)$ for each fixed channel $K$.
The statement $q\succeq_Kp$ means that $q$ is weakly preferred to $p$ in environment $K$.  Eight axioms are
imposed on this family.  To compare lotteries governed
by different channels, the channels are placed side by side as blocks of one larger environment; duplicating a
channel or adding an unused block does not alter the comparison at issue.  Two processing axioms say that the agent never gains when her record is garbled or her action variable is processed noisily, and a benchmark comparison rules out indifference to trace. A recordwise sure-thing principle for records that arrive in stages supplies
the cardinal structure.  No audience appears in the model, and the
taste is content-neutral: it prices how much the consequence reveals, not what it reveals.

The main result is an if-and-only-if characterization.  A family $\{\succeq_K\}_K$ satisfies the eight axioms if and
only if there are a nonconstant payoff utility $u:O\to\mathbb R$ and $\lambda>0$ such that, for every finite joint
channel $K$ and every $q,p\in\Delta(A)$,
\[
 q\succeq_Kp
 \quad\Longleftrightarrow\quad
 \E_{qK}[u(O)]+\lambda I_{qK}(A;O,R)
 \ge
 \E_{pK}[u(O)]+\lambda I_{pK}(A;O,R),
\]
where $I_{qK}(A;O,R)$ is the mutual information between the action and its visible consequence.  We call such
families \emph{trace-tempered}.  The utility $u$ is unique up to positive affine transformations, $\lambda$ is
measured in payoff units per unit of information, and the same pair prices comparisons across channels, embedded
as blocks of one larger environment.  A sign trichotomy completes the picture: reversing the axioms that carry the
taste for trace characterizes $\lambda<0$---a preference for deniability---while weakening them to indifference
yields expected utility, $\lambda=0$. 



The information term has a simple interpretation.  Observing $(O,R)$ induces a posterior belief about the realised
action; traceability is the expected reduction in that uncertainty, measured by Shannon entropy $\Sh$.  If visible
consequences are independent of actions, the value is zero.  If the action is known in advance because
$q$ is degenerate, the value is also zero.  If each action in the support of $q$ produces a distinct visible
consequence, the value is $\Sh(q)$---Dostoevsky's case, in which conduct is observed directly.\footnote{If the
visible-consequence supports of all channel rows are pairwise disjoint, the unique optimal lottery has
multinomial-logit weights.  Overlap can make independence of irrelevant alternatives fail; see
Remark~\ref{rem:logit}.}  Thus the criterion is
not a taste for randomisation as such.  Randomisation matters only because it creates uncertainty about agency for the
visible consequence to resolve.



Traceable agency is thus a taste for the coupling of act and consequence, not for consequences alone.  A
four-action channel turns this into a direct test.  Every action pays the same amount and writes one of two
symbols on the record.  Action $a_1$ always writes the first symbol and $a_2$ the second; $a_3$ and $a_4$ each
write either symbol with equal probability.  The choice is between a fair coin over $\{a_1,a_2\}$ and a fair
coin over $\{a_3,a_4\}$.  The two coins randomise identically and induce the same joint distribution of payoff
and record.  Under the first coin the record identifies the realised action; under the second it reveals
nothing.  Any criterion that reads only the distribution of consequences is indifferent between the coins; so is
a taste for randomisation as such.  Traceable agency predicts a strict preference for the first.  Because the
record affects neither payoffs nor any later decision, the preference cannot be
instrumental.\footnote{A second treatment completes the design: repeat
the choice in a channel that writes a fair record independently of the realised action.  Every representation
then assigns both flips the same value, whatever $u$ and $\lambda$, so a strict preference there rejects the
model.  Section~\ref{sec:implications} develops the test and what it excludes.}  The comparison identifies a
value of act--consequence dependence; by itself it does not distinguish trace from control or identify the
psychological source of that value.\footnote{The information term admits a control interpretation. Read backward: it measures trace. Read forward, it measures how strongly the action controls the distribution of visible consequences. We return to this in Section~\ref{sec:interpretation}.}



This comparison has not been run.  The nearest evidence comes from two experiments.  Bartling, Fehr, and Herz elicit each principal's point of indifference between retaining and delegating the right to choose a project and the effort devoted to it.  At that point, the delegation lottery has a certainty equivalent 16.7\% higher on average than the lottery produced by the principal's own choices, when both are valued outside the delegation setting \cite{BartlingFehrHerz2014}.  This is consistent with a preference for trace, since retaining control makes the principal's choices determine the payoff lottery, but it may instead reflect autonomy or some other value of deciding for oneself.  Mistry and Liljeholm, and later Norton and Liljeholm, ask subjects to select a room containing two slot machines and then gamble in that room.  Subjects are more likely to choose a high-divergence room over a zero-divergence room when they choose between its machines themselves than when a computer chooses for them, holding current expected payoffs and outcome entropies fixed \cite{MistryLiljeholm2016,NortonLiljeholm2020}.  With equal weight on the two machines, the divergence between their token distributions is the mutual information between the machine and the token.  This also admits a trace interpretation: in self-play, the token is evidence about a voluntarily selected machine.  But token values could change after the room was selected, allowing subjects in self-play to adapt their machine choices.  Thus trace is confounded with the value of deciding in the first design and with that of adapting in the second.  The test proposed above removes both confounds.

Section~\ref{sec:model} sets out the domain and the eight axioms.  Section~\ref{sec:theorem} states the representation theorem,
with a proof sketch, and derives uniqueness, the sign trichotomy, and the
independence of the axioms.
Section~\ref{sec:implications} presents a two-treatment test that can reject the model; derives deliberate randomisation, with multinomial logit as the boundary case when every action leaves its own mark; prices garblings of the record; measures $\lambda$ from choices; and characterises when one agent is more trace-seeking than another.
Section~\ref{sec:literature} reviews the literature, in particular rational inattention, where mutual information is a cost of attention rather than an object of
taste, and the axiomatic characterisations of information indices.
Section~\ref{sec:interpretation} discusses interpretation and scope.  The proofs are in the appendices:
the pure-trace characterisation in Appendix~A, the completion of Theorem~1 in Appendix~B, and the remaining
results in Appendix~C.  %Theorem~1 has also been verified in Lean~4.

% ===============================================================
\section{The model}\label{sec:model}
% ===============================================================

Fix a finite set $O$ of \emph{payoff-bearing outcomes}, with $|O|\ge2$.  Let $A$ be a nonempty finite set of
\emph{actions} and $R$ a nonempty finite set of \emph{explicit visible records}.  A joint channel is a stochastic
kernel
\[
  K:A\longrightarrow\D(O\times R).
\]
For each finite joint channel $K$, the primitive behavioural datum is a binary relation
\[
 \succeq_K\quad\text{on}\quad\D(A).
\]
The interpretation is that $q$ is a lottery over actions used by the agent and $q\succeq_Kp$ means that, in the fixed
environment $K$, $q$ is evaluated at least as highly as $p$.  Within that environment, the lottery is the
procedure being evaluated and the channel is fixed.\footnote{Why not a single
preference over lotteries on $A\times O\times R$?  Every joint law is $q\otimes K$ for some pair, so the domain
is rich enough.  But in any one choice situation the channel is fixed: the agent moves the marginal on $A$,
never the coupling.  A preference over all joint laws would rank couplings, and no choice reveals those
rankings.} % The block environments introduced below add tagged actions that select between component technologies.

\bigskip
Axioms are imposed on the family $\{\succeq_K\}_K$.
The first two axioms are standard, and Axiom~\textup{(A3)} is the usual nondegeneracy condition on material
outcomes.  Axioms~\textup{(A4)}, \textup{(A6)}, and \textup{(A7)} carry the preference for trace.
Axiom~\textup{(A5)} puts the separate within-channel rankings on a common scale, while
Axiom~\textup{(A8)} imposes consistency across stages; neither fixes the sign of the trace term.

\begin{description}[leftmargin=*,style=sameline]

\item[(A1) Weak order.]
For every finite joint channel $K$, $\succeq_K$ is complete and transitive on $\D(A)$.

\item[(A2) Continuity.]
For each fixed triple of alphabets $(A,O,R)$, the set
\[
 \{(K,q,p):q\succeq_Kp\}
 \subseteq \D(O\times R)^A\times\D(A)^2
\]
is closed.


\item[(A3) Material relevance.]
There are distinct outcomes $o^+,o^-\in O$ and a two-action channel $K$ with a one-point record set in which $a^+$ delivers
$o^+$ for sure and $a^-$ delivers $o^-$ for sure.  For this channel,
\[
 \delta_{a^+}\succ_K\delta_{a^-}.
\]
The agent strictly ranks at least two outcomes when each is delivered for sure and the two actions leave the same record.

\item[(A4) Trace relevance.]
There is an outcome $o^\ast\in O$ and a four-action channel $K$ in which every action delivers $o^\ast$ for
sure.  Actions $a_1$ and $a_2$ leave distinct records $r_1$ and $r_2$, respectively, while each of $a_3$ and
$a_4$ leaves either record with equal probability.  For this channel,
\[
 \frac12(\delta_{a_1}+\delta_{a_2})
 \succ_K
 \frac12(\delta_{a_3}+\delta_{a_4}).
\]
Both are fair lotteries and produce the same distribution of visible consequences. The axiom requires a strict
preference for the first, whose record identifies the realised action, over the second, whose record does not.

\end{description}

\paragraph{Comparison environments.}
Some axioms compare lotteries carried by different channels.  Each relation $\succeq_K$ ranks lotteries within
one channel, so such comparisons need a domain in which two channels appear in a single preference statement.  A
larger channel that contains both as labelled blocks provides it.

For
$K:A\to\D(O\times R)$ and $L:B\to\D(O\times R_L)$, let $K\sqcup L$ be the block-diagonal joint channel---the
\emph{block environment}---with action set
$(A\times\{0\})\cup(B\times\{1\})$, payoff set $O$, and record set
$(R\times\{0\})\cup(R_L\times\{1\})$: each block operates its own channel, and all cross-block probabilities
are zero.  For $q\in\D(A)$ and $p\in\D(B)$,
let $q^0$ and $p^1$ denote their
block-supported copies, defined on the tagged action set by
\[
\begin{aligned}
 q^0(a,0)&:=q(a), & q^0(b,1)&:=0,\\
 p^1(a,0)&:=0,    & p^1(b,1)&:=p(b),
\end{aligned}
\qquad(a\in A,\ b\in B).
\]

Under \(q^0\), the visible consequence is generated by \(K\) after the action is drawn according to \(q\); under \(p^1\), it is generated by \(L\) after the action is drawn according to \(p\). So a choice between them is
also a choice between the two channels, made inside a single environment.  %When the blocks have the same action alphabet, $q^0$ and $q^1$ are the two tagged copies of the same lottery.  
Write
\[
 (q,K)\succeq(p,L)
 \quad\Longleftrightarrow\quad
 q^0\succeq_{K\sqcup L}p^1.
\]
Each instance of this shorthand is a single comparison inside the two-block environment $K\sqcup L$, which is
built from the two channels; $\succ$ and $\sim$ denote the asymmetric and symmetric parts of this shorthand
relation.  The pair $(q,K)$
has no meaning in isolation.\footnote{Formally, $(K\sqcup L)(o,(r,0)\mid(a,0))=K(o,r\mid a)$ and
$(K\sqcup L)(o,(r_L,1)\mid(b,1))=L(o,r_L\mid b)$, with every other entry zero.  In the binary case, with
$A=\{a_1,a_2\}$ and $B=\{b_1,b_2\}$, the block environment has the four tagged actions $(a_1,0)$, $(a_2,0)$,
$(b_1,1)$, $(b_2,1)$; a lottery $q=(q(a_1),q(a_2))$ becomes $q^0=(q(a_1),q(a_2);0,0)$ and $p=(p(b_1),p(b_2))$
becomes $p^1=(0,0;p(b_1),p(b_2))$, the semicolon separating the blocks.  Then $(q,K)\succeq(p,L)$ abbreviates
$q^0\succeq_{K\sqcup L}p^1$: one ordinary comparison of two lotteries in this one four-action environment.}  The construction extends to any finite family
$\{K_i:A_i\to\D(O\times R_i)\}_{i\in I}$: write $\bigsqcup_{i\in I}K_i$ for the analogous block environment and
$q_i^i$ for the lottery that assigns $q_i(a)$ to $(a,i)$ and zero to every other block.

\begin{description}[leftmargin=*,style=sameline]

\item[(A5) Block-comparison coherence.]
\emph{Duplication:} for every $K:A\to\D(O\times R)$ and all $q,p\in\D(A)$,
\[
 q\succeq_Kp
 \quad\Longleftrightarrow\quad
 (q,K)\succeq(p,K).
\]
\emph{Irrelevant blocks:} for every finite block environment $\bigsqcup_{k\in I}K_k$, every ordered pair of
distinct blocks $i,j\in I$, and all $q_i\in\D(A_i)$ and $q_j\in\D(A_j)$,
\[
 q_i^i\succeq_{\bigsqcup_{k\in I}K_k}q_j^j
 \quad\Longleftrightarrow\quad
 (q_i,K_i)\succeq(q_j,K_j).
\]

Comparing $q$ and $p$ across two tagged copies of $K$ is just comparing them in $K$; adding blocks not involved
in a comparison leaves it unchanged.

\end{description}

\paragraph{Processing.}
The next two axioms use two ways of transforming an environment.  A \emph{record processor} is a stochastic
kernel $T:O\times R\to\D(O\times R')$ that garbles the record without changing the payoff:
$T(o',r'\mid o,r)=0$ whenever $o'\neq o$.  The garbled channel is simply the matrix product
$KT:A\to\D(O\times R')$,
\[
 (KT)(o',r'\mid a)
 :=\sum_{o,r}K(o,r\mid a)T(o',r'\mid o,r).
\]

An \emph{action processor} is a stochastic kernel $S:A\to\D(B)$ that replaces the action variable $A$ by a
processed action variable $B$.  Given $q\in\D(A)$, it induces the lottery $qS\in\D(B)$,
\[
 (qS)(b):=\sum_aq(a)S(b\mid a).
\]
A \emph{Bayesian pushforward} of $K$ by $S$ at $q$ is any channel $\widehat K:B\to\D(O\times R)$
satisfying\footnote{If $(qS)(b)=0$, the
equation leaves that row unrestricted; every statement below applies to every consistent completion, making
unreached processed actions irrelevant.}
\begin{equation}
 (qS)(b)\,\widehat K(o,r\mid b)
 =\sum_a q(a)S(b\mid a)K(o,r\mid a)
 \quad\text{for every }(b,o,r).
 \label{eq:action-processing}
\end{equation}
In the processed environment, the alternatives are the noisy names $b$, while the visible consequences remain
$(O,R)$.  Choosing $qS$ under $\widehat K$ replaces dependence on the original action with dependence on its
noisy name, without changing the marginal distribution of visible consequences.

\begin{description}[leftmargin=*,style=sameline]

\item[(A6) Record data processing.]
For every $K:A\to\D(O\times R)$, every record processor $T$, and every $q\in\D(A)$,
\[
 (q,K)\succeq(q,KT).
\]
The agent never gains when her record is garbled.

\item[(A7) Action data processing.]
For every $K:A\to\D(O\times R)$, every action processor $S:A\to\D(B)$, every $q\in\D(A)$, and every Bayesian
pushforward $\widehat K$ of $K$ by $S$ at $q$,
\[
 (q,K)\succeq(qS,\widehat K).
\]
Nor does she gain when the action variable is processed, holding the distribution of visible consequences fixed.

\end{description}

\paragraph{Compounding.}
For $k\ge1$, let $P:A\to\D(\{1,\ldots,k\})$ be a first-stage record channel and let
$(K_1,\ldots,K_k)$ be an ordered list of continuation channels $K_i:A\to\D(O\times R)$.  Their
\emph{compound environment} is
\begin{equation}
 \bigl(P\triangleright(K_1,\ldots,K_k)\bigr)(o,(i,r)\mid a)
 :=P(i\mid a)K_i(o,r\mid a),
 \qquad i=1,\ldots,k.
 \label{eq:compound}
\end{equation}
In words, one action is realised and the environment responds in two stages.  First, the action $a$ generates only
an interim record $i$ through $P$; no payoff is delivered at this stage.  Record $i$ selects the continuation
channel $K_i$, through which the same action generates the payoff $o$ and a further record $r$.  The agent does not
act again.  The interim record remains visible, so the final record is $(i,r)$.  Under $q$ and $P$, record $i$ is
\emph{reached} if $\sum_aq(a)P(i\mid a)>0$.  For every reached record, let $q_i$ denote the posterior distribution
of the action.\footnote{Thus $q_i(a):=q(a)P(i\mid a)/\sum_{a'}q(a')P(i\mid a')$.}

\begin{description}[leftmargin=*,style=sameline]

\item[(A8) Recordwise sure-thing principle.]
For every $q\in\D(A)$, every binary first-stage record channel $P:A\to\D(\{1,2\})$ for which record $1$ is
reached, and every three continuation channels $K,L,M:A\to\D(O\times R)$,
\begin{equation}
 (q_1,K)\succeq(q_1,L)
 \quad\Longleftrightarrow\quad
 \bigl(q,P\triangleright(K,M)\bigr)
 \succeq
 \bigl(q,P\triangleright(L,M)\bigr).
 \label{eq:recordwise-sure-thing}
\end{equation}

Fix $q$ and a binary first-stage record channel $P$, and suppose record $1$ is reached, inducing posterior $q_1$.
The two compound environments share the continuation $M$ after record $2$ and differ only after record $1$, where
one continues with $K$ and the other with $L$.  The axiom says that the compounds are ranked exactly as the lottery
$q_1$ in channel $K$ and the same lottery in channel $L$ are ranked in their block environment.  The common
continuation on the other branch cannot affect the comparison.  This is a version of Savage's sure-thing principle
with the event given by an interim record \cite{Savage1954}.  Unlike the exogenous event in Savage's formulation,
reaching the record can be informative here: $P(1\mid a)$ may vary with the action.  This is why the conditional
comparison is evaluated at $q_1$.  If record $1$ is never reached, the axiom is silent, as on Savage's null
events.  When $P$ is an action-independent public coin, $q_1=q$, and
Axiom~\textup{(A8)} reduces to ordinary mixture independence with the coin retained in the record.

\end{description}

% ===============================================================
\section{The representation theorem}\label{sec:theorem}
% ===============================================================

For $q\in\D(A)$ and $K:A\to\D(O\times R)$, let
\[
 p_{qK}(o,r):=\sum_a q(a)K(o,r\mid a)
\]
be the induced distribution of visible consequences, and define
\begin{equation}
 \I_{qK}(A;O,R)
 :=\sum_{a,o,r}q(a)K(o,r\mid a)
 \log\frac{K(o,r\mid a)}{p_{qK}(o,r)}.
 \label{eq:mi}
\end{equation}
Summands with $q(a)K(o,r\mid a)=0$ are set to zero, the logarithm has any fixed base greater than one, and expectations below are
taken under the joint distribution $q(a)K(o,r\mid a)$.

\begin{theorem}\label{thm:main}
For a family $\{\succeq_K\}_K$ as described above, the following are equivalent.
\begin{enumerate}[label=\textup{(\roman*)},leftmargin=*]
\item The family satisfies Axioms~\textup{(A1)--(A8)}.
\item There are a nonconstant function $u:O\to\mathbb R$ and a number $\lambda>0$ such that, for every nonempty
finite $A,R$, every $K:A\to\D(O\times R)$, and every $q,p\in\D(A)$,
\begin{equation}
 q\succeq_Kp
 \quad\Longleftrightarrow\quad
 V(q,K)\ge V(p,K),
 \qquad
 V(q,K):=\E_{qK}[u(O)]+\lambda\I_{qK}(A;O,R).
 \label{eq:representation}
\end{equation}
\end{enumerate}
Moreover, under these equivalent conditions, block-supported cross-channel comparisons are represented on the same
scale: for every finite block environment $\bigsqcup_{k\in I}K_k$, every two distinct blocks $i,j\in I$, and every
$q_i\in\D(A_i)$ and $q_j\in\D(A_j)$,
\[
 q_i^i\succeq_{\bigsqcup_{k\in I}K_k}q_j^j
 \quad\Longleftrightarrow\quad
 V(q_i,K_i)\ge V(q_j,K_j).
\]
\end{theorem}

We call a family satisfying \textup{(A1)--(A8)} \emph{trace-tempered}, and likewise the criterion that
represents it.  Its central restriction is that payoff and trace enter separately.  The agent may value good
outcomes and may value evidence about her action, but she cannot place an extra value on a good payoff being the
revealing consequence.  Thus, if under the same lottery two channels give the same expected utility from payoffs
and induce the same distribution of posteriors over actions, the agent is indifferent between them, regardless of
how payoffs and records are paired (Corollary~\ref{cor:matching}, Appendix~\ref{app:auxiliary-results}).  This
separation is derived, not assumed.

\subsection{Proof sketch}
That the representation satisfies the axioms is a direct verification.  The converse has four steps.
Appendix~\ref{app:main-proof} gives the complete proof; Appendix~\ref{pt:sec:environment} supplies the pure-trace
result used in Step~3.

\emph{Step 1: reduction to terminal laws.}  Fix a full-support $q$.  If two channels induce the same joint
distribution of the realised payoff and the posterior over actions, each is a record processing of the other.
Block coherence and record processing therefore make them indifferent.  Action processing makes action
relabellings immaterial and, for an arbitrary $q$, permits null actions to be deleted.  Thus $(q,K)$ is identified
by its \emph{terminal law}, the joint distribution of the realised payoff and posterior
(Lemma~\ref{lf:lem:terminal-affine}\textup{(a)} and
Lemma~\ref{lf:lem:bookkeeping}\textup{(a)--(b)}).

\emph{Step 2: cardinalisation.}  Let a public coin select between two channels.  Because the coin is independent
of the action, the posterior in either branch remains $q$.  The recordwise sure-thing principle therefore gives
mixture independence over terminal laws.  With continuity, the mixture-space theorem gives an affine
representation for each $q$.  All cardinal content enters here; every other axiom is an ordinal condition,
imposed environment by environment (Lemma~\ref{lf:lem:terminal-affine}\textup{(b)}).

\emph{Step 3: payoff and trace.}  On singleton action sets there is no trace, so the affine representation reduces
to expected utility.  Material relevance makes its index $u$ nonconstant, while action processing and block
coherence make the same payoff scale apply throughout.  At the payoff $o_\ast$ that witnesses trace relevance,
the payoff term is constant.  Appendix~\ref{pt:sec:environment} shows that the remaining order is represented by
mutual information.  The finite-branch extension of Axiom~\textup{(A8)} (Lemma~\ref{pt:lem:finite-branch-extension}) supplies the chain rule; action processing supplies the corresponding
grouping property; and block coherence aligns comparisons across priors.  Together these
properties give Faddeev's recursion and hence Shannon entropy \cite{Faddeev1956}.  Trace relevance fixes the
positive sign, and the common block scale determines a single coefficient $\lambda$ relative to $u$
(Lemmas~\ref{lf:lem:material-scale} and~\ref{lf:lem:trace-scale}).

\emph{Step 4: assembly.}  The finite-branch extension of Axiom~\textup{(A8)} also determines how payoff and trace combine.
Continue to use the payoff $o_\ast$ that witnesses trace relevance.  Begin with a compound that pays $o_\ast$
after every first-stage record.  If record $y$ occurs with probability $m(y)$, changing only the payoff after $y$
from $o_\ast$ to $o$ changes value by
$
 m(y)\bigl[u(o)-u(o_\ast)\bigr].
$
This baseline compound has value
\[
 u(o_\ast)+\lambda\I(A;Y).
\]
Summing across branches gives the representation when the payoff is determined by the first-stage record.  For
an arbitrary channel, announce $(o,r)$ as an interim record and then deliver $o$.  This sequential channel and
the original channel are record processings of one another and hence are indifferent.  The representation
follows (Lemmas~\ref{lf:lem:branch-increment} and~\ref{lf:lem:deterministic-assembly}).

\subsection{Uniqueness}

A real-valued index $W$ on all finite lottery--channel pairs is a \emph{common-scale representation} if, for every channel
$K:A\to\D(O\times R)$ and every $q,p\in\D(A)$,
\[
 q\succeq_Kp
 \quad\Longleftrightarrow\quad
 W(q,K)\ge W(p,K),
\]
and, for every finite block environment $\bigsqcup_{k\in I}K_k$, every two distinct blocks $i,j\in I$, and all
$q_i\in\D(A_i)$ and $q_j\in\D(A_j)$,
\[
 q_i^i\succeq_{\bigsqcup_{k\in I}K_k}q_j^j
 \quad\Longleftrightarrow\quad
 W(q_i,K_i)\ge W(q_j,K_j).
\]

Axiom~\textup{(A8)} is ordinal: it says that a common continuation on the other branch cannot affect a conditional
comparison.  By Lemma~\ref{pt:lem:finite-branch-extension}, this is equivalent---given the other structural
axioms---to the requirement that branchwise improvements cannot be reversed by finite compounding.  Branch
additivity is the corresponding requirement on numerical differences.  A common-scale representation $W$
is \emph{branch-additive} if, for every $q\in\D(A)$, every $P:A\to\D(\{1,\ldots,k\})$, and every
two ordered lists of continuation channels $(K_1,\ldots,K_k)$ and $(L_1,\ldots,L_k)$, with
$K_i,L_i:A\to\D(O\times R)$, writing
$m(i):=\sum_aq(a)P(i\mid a)$,
\begin{equation}
 W\bigl(q,P\triangleright(K_1,\ldots,K_k)\bigr)
 -W\bigl(q,P\triangleright(L_1,\ldots,L_k)\bigr)
 =\sum_{i:m(i)>0}m(i)
 \bigl[W(q_i,K_i)-W(q_i,L_i)\bigr].
 \label{eq:branch-additive-index}
\end{equation}
A continuation change worth $x$ after record $i$ therefore changes ex-ante value by $m(i)x$, and changes on
different branches add. 

\begin{proposition}[Uniqueness]\label{prop:uniqueness}
Suppose \textup{(A1)--(A8)} hold, and fix a representation $V$, with parameters $(u,\lambda)$, as in
Theorem~\ref{thm:main}.  For any real-valued index $W$ on all finite lottery--channel~pairs:
\begin{enumerate}[label=\textup{(\roman*)},leftmargin=*]
\item $W$ is a common-scale representation if and only if $W=\varphi\circ V$ for a single strictly increasing
$\varphi:[\underline u,\infty)\to\mathbb R$, where $\underline u:=\min_{o\in O}u(o)$.
\item A common-scale representation $W$ is branch-additive if and only if
$W=\alpha V+\beta$ for some $\alpha>0$ and $\beta\in\mathbb R$.
\end{enumerate}
\end{proposition}

The proof is in Appendix~\ref{app:auxiliary-results}.

Every trace-tempered representation is a common-scale representation and, by the expected-utility and
mutual-information chain rules, branch-additive.  Part~\textup{(ii)} therefore implies that if $(u,\lambda)$ and
$(\widetilde u,\widetilde\lambda)$ are two such representations, then
\[
 \widetilde u=\alpha u+\beta,
 \qquad
 \widetilde\lambda=\alpha\lambda
\]
for some $\alpha>0$ and $\beta\in\mathbb R$.  Once two non-indifferent payoff outcomes are assigned numerical
utilities, $\lambda$ is unique for the chosen logarithm base.  In the pure-trace model of
Appendix~\ref{pt:sec:environment}, multiplying the index by a positive constant has no behavioural content.  Here
the payoff normalisation fixes that scale, and $\lambda$ measures the trade-off between payoff and trace.

\subsection{Sign of the trace weight}

Uniqueness concerns scale; a complementary observation concerns sign.  Axioms~\textup{(A4)},
\textup{(A6)}, and \textup{(A7)} orient the trace component.  Write $(\mathrm{A}4^-)$,
$(\mathrm{A}6^-)$, and $(\mathrm{A}7^-)$ for the clauses obtained by reversing their displayed comparisons,
and write $(\mathrm{A}4^0)$, $(\mathrm{A}6^0)$, and $(\mathrm{A}7^0)$ for the corresponding indifference clauses.

\begin{corollary}[Sign trichotomy]\label{cor:signduality}
Suppose Axioms~\textup{(A1)--(A3)}, \textup{(A5)}, and
\textup{(A8)} hold.  Then:
\begin{enumerate}[label=\textup{(\roman*)},leftmargin=*]
\item under Axioms~\textup{(A4)}, \textup{(A6)}, and \textup{(A7)}, the family has the representation in
Theorem~\ref{thm:main}, with $\lambda>0$;
\item under $(\mathrm{A}4^-)$, $(\mathrm{A}6^-)$, and $(\mathrm{A}7^-)$, it has the same representation with
$\lambda<0$;
\item under $(\mathrm{A}4^0)$, $(\mathrm{A}6^0)$, and $(\mathrm{A}7^0)$, it is represented by expected utility,
equivalently by the same formula with $\lambda=0$.
\end{enumerate}
Thus the trace component has exactly the positive, negative, and neutral orientations.
\end{corollary}

The proof is in Appendix~\ref{app:auxiliary-results}.

The directions of the three orientation clauses cannot be chosen independently.  For example, under the
corollary's background axioms and any orientation of Axiom~\textup{(A7)}, Axiom~\textup{(A4)} is inconsistent
with $(\mathrm{A}6^-)$.  Restrict the two lotteries in Axiom~\textup{(A4)} to their supports and relabel the
actions.  The second channel is then a total garbling of the first.  Axiom~$(\mathrm{A}6^-)$ weakly prefers the
second, whereas Axiom~\textup{(A4)} strictly prefers the first.

The negative case is quiet in unconstrained choice over lotteries.  With $\lambda<0$, every payoff-maximising
pure action is optimal.  More generally, $q$ is optimal in a fixed channel if and only if it is supported on
payoff-maximising actions and $\I_{qK}(A;O,R)=0$.  To reveal deniability in optimal choice, one must therefore
constrain the lottery or compare channels.  Under an imposed fair lottery over two actions, the agent is willing
to give up as much as $|\lambda|\log2$ of expected utility from payoffs to replace full revelation with silence.

\subsection{Independence of the axioms}

\begin{proposition}[Independence of the axioms]\label{prop:independence}
For each $i\in\{1,\ldots,8\}$ there is a family satisfying the other seven axioms and violating
Axiom~\textup{(A$i$)}.
\end{proposition}

The eight families are constructed in Appendix~\ref{app:auxiliary-results}.  Three mark the substantive frontier of the
system.  A conditional-entropy penalty satisfies every axiom except record data processing and values the removal of
extraneous record noise.  A Gini posterior-reduction criterion satisfies every axiom except action data processing.
The satiation criterion $\E[u(O)]+\min\{I(A;O,R),c\}$, for $c>0$, satisfies every axiom except the recordwise
sure-thing principle and ceases to value trace once $c$ units of information are secured.

% ===============================================================
\section{Choice implications}\label{sec:implications}
% ===============================================================

This section works out what an agent with the representation of Theorem~\ref{thm:main} does: which finite designs
falsify the model directly (\S\ref{sec:tracetest}), when she deliberately randomises (\S\ref{sec:randomisation}),
what she will pay for changes to the record technology (\S\ref{sec:record-price}), and how that price can be
measured and compared across agents (\S\ref{sec:measurement}).
Throughout, $\lambda$ is measured in payoff units per unit of information in the fixed base.  Proofs for this
section, together with full statements of the supporting results that the text presents in summary form, are
collected in Appendix~\ref{app:choice-proofs}; only one-line verifications are kept here.

\subsection{The trace test}\label{sec:tracetest}

The criterion $V$ reads the joint distribution of action and visible consequence, rather than the consequence
distribution or action lottery alone.  The following two-treatment design separates these objects. In both treatments, the lottery must be knowingly chosen and implemented by the subject, since the primitive ranks her own randomisations.\footnote{A direct value for information about which action was realised is observationally equivalent on this domain. It can be separated by assigning the same lottery externally: the information structure remains, but the subject no longer chooses the lottery.}

\begin{proposition}[The trace test]\label{prop:signature}
Four actions pay the same outcome $o_0$.  Two of them leave distinct marks: $a_1$ always produces record $r_1$,
and $a_2$ always produces the distinct record $r_2$.  The other two are indistinguishable: $a_3$ and $a_4$ each
produce $r_1$ or $r_2$ with equal probability.  Call this channel $K$, and let $q'$ be the fair coin over
$\{a_1,a_2\}$ and $q''$ the fair coin over $\{a_3,a_4\}$.
\begin{enumerate}[label=\textup{(\roman*)},leftmargin=*]
\item The two coins induce the same distribution of visible consequences, $p_{q'K}=p_{q''K}$; yet
\[
 V(q',K)-V(q'',K)=\lambda\log2>0,
 \qquad\text{so}\qquad
 q'\succ_Kq''.
\]
\item Let $K^\circ$ be the channel in which every action pays $o_0$ and produces $r_1$ or $r_2$ with equal
probability.  Then $q'\sim_{K^\circ}q''$.
\end{enumerate}
\end{proposition}

The prediction is strict preference for the revealing coin in the first treatment and exact indifference in the
second.  Part \textup{(i)} rules out any criterion that depends only on the distribution of visible consequences.
The two coins induce the same distribution over payoff--record pairs, so any such criterion is indifferent between them.

Part \textup{(i)} restricts only how a theory reads consequences; a theory that also reads the lottery can
match it.  Existing theories of deliberate randomisation are of this kind: in the hedging model
of Cerreia-Vioglio, Dillenberger, Ortoleva, and Riella \cite{CerreiaVioglioDillenbergerOrtolevaRiella2019}, and
in the preference for randomisation documented by Agranov and Ortoleva \cite{AgranovOrtoleva2017}, the value of
mixing depends only on the induced distribution and the lottery.  Call an index $W$ \emph{channel-blind} if it
reads only these two objects: there is a function $F$, fixed across channels, such that
\[
 W(q,K)=F\bigl(p_{qK},q\bigr).
\]

Part \textup{(ii)}  rules these out.  Nothing a channel-blind index reads changes between the treatments: in
all four lottery--channel combinations the distribution of visible consequences is uniform on
$\{(o_0,r_1),(o_0,r_2)\}$, and the lotteries are the same two lotteries.  A channel-blind index therefore
separates the coins in both treatments or in neither, so none reproduces the pattern of
Proposition~\ref{prop:signature}: strict preference under $K$ and indifference under
$K^\circ$.\footnote{If $W$ does not depend on how the actions are labelled---an entropy bonus, say---the first
treatment is enough: the permutation that swaps $a_1$ with $a_3$ and $a_2$ with $a_4$ turns $q'$ into $q''$ and
leaves $p_{qK}$ unchanged, so $W(q',K)=W(q'',K)$.  The second treatment is needed for label-sensitive indices.}

The first treatment also identifies the orientation of taste: it selects the revealing coin when $\lambda>0$,
the noisy coin when $\lambda<0$, and indifference when $\lambda=0$.  A strict preference between the coins in
the second treatment rejects all three cases of Corollary~\ref{cor:signduality}.


\subsection{Deliberate randomisation}\label{sec:randomisation}

Randomisation gives the record something to reveal.  A pure action leaves no
trace: if only one action can be realised, its consequence cannot reduce uncertainty about it.  The
agent may therefore give up expected utility to make her action uncertain, so that its consequence can reveal
it.  Whether she does so depends on how different the actions' visible consequences are.  Write
$K_a:=K(\cdot\mid a)$ for the distribution of visible consequences under action $a$ and
$\bar u(a):=\E_{K_a}[u(O)]$ for its expected utility.  The Kullback--Leibler divergence
$\KL(P\,\|\,Q)$ measures how far $P$ departs from $Q$.\footnote{For distributions $P$ and $Q$,
$\KL(P\,\|\,Q):=\sum_{o,r}P(o,r)\log\bigl(P(o,r)/Q(o,r)\bigr)$, with zero-mass summands set to zero.  The
divergence is $+\infty$ if $P$ assigns positive probability to a consequence and $Q$ assigns zero.}

Consider two actions $a$ and $b$, with $K_a\ne K_b$, and suppose that $a$ has expected-utility advantage
$\Delta:=\bar u(a)-\bar u(b)\ge0$.  Proposition~\ref{prop:binary} in
Appendix~\ref{app:choice-proofs} gives the exact rule:
\[
 \delta_a\text{ is optimal}
 \quad\Longleftrightarrow\quad
 \Delta\ge\lambda\,\KL(K_b\,\|\,K_a).
\]
A small shift of probability from $a$ to $b$ costs $\Delta$ per unit.  At the pure choice, its marginal trace
value is $\lambda\KL(K_b\,\|\,K_a)$.  The order in the divergence matters because the question is whether to
introduce $b$ into a lottery concentrated on $a$.  If $b$ can produce a consequence that $a$ cannot, the
divergence is infinite---the first unit of probability on $b$ is unboundedly informative---and $b$ receives
positive probability whatever its finite payoff disadvantage.

Below the threshold the optimal lottery is unique and mixes the two actions.  The inferior action receives more
probability as $\lambda$ rises and less as $\Delta$ rises.  At equal payoffs the chosen mixture maximises the
information carried by the channel; in information-theoretic terms, it is the capacity-achieving input. % Thus randomisation is instrumental to trace.  An uninformative channel supplies no reason to mix, and a taste for mixing as such is channel-blind.

With more than two actions, each action is compared with the distribution of visible consequences generated by
the lottery as a whole.

\begin{proposition}[Optimality condition]\label{prop:optimal-choice}
Let $\lambda>0$.  A lottery $q^*$ is optimal if and only if there is a constant $c$ such that
\[
 \bar u(a)+\lambda\KL\!\left(K_a\,\middle\|\,p_{q^*K}\right)=c
 \quad\text{for }a\in\supp(q^*),
 \qquad
 \le c\quad\text{for }a\notin\supp(q^*).
\]
\end{proposition}

An action is credited with its expected utility and with how far its visible consequences depart from those
generated on average.  At an optimum this total is the same for every action in use; an unused action cannot
offer more.

The condition has a sharp support implication.  Every visible consequence that can arise under some action
occurs with positive probability under any optimal lottery.  Hence, if each action can produce some consequence
that no other action can produce, every optimal lottery uses every action.  No finite payoff disadvantage can
eliminate an action with such a private consequence (Proposition~\ref{cor:support-pure},
Appendix~\ref{app:choice-proofs}).

The optimal lottery need not be unique, but every optimal lottery induces the same distribution of 
consequences.  If two actions have identical rows, only their combined probability matters: reallocating a fixed
total between them changes nothing.  Splitting an action into duplicate copies therefore leaves the set of
optimal aggregate choices unchanged and permits any division of the original mass among the copies.  A model
that gives each action its own weight, as multinomial logit does, instead raises the copies' combined share
(Proposition~\ref{prop:optimal-marginal}, Appendix~\ref{app:choice-proofs}).

\begin{remark}[Separated and overlapping traces]\label{rem:logit}
If every action leaves its own mark, the row supports are pairwise disjoint, the visible consequence identifies
the action, and $\I_{qK}(A;O,R)=\Sh(q)$.  The unique optimum is then multinomial logit,
\[
 q^*(a)
 =\frac{\exp\bigl(\bar u(a)/\lambda\bigr)}
 {\sum_{a'}\exp\bigl(\bar u(a')/\lambda\bigr)},
\]
with $1/\lambda$ as the payoff sensitivity when information is measured in natural logarithms.  Changing the
logarithm base rescales $\lambda$.  At the other extreme, suppose every action produces the same distribution of payoffs and records. The consequence then reveals nothing about the action, and all actions have the same expected utility. Hence the agent is indifferent among all lotteries.  With partial overlap, the ratio
of two actions' probabilities can depend on the rest of the menu because adding an action changes the average
consequence distribution.  Pairwise-disjoint row supports therefore imply independence of irrelevant
alternatives, while overlap permits, but does not require, its failure.  Remark~\ref{rem:overlap-iia-example} in
Appendix~\ref{app:choice-proofs} gives a three-action example in which the ratio falls from $1$ to $3/4$.
The same logit formula appears in rational inattention for the opposite reason: there the prior is fixed and the
channel is chosen; here the channel is fixed and the distribution of actions is chosen
\cite{MatejkaMcKay2015}.
\end{remark}

\subsection{The price of being on the record}\label{sec:record-price}

The representation puts different record technologies on a common scale.  Processing a record leaves payoffs
unchanged, so its price depends only on the information about the action that it removes.

\begin{proposition}[Price of a record garbling]\label{prop:garbling-price}
Let $K:A\to\D(O\times R)$, let $T$ be a record processor, and write
$L:=KT$.  Fix $q\in\D(A)$ and let $(A,O,R,R')$ carry the joint law $q(a)K(o,r\mid a)T(o,r'\mid o,r)$.  Then
\[
 V(q,K)-V(q,L)=\lambda\,\I\!\left(A;R\mid O,R'\right)\ge0,
\]
with equality if and only if, after $(O,R')$ has been observed, the original record $R$ contains no further
information about the action.
\end{proposition}

Rewriting the record costs $\lambda$ times the information about the action that the original record carried
beyond the rewritten one.  A garbling is free exactly when it combines only records that support the same
posterior about the action.  The agent does not pay for detail as such, only for detail that distinguishes among
her actions.\footnote{The same calculation applies to added information.  If $K^{+S}$ appends a signal $S$
drawn conditional on $(a,o,r)$, then
$
 V(q,K^{+S})-V(q,K)=\lambda\,\I(A;S\mid O,R).
$
Under the deniability case of Corollary~\ref{cor:signduality}, the sign is reversed.}

One application concerns the mixing device itself.  Suppose an action-independent coin $Y$ selects $K$ with
probability $t$ and $L$ otherwise.  If the coin is recorded, the resulting environment is worth
$tV(q,K)+(1-t)V(q,L)$: the agent has no taste for the draw itself.  If the coin is hidden, the resulting channel
is the rowwise mixture $M_t:=tK+(1-t)L$, and
\[
 tV(q,K)+(1-t)V(q,L)-V(q,M_t)
 =\lambda\,\I(A;Y\mid O,R)\ge0.
\]
The coin says nothing about the action by itself.  Hiding it is costly only when knowing which channel operated
would sharpen the inference about the action after the other visible consequences have been observed.  The
premium is for recording which channel operated, not for drawing the coin.  A theory that values mixing as such
has no reason to distinguish a recorded draw from an unrecorded one.  Corollary~\ref{cor:recorded-coin} in
Appendix~\ref{app:choice-proofs} gives the full statement.

Uniform rankings do not identify garbling.  Suppose two channels $K$ and $L$ generate the same payoff distribution after every action, so they
differ only in their records, and every action lottery is worth at least as much under $K$ as under $L$.  This says
exactly that $K$ carries at least as much mutual information at every lottery---that $K$ is more capable
\cite{KornerMarton1977}---but not that $L$ is a garbling of $K$: two constant-payoff binary channels can be ordered
by information at every prior although no record processor maps one into the other
(Proposition~\ref{prop:blackwell}, Appendix~\ref{app:choice-proofs}).  Uniform preference reveals capability, not
censorship.

\subsection{Measuring the trace weight}\label{sec:measurement}

The coefficient $\lambda$ is the price of one unit of information in expected-utility units.  Any indifference
between two environments with different amounts of trace reveals it.  Write $E_{qK}:=\E_{qK}[u(O)]$.

\begin{proposition}[Single information price]\label{prop:single-price}
If $(q,K)\sim(p,L)$ and $\I_{qK}\ne\I_{pL}$, then
\[
 \frac{E_{pL}-E_{qK}}{\I_{qK}-\I_{pL}}=\lambda,
\]
and $\lambda$ is the only number for which this identity holds at every such indifference.
\end{proposition}

In the information--payoff plane, indifference curves are parallel lines of slope $-\lambda$.  The slope does
not depend on the prior, the payoff level, the action alphabet, or the amount of information already present.
This constancy is the observable restriction that Axiom~\textup{(A8)} adds to the rest of the system.  A taste
for trace that satiates instead has a measured slope that falls to zero beyond the satiation point.

Consider a lottery $q$ with $\Sh(q)>0$ and two channels with the same constant payoff, one that reveals the
realised action and one that is silent.  Revelation is worth $\lambda\Sh(q)$.  Thus compensating the silent
channel until the agent is indifferent measures her willingness to pay for trace.  One implementation dilutes
both channels by the same public side branch, reached with probability $\varepsilon$.  Normalise the two
outcomes in Axiom~\textup{(A3)} so that $u(o^+)=1$ and $u(o^-)=0$.  On the revealing side the branch pays $o^-$
for sure; on the silent side it pays $o^+$ with probability $\beta$ and $o^-$ otherwise.  If $\beta$ produces
indifference, then
\[
 \lambda=\frac{\varepsilon}{1-\varepsilon}\frac{\beta}{\Sh(q)}.
\]
The common baseline payoff's utility cancels, so it need not be known.  Proposition~\ref{prop:wtp} and
Remark~\ref{rem:elicitation} in Appendix~\ref{app:choice-proofs} give the construction, boundary cases, and
controls.

Varying the prior tests the model rather than merely estimating $\lambda$: the premia at $q=(1/2,1/2)$ and
$q=(1/4,3/4)$ must stand in the ratio
\[
 \frac{\log 2}{\frac14\log4+\frac34\log\frac43}
 \approx1.23.
\]
A quadratic posterior-reduction criterion gives $4/3$.  Agreement across the two designs overidentifies
$\lambda$; disagreement rejects the Shannon specification even if each design separately admits an information
price.  The switching threshold in \S\ref{sec:randomisation} supplies a third estimate of $\lambda$, from
optimal choices rather than stated indifferences: the payoff advantage at which the agent abandons mixing equals
$\lambda$ times a known divergence.  It must return the same number.

Within a fixed channel, $\lambda$ governs the demand for trace.  Holding payoff utility fixed, raising $\lambda$
leads the agent to choose a lottery that carries weakly more information about her action and yields weakly less
expected utility from payoffs.  As $\lambda$ approaches zero, payoff is maximised first and trace breaks ties; as
$\lambda$ becomes large, trace is maximised first and payoff breaks ties
(Proposition~\ref{prop:trace-demand} and Corollary~\ref{cor:lambda-limits},
Appendix~\ref{app:choice-proofs}).

The coefficient also permits comparisons across agents.  Say agent 2 is more trace-seeking than agent 1 if she
weakly prefers an environment to a silent benchmark whenever agent 1 does.  This ordering holds exactly when
their representations can use the same payoff utility, after a common normalisation, with
$\lambda_2\ge\lambda_1$.  Agents with different payoff tastes are not ordered
(Proposition~\ref{prop:comparative}, Appendix~\ref{app:choice-proofs}; compare
\cite{GhirardatoMarinacci2002}).


% ===============================================================
\section{Related literature}\label{sec:literature}
% ===============================================================

This paper intersects several literatures that differ in what is ranked, which way information flows, and
whether an information functional is assumed or derived.  In most models of information, a signal tells the
agent about an exogenous state; here a visible consequence carries information about the agent's realised
action.

Kreps \cite{Kreps1979} and Dekel, Lipman, and Rustichini \cite{DekelLipmanRustichini2001} use preferences over
menus to represent flexibility and uncertainty about future tastes; Pattanaik and Xu \cite{PattanaikXu1990} rank
opportunity sets by the freedom they provide.  Gul and Pesendorfer \cite{GulPesendorfer2001} use the same domain
to separate temptation, commitment, and self-control, while Kreps and Porteus \cite{KrepsPorteus1978} use
temporal lotteries to study preferences over the timing of uncertainty resolution.  In this paper, a choice
problem fixes the feasible actions, the decision rights, and the channel from actions to visible consequences.
The primitive preference ranks lotteries over actions and therefore reads the dependence between the realised
action and its visible consequence, rather than the size of a menu, the value of commitment, or the timing of
resolution.

Grant, Kajii, and Polak \cite{GrantKajiiPolak1998} characterise an intrinsic preference for information
independent of its use in planning; Caplin and Leahy \cite{CaplinLeahy2001} model anticipatory feelings generated
by beliefs, while Ely, Frankel, and Kamenica \cite{ElyFrankelKamenica2015} study preferences over paths of beliefs
and the resulting value of suspense and surprise.  These papers show that information need not be instrumental,
but the relevant beliefs concern uncertainty faced by the agent, whereas the posterior here concerns her
realised action.  The axioms also restrict its value to expected entropy reduction and fix its rate of
substitution with material utility.

Cerreia-Vioglio, Dillenberger, Ortoleva, and Riella
\cite{CerreiaVioglioDillenbergerOrtolevaRiella2019} study deliberate randomisation and characterise a model in
which mixing hedges uncertainty about how to evaluate lotteries; Agranov and Ortoleva
\cite{AgranovOrtoleva2017} document stochastic choice under announced repetition and evidence that much of it is
deliberate.  Trace-tempered choice can also make a nondegenerate lottery optimal.
Proposition~\ref{prop:signature}, however, holds fixed the action lotteries and their distributions of visible
consequences while changing whether the consequence identifies the realised action; its second treatment removes
that dependence.  Together the two treatments rule out every channel-blind criterion that reads only the lottery
and its consequence distribution.  Randomisation is an implication of traceable agency, not its primitive
motive.

In signaling models following Spence \cite{Spence1973}, actions convey information about hidden characteristics.
Glazer and Konrad \cite{GlazerKonrad1996} study charitable giving as a signal of wealth, and Andreoni and Bernheim
\cite{AndreoniBernheim2009} study social image and audience effects.  B\'enabou and Tirole
\cite{BenabouTirole2006} allow concerns for social reputation or self-respect.  Their earlier paper
\cite{BenabouTirole2002} studies self-serving beliefs about ability, while Bernheim \cite{Bernheim1994} shows how
status concerns can induce agents with heterogeneous preferences to choose the same action, making behaviour less
informative about predispositions---an instrumental counterpart to the deniability case in
Corollary~\ref{cor:signduality}.  In these models, inference matters through esteem, sanction, reputation, or
self-regard.  Here no such response enters utility.  Beliefs serve only as the statistical yardstick for how
clearly the consequence identifies the act.

The proof also relates to axiomatic characterisations of entropy.  Shannon \cite{Shannon1948} characterised
entropy from continuity, monotonicity on uniform distributions, and a grouping recursion; Faddeev
\cite{Faddeev1956} reduced these to symmetry, continuity of the binary entropy, and the recursion itself.  The
proof uses the strong-additivity form of Faddeev's theorem given by Baez, Fritz, and Leinster
\cite{BaezFritzLeinster2011}; other formulations and extensions include Khinchin \cite{Khinchin1957}, R\'enyi
\cite{Renyi1961}, Acz\'el and Dar\'oczy \cite{AczelDaroczy1975}, and Csisz\'ar \cite{Csiszar2008}.  Mutual
information can also be obtained through relative entropy, which Baez and Fritz \cite{BaezFritz2014}
characterise through properties of Bayesian inference.

In those characterisations, a numerical information functional is primitive and structural conditions are
imposed directly on it; here the primitive is an ordinal family of weak orders over action lotteries.  The
information chain rule and Faddeev's recursion are consequences of behavioural axioms, not postulates on a given
function.  The theorem also calibrates the information term against material utility.  The result is not another
characterisation of entropy; it derives entropy reduction as the representation of a preference over trace.

Blackwell \cite{Blackwell1953} compares experiments by their value in all downstream decision problems; Gilboa
and Lehrer \cite{GilboaLehrer1991} characterise which functions on partitions can arise as values of information
in Bayesian decision problems, and Azrieli and Lehrer \cite{AzrieliLehrer2008} extend the analysis to stochastic
information structures.  Cabrales, Gossner, and Serrano \cite{CabralesGossnerSerrano2013} obtain entropy
reduction from the decisions of ruin-averse investors.  For Blackwell's decision-maker, information is an input:
she can use a finer signal or ignore it, so more information cannot hurt.  Here information is an output: the
record is evidence about the agent, and no decision follows it.  Uniform preference between two record
technologies identifies the more-capable order, not Blackwell garbling (Proposition~\ref{prop:blackwell},
Appendix~\ref{app:choice-proofs}).  Whether a finer record is desirable is therefore a matter of taste---positive
for the trace-seeker, negative under the deniability dual of Corollary~\ref{cor:signduality}, and zero under
expected utility.  Being better informed is a theorem; being better known is a preference.

The closest formal comparison is with rational inattention and information costs.  Sims \cite{Sims2003} imposes
a finite Shannon-capacity constraint on information flow, while Mat\v{e}jka and McKay
\cite{MatejkaMcKay2015} show that a cost proportional to mutual information yields a generalized
multinomial-logit rule.  Rational inattention fixes a prior over states and allows the agent to choose how much
to learn before acting, whereas here the channel from actions to consequences is fixed and the agent chooses a
distribution over actions.  Pairwise-disjoint row supports yield multinomial logit.  With overlapping rows,
logit can still obtain, but independence of irrelevant alternatives can fail, as Remark~\ref{rem:logit} shows.

Caplin, Dean, and Leahy \cite{CaplinDeanLeahy2022} characterise Shannon rational inattention through Invariance
under Compression and develop posterior-separable generalisations; Denti \cite{Denti2022} gives testable
conditions under which state-dependent stochastic choice is rationalised by posterior-separable information
costs.  Frankel and Kamenica \cite{FrankelKamenica2019} characterise measures of information that arise as the
value of news in a decision problem and show that the corresponding expected reduction in uncertainty equals
expected information.  In these papers, the information cost or the underlying optimisation problem supplies
the cardinal scale; the present primitive is ordinal, and the common scale for payoff and trace is derived.

Pomatto, Strack, and Tamuz \cite{PomattoStrackTamuz2023} take a cardinal, prior-independent cost functional on
Blackwell experiments and show that monotonicity, continuity, product additivity, and constant marginal cost
characterise weighted sums of Kullback--Leibler divergences between state-conditional signal distributions.
Their costs and the trace value studied here are not nested.  The trace component
$(q,K)\mapsto I_{qK}(A;O,R)$ depends on the prior and is not a prior-free additive cost of an experiment.  Full
revelation has finite trace value $\Sh(q)$; under every nonzero cost in their class it is excluded from the
finite-cost domain and would have infinite cost.  Record and action data processing are behavioural assumptions
about whether the agent values evidence of her action; reversing them, with the taste for revelation, yields a
preference for deniability, not an information-acquisition cost.  The information-cost literature prices
learning about the world; the present paper prices being learned about.

% PARKED 2026-08-11 (referee-verified experimental-evidence paragraphs; restore here if wanted, and
% reinstate the four commented bibitems plus a "(Section~\ref{sec:literature})" pointer in the intro's
% evidence paragraph):
%
% Two experimental literatures come close to the object priced here.  In Bartling, Fehr, and Herz, retaining
% control makes the principal's project and effort choices determine the payoff lottery, whereas delegation
% makes another person's choices determine it.  At the elicited point of indifference, the delegation lottery
% has a certainty equivalent 16.7 percent higher on average than the control lottery; the difference is positive
% for 83 percent of principals, and the finding replicates \cite{BartlingFehrHerz2014,BuffatPraxmarerSutter2023}.
% Similar premia arise when subjects bet on their own performance rather than another person's
% \cite{OwensGrossmanFackler2014}.  These findings admit a trace interpretation: people may value consequences
% generated by their own acts.  They do not identify the mutual-information term, however, because they change
% whose act generates the consequences rather than varying, at a fixed lottery over one person's acts, how
% informative the consequences are about the realised act.
%
% Norton and Liljeholm provide the closest statistical antecedent.  Subjects first choose a pair of slot
% machines and then make several machine choices from that pair; each machine produces a distribution over
% coloured tokens.  Subjects prefer pairs in which the two machines produce more distinct token distributions,
% holding expected payoff and the machines' outcome entropies fixed \cite{MistryLiljeholm2016,NortonLiljeholm2020}.
% Their measure of distinctness is Jensen--Shannon divergence, which, with the machines weighted equally, is the
% mutual information between the machine chosen and the token received.  Thus choice responds to the same
% statistic used here.  But subjects knew that token values could change after the pair was selected and before
% the machines were chosen, making this information useful for adapting the later choices.  Their estimate
% therefore combines any direct preference for mutual information with its option value; consistent with this
% reading, a related experiment that made the value changes explicit finds that a forward-looking
% expected-utility model fits choices better than a model assigning divergence direct utility \cite{Liljeholm2022}.
% The trace test of Section~\ref{sec:tracetest} removes both confounders: the same person acts throughout, and
% the information has no later use.

% ===============================================================
\section{Interpretation and scope}\label{sec:interpretation}
% ===============================================================

\paragraph{Control and empowerment.}
The information term reads in two directions.  Backward, it measures trace: how much the visible consequence
reveals about the realised action.  Forward, it measures control: how strongly the action shapes the
distribution of visible consequences.  Mutual information is symmetric, but the channel is
directional---$K$ maps actions into distributions over consequences---so the two readings describe the same
dependence.  No choice among lotteries separates them.  The control in question is informational: how
distinguishable actions are in their consequences, not the agent's power to obtain the consequences she
prefers.

The forward reading has an exact counterpart.  Suppose $\E[u(O)\mid a]=\bar u$ for every action.  Since
$\lambda>0$, maximising $V$ over $q$ gives
\[
 \max_{q\in\Delta(A)}V(q,K)=\bar u+\lambda C(K),
 \qquad
 C(K):=\max_{q\in\Delta(A)}I_{qK}(A;O,R),
\]
the Shannon capacity of the channel.  When the visible consequence is read as the agent's sensor state, $C(K)$
is the one-step empowerment index of Klyubin, Polani, and Nehaniv \cite{KlyubinPolaniNehaniv2005}; see also
\cite{SalgeGlackinPolani2014}.  Under this supposition, the trace-tempered agent solves exactly the problem
that defines one-step empowerment: her optimal lotteries are the capacity-achieving action distributions, and
the size of $\lambda$ plays no role in the ranking---the price is visible only against payoffs.  Empowerment
summarises the channel by the optimised value $C(K)$; the representation recovers the ranking beneath it,
$q\mapsto I_{qK}(A;O,R)$, from preferences and prices it against material utility.  Empowerment is thus the
trace part of indirect utility when expected payoffs are equal; once they differ across actions, the ranking
also reflects their material value, which the empowerment index does not encode.

\paragraph{Evidence, not authorship.}
The trace term prices evidence of agency, not authorship as such.  A pure action has zero trace value even when
its consequence names it perfectly: an act settled in advance leaves no uncertainty for the consequence to
resolve.  A coin the agent flips herself can do better.  The model does not say that surrendering a choice to
chance creates more agency; it says that an unsettled act can leave more evidence than a predetermined one.  The
lottery must be hers---knowingly chosen and implemented, as in the test of \S\ref{sec:tracetest}.  This is the
Underground Man's problem: the proof he wants cannot come from an act fixed in advance.

\paragraph{What the axioms preclude.}
A principal substantive exclusion is selective legibility.  Mutual information responds to an act's
probability and to the distinctiveness of its consequences, not to what the act is.  An agent who wants her generous
acts recognised and her mean acts hidden may therefore prefer an informative record in one channel and a garbling
in another.  Such preferences violate the global data-processing requirement in Axiom~\textup{(A6)} and fall
outside all three cases of the sign trichotomy.  They are not irrational; they are simply not represented here.
A model of selective legibility would have to make the value of evidence depend on its content, not only on the
amount of uncertainty it removes.

\paragraph{The payoff--record distinction.}
Although the information term treats $O$ and $R$ jointly, the axioms cannot.  Axiom~\textup{(A6)} permits the
record to be garbled while the payoff is held fixed; this separates a change in evidence from a change in material
consequences.  Axiom~\textup{(A8)} permits records to arrive in stages while the payoff is delivered once.  Its
recordwise sure-thing restriction makes branchwise comparisons commensurable and supplies the cardinal structure used to
place payoff utility and information on a common scale.  Thus either $O$ or $R$ may reveal the action, but only
$O$ enters material utility.  The record is also assumed to have no utility beyond what it reveals: a record
that is itself pleasant or shameful belongs in $O$, or the domain must be enlarged.  Once $u$ is normalised,
willingness to sacrifice payoff utility for additional mutual information identifies $\lambda$.

\paragraph{One shot, with no audience.}
The model describes one episode.  It specifies neither an audience nor any response to the beliefs induced by
the record; the posterior is a statistic of the implemented lottery and the channel, not anyone's belief.  The stages in
Axiom~\textup{(A8)} divide a single interaction: they create neither a reputational relationship nor a feedback
policy, so dynamic empowerment and the value of keeping future options open lie outside the model.  An
experiment could vary whether an audience observes the record while holding the agent's payoffs and information
fixed.  A response to that treatment would reveal an external component at work---praise, sanction, reputation;
where such responses change material consequences, they belong in $O$ and the channel.  Persistence in private
would be consistent with an internal mechanism, such as self-image, but would not identify it.  Distinguishing
these sources requires additional primitives.

\clearpage
\appendix

\part*{Proofs}
\addcontentsline{toc}{part}{Proofs}

\section{Pure trace on the constant-payoff domain}
\label{pt:sec:environment}\label{lf:app:pure-trace}

\subsection{Binary and finite branching}

We first record the finite-branch form of Axiom~\textup{(A8)} used throughout the proofs.  The result also makes
explicit where the other structural axioms enter the equivalence.

\begin{lemma}[Finite-branch extension]\label{pt:lem:finite-branch-extension}
Suppose Axioms~\textup{(A1)} and \textup{(A5)--(A7)} hold.  Then Axiom~\textup{(A8)} is equivalent to the
following property: for every $k\ge1$, every first-stage record channel
$P:A\to\D(\{1,\ldots,k\})$, every $q\in\D(A)$, and every
two ordered lists of continuation channels $(K_1,\ldots,K_k)$ and $(L_1,\ldots,L_k)$, with
$K_i,L_i:A\to\D(O\times R)$, if
\[
 (q_i,K_i)\succeq(q_i,L_i)
 \quad\text{for every reached }i,
\]
then
\begin{equation}
 \bigl(q,P\triangleright(K_1,\ldots,K_k)\bigr)
 \succeq
 \bigl(q,P\triangleright(L_1,\ldots,L_k)\bigr),
 \label{pt:eq:finite-branch-extension}
\end{equation}
strictly if one reached branch comparison is strict.
\end{lemma}

\begin{proof}
Assume first Axiom~\textup{(A8)}.  Put
\[
 m(i):=\sum_aq(a)P(i\mid a),
 \qquad
 I^+:=\{i\in\{1,\ldots,k\}:m(i)>0\}.
\]
If $m(i)=0$, then $P(i\mid a)=0$ for every $a\in\supp q$.  Consequently, changing the continuation after such a
record changes the compound channel only on action rows outside $\supp q$.  Apply Axiom~\textup{(A7)} in both
directions using the identity action processor: its Bayesian pushforwards agree on supported rows and allow
arbitrary completions elsewhere.  Continuations following unreached records are therefore immaterial.  We may
hold them fixed and enumerate the reached records as $I^+=\{i_1,\ldots,i_n\}$.

If $n=1$, the first-stage record is constant on $\supp q$.  Deleting or restoring that constant record is an
invertible record relabelling, up to rows outside $\supp q$; Axioms~\textup{(A6)} and \textup{(A7)} therefore
identify the compound comparison with the sole branch comparison.  Suppose henceforth that $n\ge2$.

For $j=0,\ldots,n$, define hybrid continuation channels by
\[
 H_i^j:=
 \begin{cases}
 K_i,&i\in\{i_1,\ldots,i_j\},\\
 L_i,&i\in I^+\setminus\{i_1,\ldots,i_j\},
 \end{cases}
\]
and use a common arbitrary continuation on $\{1,\ldots,k\}\setminus I^+$.  Write
$E_j:=P\triangleright(H_1^j,\ldots,H_k^j)$.  Fix $j\ge1$ and collapse the first-stage record to the binary channel
\[
 B_j(1\mid a):=P(i_j\mid a),
 \qquad
 B_j(2\mid a):=\sum_{i\ne i_j}P(i\mid a).
\]
Both binary records are reached.  Put all binary continuations on the common record set
$\{1,\ldots,k\}\times R$.  Pad the
two continuations that may differ by
\[
 \widetilde K_j(o,(i,r)\mid a):=\1_{\{i=i_j\}}K_{i_j}(o,r\mid a),
 \qquad
 \widetilde L_j(o,(i,r)\mid a):=\1_{\{i=i_j\}}L_{i_j}(o,r\mid a).
\]
On record $2$, collect all unchanged branches in the continuation channel
\begin{equation}
 M_j(o,(i,r)\mid a)
 :=\frac{P(i\mid a)}{B_j(2\mid a)}H_i^j(o,r\mid a),
 \qquad i\ne i_j,
 \label{pt:eq:complement-continuation}
\end{equation}
with zero probability on records $(i_j,r)$, whenever $B_j(2\mid a)>0$, and complete the other rows arbitrarily.
Since $H_i^j=H_i^{j-1}$ for $i\ne i_j$, this continuation is common to the two comparisons.
Binary record $1$ has probability
$m(i_j)$ and carries the conditional action lottery $q_{i_j}$; record $2$ carries the complementary conditional
lottery, and their mixture is $q$.  Appending or deleting the constant label $i_j$ is an invertible record
relabeling, so Axiom~\textup{(A6)} identifies the comparison of $\widetilde K_j$ and $\widetilde L_j$ with that of
$K_{i_j}$ and $L_{i_j}$.  Hence Axiom~\textup{(A8)} gives
\[
 \bigl(q,B_j\triangleright(\widetilde K_j,M_j)\bigr)
 \succeq
 \bigl(q,B_j\triangleright(\widetilde L_j,M_j)\bigr),
\]
strictly whenever the comparison at $q_{i_j}$ is strict.

The nested records $(1,(i_j,r))$ and $(2,(i,r))$ correspond bijectively to the flat records $(i,r)$ through
\[
 (1,(i_j,r))\longmapsto(i_j,r),
 \qquad
 (2,(i,r))\longmapsto(i,r)\quad(i\ne i_j).
\]
This relabelling is a payoff-preserving record processing in both directions.  Axiom~\textup{(A6)} therefore
makes the two binary compounds indifferent to $E_j$ and $E_{j-1}$, respectively.  Thus
$(q,E_j)\succeq(q,E_{j-1})$, strictly at any strict replacement.  Place the finite collection
$E_0,\ldots,E_n$ in one block environment.  Axiom~\textup{(A5)} transports each pairwise comparison to that
environment, and transitivity in Axiom~\textup{(A1)} yields $(q,E_n)\succeq(q,E_0)$, strictly if at least one
step is strict.  Restoring the immaterial null branches proves \eqref{pt:eq:finite-branch-extension}.

Conversely, assume the displayed finite-branch property, including its strict clause.  In the binary environment
of Axiom~\textup{(A8)}, reflexivity and duplication coherence make the common continuation weakly comparable to
itself, so the finite-branch property gives the forward implication in
\eqref{eq:recordwise-sure-thing}.  If the aggregate comparison in that equation held while
$(q_1,K)\not\succeq(q_1,L)$, completeness would give $(q_1,L)\succ(q_1,K)$.  Applying the strict
finite-branch property to the reversed lists would give the strict reverse aggregate comparison, a contradiction.  This
proves the converse implication and hence Axiom~\textup{(A8)}.
\end{proof}

\begin{corollary}[One-branch insertion]\label{pt:cor:one-branch-a8}
Under the hypotheses of Lemma~\ref{pt:lem:finite-branch-extension}, fix $q\in\D(A)$, a first-stage record channel
$P:A\to\D(\{1,\ldots,k\})$, a reached record $i^*$, and continuation channels
$K_1,\ldots,K_k,L_1,\ldots,L_k:A\to\D(O\times R)$ such that $K_i=L_i$ for every $i\ne i^*$.  Then
\[
 (q_{i^*},K_{i^*})\succeq(q_{i^*},L_{i^*})
 \quad\Longleftrightarrow\quad
 \bigl(q,P\triangleright(K_1,\ldots,K_k)\bigr)
 \succeq
 \bigl(q,P\triangleright(L_1,\ldots,L_k)\bigr),
\]
and the same equivalence holds for strict preference.
\end{corollary}

\begin{proof}
The forward implication follows from Lemma~\ref{pt:lem:finite-branch-extension}, using duplication coherence on
the common branches.  If the aggregate comparison held but the branch comparison failed, completeness would make
the reversed branch comparison strict; the strict part of the finite-branch extension would then give a strict
reverse aggregate comparison.  The strict equivalence follows by applying the weak equivalence in both directions.
\end{proof}

In the appendices, references to Axiom~\textup{(A8)} with more than two first-stage records mean its
finite-branch extension in Lemma~\ref{pt:lem:finite-branch-extension}; when only one branch changes, we cite
Corollary~\ref{pt:cor:one-branch-a8}.

The relevance axioms in the main text are stated inside fixed channels.  The
proofs below use their equivalent comparison-environment forms.

For $o\in O$, write $K_o:\{\ast\}\to\D(O\times\{\ast\})$ for the channel that
delivers $o$ for sure.  On a finite action set $A$, write
$K^{\mathrm{id}}_o:A\to\D(O\times A)$ for the channel that delivers $o$ and
reports the action, and $K^0_o:A\to\D(O\times\{\ast\})$ for the channel that
delivers $o$ and no record.

\begin{lemma}[Fixed-channel and environment relevance]\label{lem:relevance-bridge}
Given Axioms~\textup{(A1)}, \textup{(A5)}, \textup{(A6)},
\textup{(A7)}, and \textup{(A8)}, the following equivalences hold.
\begin{enumerate}[label=\textup{(\alph*)},leftmargin=*]
\item Axiom~\textup{(A3)} is equivalent to the existence of
  $o^+,o^-\in O$ such that
  \[
    (\delta_\ast,K_{o^+})\succ(\delta_\ast,K_{o^-}).
  \]
\item Axiom~\textup{(A4)} is equivalent to the following condition: there is
  an outcome $o_\ast\in O$ such that, for every finite $A$ with $|A|\ge2$
  and every full-support $q\in\D(A)$,
  \[
    (q,K^{\mathrm{id}}_{o_\ast})\succ(q,K^0_{o_\ast}).
  \]
\end{enumerate}
\end{lemma}

\begin{proof}
For part \textup{(a)}, apply action processing to either pure lottery in the
two-action channel of Axiom~\textup{(A3)}.  Collapsing its support to
a singleton gives $K_{o^+}$ or $K_{o^-}$.  Conversely, the singleton can be
processed back to the corresponding action; the unrestricted row at the
unreached action completes the pushforward to the original channel.  Thus
each pure lottery is indifferent to the corresponding singleton experiment.
Weak order and block coherence transport the strict comparison in either
direction.

For part \textup{(b)}, let $o_\ast$ witness Axiom~\textup{(A4)} and first
restrict its two lotteries to their supports.  Action processing in both
directions identifies the first with fair binary full revelation and the
second with a fair binary channel whose two rows are the same.  Record
processing in both directions identifies the latter channel with silence.
Hence the fixed-channel axiom is equivalent to
\begin{equation}
  ((\tfrac12,\tfrac12),K^{\mathrm{id}}_{o_\ast})
  \succ
  ((\tfrac12,\tfrac12),K^0_{o_\ast})
  \label{eq:bridge-fair-binary}
\end{equation}
on a binary action set.

Now let $p$ have full support on a binary set.  Choose
$0<\varepsilon\le 2\min_i p_i$ and an interim record $y$ with
$P(y\mid i)=\varepsilon/(2p_i)$.  The branch has probability $\varepsilon$
and posterior $(\tfrac12,\tfrac12)$.  Compare two continuation plans: after that branch the
first reveals the action and the second is silent; after every other branch
both are silent.  Equation~\eqref{eq:bridge-fair-binary} and
Corollary~\ref{pt:cor:one-branch-a8} give a strict comparison of the two compounds.  Full revelation
can be garbled into the first compound, and the second compound can be
garbled into silence.  Record processing and transitivity therefore give
$(p,K^{\mathrm{id}}_{o_\ast})\succ(p,K^0_{o_\ast})$.

Finally, let $q$ have full support on any finite $A$ with $|A|\ge2$, and
partition $A$ into two nonempty cells.  The channel that reports only the cell
is, by action processing in both directions, indifferent to binary full
revelation under the induced binary prior; silence on $A$ is similarly
indifferent to binary silence.  The cell-reporting channel is therefore
strictly preferred to silence, and full revelation is weakly preferred to
the cell report by record processing.  This proves the displayed
comparison for every $q$ and $A$ at $o_\ast$.  The converse follows by taking
a fair binary prior and reversing the support identifications used above.
\end{proof}

Whenever Axiom~\textup{(A4)} is imposed below, fix one of its witnessing
outcomes and denote it by $o_\ast$.  Every channel between nonempty finite sets
$P:A\to\D(X)$ has the constant-payoff lift
\begin{equation}
K_P(o,x\mid a):=\1_{\{o=o_\ast\}}P(x\mid a).
\label{pt:eq:constant-payoff-lift}
\end{equation}
Throughout this appendix, comparisons of pure record channels mean comparisons
of these lifts in the model already defined in the main text:
\[
q\succeq_Pq'\quad:\Longleftrightarrow\quad q\succeq_{K_P}q'.
\]
We similarly write $(q,P)\succeq(r,Q)$ for the main-text two-block
comparison of $(q,K_P)$ and $(r,K_Q)$.
Likewise, $P\sqcup Q$, block-supported lotteries, action and record
processing, ordered lists of continuation channels, and compounds
$P_1\triangleright(Q_1,\ldots,Q_k)$ are shorthand for the
corresponding main-text constructions after applying
\eqref{pt:eq:constant-payoff-lift}.  The lift commutes with each construction
up to canonical tagged or singleton-record bijections; those bijections are
neutral by Axiom~\textup{(A6)} applied in both directions.  No new primitive
relation or behavioural condition is introduced here.

Two pieces of notation will be used repeatedly.  If
$S:A\to\Delta(B)$ is an action processor, let $S^qP:B\to\Delta(X)$ denote any
Bayesian pushforward satisfying
\[
 (qS)(b)(S^qP)(x\mid b)
 =\sum_a q(a)S(b\mid a)P(x\mid a).
\]
Its rows at zero-mass reports are arbitrary, exactly as in
\eqref{eq:action-processing}.  For channels $P:A\to\Delta(X)$ and
$Q:B\to\Delta(Y)$, their independent product is
\[
 (P\otimes Q)((x,y)\mid(a,b)):=P(x\mid a)Q(y\mid b).
\]

For $q\in\D(A)$ and $P:A\to\D(X)$, write
\[
m_{qP}(x):=\sum_a q(a)P(x\mid a),
\qquad
r_x^{qP}(a):=\frac{q(a)P(x\mid a)}{m_{qP}(x)}
\quad\text{when }m_{qP}(x)>0,
\]
and define the finite posterior law
\[
\mu_{qP}:=
\sum_{x:m_{qP}(x)>0}m_{qP}(x)\,\delta_{r_x^{qP}}.
\]
Let $\Id_A$ be full revelation and $U_A$ the one-record uninformative
channel.  With logarithms in the fixed base used in the main text, put
\[
I_{qP}(A;X)
:=\Sh(m_{qP})-\sum_a q(a)\Sh(P(\cdot\mid a))
=\Sh(q)-\int_{\D(A)}\Sh(r)\,d\mu_{qP}(r).
\]

\begin{proposition}[Pure-trace lemma]\label{pt:prop:main}
Suppose the preference family satisfies main-text Axioms
\textup{(A1)}, \textup{(A2)}, and \textup{(A4)--(A8)}.  Then, for every record channel between
nonempty finite sets
$P:A\to\D(X)$ and every $q,q'\in\D(A)$,
\[
q\succeq_Pq'
\quad\Longleftrightarrow\quad
I_{qP}(A;X)\ge I_{q'P}(A;X).
\]
The same numerical scale represents every block-supported comparison: for
every nonempty finite family $\{P_i:A_i\to\D(X_i)\}_{i\in J}$ of channels
between nonempty finite sets, distinct $i,j\in J$,
and priors $q_i\in\D(A_i)$ and $q_j\in\D(A_j)$,
\[
q_i^i\succeq_{\bigsqcup_{k\in J}P_k}q_j^j
\quad\Longleftrightarrow\quad
I_{q_iP_i}(A_i;X_i)\ge I_{q_jP_j}(A_j;X_j).
\]
\end{proposition}

\subsection{Proof}

Assume the hypotheses of Proposition~\ref{pt:prop:main}.  Every appeal below
to weak order, continuity, block coherence, record processing, action
processing, the finite-branch extension of the recordwise sure-thing principle, or positive trace orientation is an
appeal to Axiom \textup{(A1)}, \textup{(A2)}, \textup{(A5)}, \textup{(A6)},
\textup{(A7)}, \textup{(A8)}, or \textup{(A4)}, respectively, on the
constant-payoff lifts.  Thus the proof has no auxiliary axiom system.

\paragraph{How to read the proof.}
The argument has five interfaces.  Each stage ends with exactly the object
needed by the next one:
\begin{enumerate}[label=\textup{Stage \arabic*:},leftmargin=*]
\item replace channels by Bayes-plausible posterior laws, derive public-mixture
  independence, and cardinalise each fixed-prior quotient using only the
  segment closedness supplied by primitive continuity;
\item make relabelling, undetectable distinctions, and restriction to the
  positive support exact consequences of action processing;
\item transport those affine values through independent products and into one
  cross-prior block scale, allowing provisionally a bilinear interaction;
\item identify reached-branch coefficients, prove their cocycle, and align
  the scales on nested support faces;
\item compare product and sequential decompositions, eliminate the
  interaction, and invoke Faddeev's recursion to obtain Shannon entropy and
  hence mutual information.
\end{enumerate}
The product-coefficient and branch tangent-space calculations are the two
technical subarguments.  A reader interested in the logical spine may use
their lemma statements as interfaces.  Product calibration is
placed before branch calibration because this makes the scalar algebra easier
to read.  The early full-revelation normalisation and exact relabelling lemma
below supply the coherence needed by that order of presentation.

\subsubsection{Stage 1: posterior-law quotient and affine fibre values}

For a finite action set $A$, let $\mathsf X_A$ be the set of finitely
supported probability measures on $\Delta(A)$, and let
\[
 b(\mu):=\int_{\Delta(A)}s\,d\mu(s),
 \qquad
 \mathcal M_q:=\{\mu\in\mathsf X_A:b(\mu)=q\}.
\]
Thus $\mathcal M_q$ is the mixture set of finite Bayes-plausible posterior
laws with prior $q$, and $\mu_{qP}\in\mathcal M_q$ for every finite channel
$P$.

\begin{lemma}[Coherent comparisons across blocks]\label{pt:lem:blockcoh}
The relation on prior--channel pairs defined above is complete and transitive.
It is unaffected by adding or deleting blocks other than the two being
compared.  In particular, for fixed $q$ the induced relation on channels is a
weak order, and
\[
 q\succeq_P r
 \quad\Longleftrightarrow\quad
 (q,P)\succeq(r,P).
\]
For fixed alphabets, its graph is closed under entrywise convergence of the
two channels and convergence of the two priors.
\end{lemma}

\begin{proof}
Completeness is Axiom~\textup{(A1)} in the two-block environment.  To prove
transitivity, suppose
$(q_1,P_1)\succeq(q_2,P_2)$ and
$(q_2,P_2)\succeq(q_3,P_3)$.  Put the three pairs in the common environment
$P_1\sqcup P_2\sqcup P_3$.  Axiom~\textup{(A5)} transports the two
hypotheses to that environment, Axiom~\textup{(A1)} composes them there, and
Axiom~\textup{(A5)} deletes the middle block.  Applying Axiom~\textup{(A5)}
with the ordered pair of block labels reversed gives
\[
 (q_2,P_2)\succeq(q_1,P_1)
 \quad\Longleftrightarrow\quad
 q_2^1\succeq_{P_1\sqcup P_2}q_1^0,
\]
so reversal is also canonical.  The displayed within-channel equivalence is
the duplication clause of Axiom~\textup{(A5)}.  Finally, a block channel is a
continuous affine function of its component channels.  Closedness on fixed
alphabets is therefore exactly Axiom~\textup{(A2)} applied to the common block
environment.
\end{proof}

\begin{lemma}[Posterior-law quotient]\label{pt:lem:blackwell}\label{pt:lem:plsuff}
Fix a full-support $q\in\Delta(A)$.  If finite channels
$P:A\to\Delta(X)$ and $Q:A\to\Delta(Y)$ satisfy
\[
 \mu_{qP}=\mu_{qQ},
\]
then they are mutually obtainable by record post-processing.  Consequently
\[
 (q,P)\sim(q,Q),
\]
and either channel may be substituted for the other in every block
comparison at prior $q$.  Hence the fixed-$q$ order of channels descends to a
well-defined weak order $\succeq_q$ on $\mathcal M_q$.
\end{lemma}

\begin{proof}
We give the finite matching argument because it also handles repeated and
zero-probability records.  Put
\[
 m_P(x):=\sum_aq(a)P(x\mid a),
 \qquad
 m_Q(y):=\sum_aq(a)Q(y\mid a).
\]
For every posterior $s$ occurring with positive probability, let
\[
 X_s:=\{x:m_P(x)>0,\ r_x^{qP}=s\},
 \qquad
 Y_s:=\{y:m_Q(y)>0,\ r_y^{qQ}=s\}.
\]
Equality of posterior laws gives the common positive mass
\[
 \lambda_s:=\sum_{x\in X_s}m_P(x)
            =\sum_{y\in Y_s}m_Q(y).
\]
For $x\in X_s$, define a stochastic record kernel by
\[
 T(y\mid x):=
 \begin{cases}
   m_Q(y)/\lambda_s,&y\in Y_s,\\
   0,&y\notin Y_s.
 \end{cases}
\]
Choose arbitrary stochastic rows at records $x$ with $m_P(x)=0$.  Full support
of $q$ implies that such a record has $P(x\mid a)=0$ for every $a$.

If $y\in Y_s$, Bayes' rule gives
\[
 \sum_{x\in X_s}P(x\mid a)
 =\frac{\lambda_s s(a)}{q(a)},
 \qquad
 Q(y\mid a)=\frac{m_Q(y)s(a)}{q(a)}.
\]
It follows that
\[
 (PT)(y\mid a)
 =\sum_{x\in X_s}P(x\mid a)\frac{m_Q(y)}{\lambda_s}
 =Q(y\mid a).
\]
Both sides vanish when $m_Q(y)=0$, again by full support.  Thus $PT=Q$.
Interchanging $P$ and $Q$ constructs $T'$ with $QT'=P$.

Axiom~\textup{(A6)} applied to $T$ and $T'$ yields the two weak comparisons;
Lemma~\ref{pt:lem:blockcoh} puts the reverse comparison into the same block
order and gives indifference.  If $P,P'$ have the same posterior law, place
$(q,P),(q,P')$ and the two objects being compared in one finite block
environment.  Transitivity then permits replacement of $P$ by $P'$ on either
side.  The quotient order is consequently independent of the implementing
channel and inherits completeness and transitivity from
Lemma~\ref{pt:lem:blockcoh}.
\end{proof}

Every element of $\mathcal M_q$ is implementable when $q$ has full support.
Indeed, if
\[
 \mu=\sum_{k=1}^n\lambda_k\delta_{s_k},
 \qquad \lambda_k>0,
 \qquad \sum_k\lambda_ks_k=q,
\]
then
\begin{equation}
 C_\mu(k\mid a):=\frac{\lambda_ks_k(a)}{q(a)}
 \label{pt:eq:canonical-implementation}
\end{equation}
defines a channel and $\mu_{qC_\mu}=\mu$.

For channels $P:A\to\Delta(X)$ and $R:A\to\Delta(Z)$ and
$t\in[0,1]$, let $tP\oplus(1-t)R$ be the publicly tagged mixture: it
reports which component was selected as well as that component's record.  Its
posterior law is
\begin{equation}
 \mu_{q,tP\oplus(1-t)R}
 =t\mu_{qP}+(1-t)\mu_{qR}.
 \label{pt:eq:public-posterior-mixture}
\end{equation}

\begin{lemma}[Public-mixture independence]\label{pt:lem:publiccoin}
Let $q$ have full support.  For all finite channels $P,Q,R$ on the action set
$A$ and every $t\in(0,1)$,
\[
 (q,P)\succeq(q,Q)
 \quad\Longleftrightarrow\quad
 (q,tP\oplus(1-t)R)
 \succeq
 (q,tQ\oplus(1-t)R).
\]
Equivalently, the quotient order on $\mathcal M_q$ satisfies mixture
independence.
\end{lemma}

\begin{proof}
If necessary, embed the record alphabets of $P,Q,R$ into one tagged disjoint
union and fill the unused coordinates with zeros.  This does not change any
posterior law, so Lemma~\ref{pt:lem:plsuff} makes the padding neutral.

Let $B:A\to\Delta(\{1,2\})$ be the uninformative first-stage channel with
$B(1\mid a)=t$ for every $a$.
Both branches are reached, and their posterior is $q$.  The compounds
\[
 B\triangleright(P,R),
 \qquad
 B\triangleright(Q,R)
\]
have the posterior laws on the right-hand side of
\eqref{pt:eq:public-posterior-mixture}; they differ from the displayed public
mixtures only by record labels.  Axiom~\textup{(A8)} gives the desired
biconditional directly.  Posterior-law substitution transfers it between the
compounds and the tagged mixtures.
\end{proof}

\begin{lemma}[Affine posterior-law representation]\label{pt:lem:postsep}
For every full-support $q\in\Delta(A)$ there is an affine functional
\[
 F_q:\mathcal M_q\longrightarrow\mathbb R
\]
such that, for all finite channels $P,Q$ on $A$,
\begin{equation}
 (q,P)\succeq(q,Q)
 \quad\Longleftrightarrow\quad
 F_q(\mu_{qP})\ge F_q(\mu_{qQ}).
 \label{pt:eq:fibre-representation}
\end{equation}
If $|A|\ge2$, $F_q$ is unique up to a positive affine transformation.  If
$|A|=1$, $\mathcal M_q$ is a singleton and we take $F_q\equiv0$.
\end{lemma}

\begin{proof}
By Lemma~\ref{pt:lem:plsuff} and
\eqref{pt:eq:canonical-implementation}, $\succeq_q$ is a weak order on the
mixture set $\mathcal M_q$.  Lemma~\ref{pt:lem:publiccoin} gives mixture
independence.  It remains only to verify Herstein--Milnor's segment-continuity
condition.  This is where primitive
Axiom~\textup{(A2)} is used; no topology on the collection of all finite
posterior laws is needed.

Fix $\mu,\nu,\rho\in\mathcal M_q$ and choose finite implementing channels
$P,Q,R$.  For $t\in[0,1]$ define
\[
 M_t:=tP\oplus(1-t)Q.
\]
All $M_t$ have the single fixed record alphabet given by the tagged union of
the record alphabets of $P$ and $Q$.  If $t_n\to t$, then $M_{t_n}\to M_t$
entrywise.  After adjoining the fixed block $R$, all comparisons occur in one
fixed finite action and record alphabet.  Axiom~\textup{(A2)} therefore makes
both sets
\[
 \{t\in[0,1]:M_t\succeq_q R\},
 \qquad
 \{t\in[0,1]:R\succeq_q M_t\}
\]
closed.  By \eqref{pt:eq:public-posterior-mixture}, these are exactly
\[
 \{t:t\mu+(1-t)\nu\succeq_q\rho\},
 \qquad
 \{t:\rho\succeq_qt\mu+(1-t)\nu\}.
\]
Thus the Herstein--Milnor mixture-space theorem
\cite{HersteinMilnor1953} applies and supplies an affine representation
$F_q$.  When $|A|\ge2$, Axiom~\textup{(A4)} makes the order nonconstant, so
the usual positive-affine uniqueness conclusion applies.  For a singleton
action set, Bayes plausibility forces every posterior law to equal
$\delta_q$.
\end{proof}

The preceding proof uses only closedness along each particular public-mixture
segment.  It neither assumes nor proves closedness under arbitrary weakly
convergent sequences whose support sizes grow.  It also requires no
pointwise integrand on $\Delta(A)$: every subsequent calculation uses only the
affinity of $F_q$ on $\mathcal M_q$ and its linear part on the affine hull.

Let
\[
 \delta_q\in\mathcal M_q
\]
be the posterior law of the uninformative channel and let
\[
 \chi_q:=\sum_{a\in A}q(a)\delta_{e_a}\in\mathcal M_q
\]
be the posterior law of full revelation.

\begin{lemma}[Canonical fibre normalization]\label{pt:lem:convnorm}
The representatives in Lemma~\ref{pt:lem:postsep} may be chosen so that
\begin{equation}
 F_q(\delta_q)=0,
 \qquad
 F_q(\mu)\ge0\quad(\mu\in\mathcal M_q),
 \label{pt:eq:zero-minimum}
\end{equation}
and, whenever $|A|\ge2$,
\begin{equation}
 F_q(\chi_q)=1.
 \label{pt:eq:full-revelation-anchor}
\end{equation}
These two anchors select a unique affine representative on every
nondegenerate full-support fibre.  On a singleton fibre we retain
$F_q\equiv0$.
\end{lemma}

\begin{proof}
Every channel $P$ can be record-postprocessed to the uninformative channel by
deleting its record.  Axiom~\textup{(A6)} and
\eqref{pt:eq:fibre-representation} therefore imply
\[
 F_q(\mu_{qP})\ge F_q(\delta_q).
\]
Every law in $\mathcal M_q$ is implemented by
\eqref{pt:eq:canonical-implementation}, so $\delta_q$ is a minimum of $F_q$
on the whole fibre.  Subtract $F_q(\delta_q)$ to obtain
\eqref{pt:eq:zero-minimum}.

If $|A|\ge2$, Axiom~\textup{(A4)} says that full revelation is strictly
preferred to no information at every full-support prior.  Hence
$F_q(\chi_q)>0$.  Dividing by this positive number gives
\eqref{pt:eq:full-revelation-anchor} without changing any comparison.  If two
affine representatives have both anchor values, positive-affine uniqueness
forces the additive constant to be zero and the positive multiplier to be
one.  The singleton convention is immediate.
\end{proof}

\subsubsection{Stage 2: transport across priors and support faces}

The remaining front-end issue is transport.  The same argument handles
renaming actions, deleting distinctions that the channel cannot detect, and
passing to the positive support of a boundary prior.  Keeping these facts in
one lemma makes clear that all three are consequences of action processing,
not extra invariance conventions.

For a bijection $\pi:A\to A'$, write $q\pi$ for the transported prior and
$\pi_\#\mu$ for the posterior law obtained by transporting every atom of
$\mu$.  If $B=\operatorname{supp}(q)$, write $q_B$ for the full-support
restriction of $q$ to $B$, $P_B$ for the restricted channel, and
$\pi_B\mu$ for the law obtained by restricting every posterior atom to $B$.

\begin{lemma}[Relabelling, undetectable distinctions, and support transport]
\label{pt:lem:transport}\label{pt:lem:actionbase}\label{pt:lem:neutrality}\label{pt:lem:supprestrict}
The canonically normalized fibre values have the following properties.

\begin{enumerate}[label=\textup{(\alph*)},leftmargin=*]
\item \emph{Exact relabelling.}
If $q$ has full support and $\pi:A\to A'$ is a bijection, then
\begin{equation}
 F_{q\pi}^{A'}(\pi_\#\mu)=F_q^A(\mu)
 \qquad(\mu\in\mathcal M_q).
 \label{pt:eq:exact-relabel}
\end{equation}

\item \emph{Undetectable distinctions.}
Let $S:A\to B$ be an onto deterministic map, let $q$ have full support, and
suppose that $P:A\to\Delta(X)$ is constant on every fibre of $S$.  Then
\begin{equation}
 (q,P)\sim(qS,S^qP).
 \label{pt:eq:undetectable}
\end{equation}
In particular, if $G_S$ reports only $S(a)$, then
$(q,G_S)\sim(qS,\Id_B)$.

\item \emph{Support restriction.}
For an arbitrary prior $q\in\Delta(A)$ with nonempty support
$B:=\operatorname{supp}(q)$ and every channel $P$,
\begin{equation}
 (q,P)\sim(q_B,P_B).
 \label{pt:eq:support-indifference}
\end{equation}
Consequently comparisons at $q$ depend only on the rows indexed by $B$, and
\begin{equation}
 (q,P)\succeq(q,Q)
 \quad\Longleftrightarrow\quad
 (q_B,P_B)\succeq(q_B,Q_B).
 \label{pt:eq:support-comparison}
\end{equation}
Every atom of a law $\mu\in\mathcal M_q$ is supported on $B$, and the
definition
\begin{equation}
 F_q^A(\mu):=F_{q_B}^{B}(\pi_B\mu)
 \label{pt:eq:boundary-representative}
\end{equation}
represents all continuation comparisons at the boundary prior $q$.
\end{enumerate}
\end{lemma}

\begin{proof}
\emph{Part (a).}
Regard $\pi$ as a deterministic action processor.  Axiom~\textup{(A7)}
gives
\[
 (q,P)\succeq(q\pi,P\pi^{-1}),
\]
where $(P\pi^{-1})(x\mid a')=P(x\mid\pi^{-1}(a'))$.  Applying the same axiom
to $\pi^{-1}$ gives the reverse comparison, so the two pairs are indifferent.
Apply this simultaneously to channels implementing arbitrary
$\mu,\nu\in\mathcal M_q$ and use Lemma~\ref{pt:lem:blockcoh} in a common
four-block environment.  One obtains
\[
 F_q^A(\mu)\ge F_q^A(\nu)
 \quad\Longleftrightarrow\quad
 F_{q\pi}^{A'}(\pi_\#\mu)
 \ge F_{q\pi}^{A'}(\pi_\#\nu).
\]
Thus $G(\mu):=F_{q\pi}^{A'}(\pi_\#\mu)$ is an affine representation of the
same nonconstant order as $F_q^A$.  Relabelling sends
$\delta_q$ to $\delta_{q\pi}$ and $\chi_q$ to $\chi_{q\pi}$.  The two anchor
equalities in Lemma~\ref{pt:lem:convnorm}, followed by affine uniqueness, give
$G=F_q^A$.  If $A$ is a singleton, both sides vanish.

\emph{Part (b).}
Since $q$ has full support and $S$ is onto, $qS$ has full support.  One
direction of \eqref{pt:eq:undetectable} is Axiom~\textup{(A7)}.  For the
reverse direction, define the Bayesian refinement kernel
\[
 R(a\mid b):=
 \begin{cases}
 q(a)/(qS)(b),&S(a)=b,\\
 0,&S(a)\ne b.
 \end{cases}
\]
Then $(qS)R=q$.  Since $P$ is constant on each fibre of $S$, the pushforward
$S^qP$ has at $b$ precisely the common row of $P$ on $S^{-1}(b)$, and pushing
it back through $R$ recovers $P$.  Axiom~\textup{(A7)} applied to
$(qS,S^qP)$ and $R$ therefore gives the reverse weak comparison.  For the
last assertion, $G_S$ is $S$-measurable and its pushforward is $\Id_B$.

\emph{Part (c).}
Choose $b_0\in B$ and let $S:A\to B$ fix every $b\in B$ and send every
$a\notin B$ to $b_0$.  Since $q$ assigns zero mass outside $B$,
\[
 qS=q_B,
 \qquad
 S^qP=P_B.
\]
Axiom~\textup{(A7)} gives $(q,P)\succeq(q_B,P_B)$.  Conversely, let
$\iota:B\hookrightarrow A$ be the inclusion.  Then $q_B\iota=q$.  The exact
Bayesian completion equation pins down the processed channel on $B$ and leaves
all zero-mass rows in $A\setminus B$ arbitrary.  The all-completions clause of
Axiom~\textup{(A7)} therefore permits us to choose the completion to be the
original $P$, yielding $(q_B,P_B)\succeq(q,P)$.  This proves
\eqref{pt:eq:support-indifference}; applying it to $P$ and $Q$ in a common
block environment proves \eqref{pt:eq:support-comparison}.  It also shows
immediately that channels agreeing on $B$ are indifferent at $q$.

Finally, if $\mu=\sum_k\lambda_k\delta_{s_k}\in\mathcal M_q$ with
$\lambda_k>0$, then for $a\notin B$,
\[
 0=q(a)=\sum_k\lambda_ks_k(a).
\]
Nonnegativity forces $s_k(a)=0$ for every $k$.  Thus every atom lies on the
face $\Delta(B)$ and $\pi_B\mu\in\mathcal M_{q_B}$.  The comparison
equivalence just proved and the full-support representation on $B$ show that
\eqref{pt:eq:boundary-representative} represents precisely the boundary
continuation order.  This definition is also compatible with relabelling by
part~(a).
\end{proof}

At this point every later calculation may work entirely with the affine
functionals $F_q$, their linear parts on the affine hulls of
$\mathcal M_q$, and the exact transport identities above.  No global
posterior-law continuity statement and no posterior-separable integrand are
needed.  From now on ambient action-set superscripts are suppressed unless
two action sets occur in the same formula; a boundary-prior subscript always
means the representative transported from its positive support.  The next two
stages derive product and branch scale coherence, with Axiom~\textup{(A8)} and its finite-branch extension providing
the essential cardinal alignment across reached
branches.

\subsubsection{Stage 3: products and a common comparison scale}

The value $F_q$ compares experiments only within the fibre of one prior.
Embedding experiments from two priors into their product fibre will put both
values on one cardinal comparison scale.  We first derive the only
separability fact about independent products needed for that construction.
Write $P\succeq_qQ$ for $(q,P)\succeq(q,Q)$, and similarly for strict
preference.  For a channel $P:A\to\Delta(X)$, let $\widetilde P$ denote its
lift to $A\times B$, constant in the $B$-coordinate; for
$R:B\to\Delta(Y)$, let $\widetilde R$ denote the symmetric lift, constant in
the $A$-coordinate.
Posterior laws may be substituted throughout by Lemma~\ref{pt:lem:plsuff};
in particular, adding null records and permuting record labels has no effect
on a comparison.

\begin{lemma}[Cancellation of a common independent background]
\label{pt:lem:background}
Let $p\in\Delta(A)$ and $r\in\Delta(B)$ have full support.  For channels
$P,Q$ on $A$ and a channel $R$ on $B$,
\begin{equation}
 P\otimes R\succeq_{p\otimes r}Q\otimes R
 \quad\Longleftrightarrow\quad
 P\succeq_p Q .
\tag{BI}\label{pt:eq:backgroundinert}
\end{equation}
The symmetric assertion holds with the two coordinates interchanged.  The
equivalence, including its strict form, remains valid when a posterior reached
in the background coordinate lies on a proper support face.
\end{lemma}

\begin{proof}
Pad $P$ and $Q$ by null records so that they have the same record alphabet.
Regard the lift $\widetilde R$ of $R$ as a first-stage experiment on
$A\times B$, followed on every reached branch by $\widetilde P$, respectively
$\widetilde Q$.  Up to the harmless permutation of the two record
coordinates, the resulting compounds are $P\otimes R$ and $Q\otimes R$.

If $y$ is reached under $R$, its posterior on $A\times B$ is
\[
 p\otimes r_y,
 \qquad r_y:=r_y^{\,r,R}.
\]
The coordinate projection $A\times\operatorname{supp}(r_y)\to A$ and the
support-transport clauses of Lemmas~\ref{pt:lem:neutrality} and
\ref{pt:lem:supprestrict} give
\[
 (p\otimes r_y,\widetilde P)\sim(p,P),
 \qquad
 (p\otimes r_y,\widetilde Q)\sim(p,Q).
\tag{BI$'$}\label{pt:eq:bg-branch-transport}
\]
Thus every reached branch has the same weak comparison as $P$ versus $Q$ at
$p$.  If $P\succeq_pQ$, Lemma~\ref{pt:lem:finite-branch-extension} applied to
these branches proves the forward implication in~\eqref{pt:eq:backgroundinert}.

For the converse, suppose that $P\not\succeq_pQ$.  Completeness gives
$Q\succ_pP$.  By~\eqref{pt:eq:bg-branch-transport}, the reversed comparison is
strict on every reached branch; at least one such branch has positive mass.
The strict part of Lemma~\ref{pt:lem:finite-branch-extension} therefore gives
$Q\otimes R\succ_{p\otimes r}P\otimes R$, contradicting
$P\otimes R\succeq_{p\otimes r}Q\otimes R$.  This proves the converse and
also records why weak branch monotonicity alone would not suffice.  Interchange
the coordinates for the symmetric assertion.
\end{proof}

For finite posterior laws define
\[
 \left(\sum_i m_i\delta_{u_i}\right)\otimes
 \left(\sum_j n_j\delta_{v_j}\right)
 :=\sum_{i,j}m_i n_j\delta_{u_i\otimes v_j}.
\]
Bayes' rule gives
\begin{equation}
 \mu_{p\otimes r,P\otimes R}=\mu_{pP}\otimes\mu_{rR}.
\tag{PL}\label{pt:eq:productlaw}
\end{equation}

\begin{lemma}[Product gauge]
\label{pt:lem:coherentnorm}
There are positive rescalings $G_q=c_qF_q$ of the zero-normalised affine
representatives and a single number $\kappa\in\mathbb R$ such that
\begin{equation}
 G_{p\otimes r}(\mu\otimes\nu)
 =G_p(\mu)+G_r(\nu)+\kappa G_p(\mu)G_r(\nu)
\tag{QA}\label{pt:eq:QA}
\end{equation}
for all full-support priors $p,r$ and all
$\mu\in\mathcal M_p$, $\nu\in\mathcal M_r$.  These rescalings can be chosen
exactly invariant under bijections.  If a factor is a singleton, its
representative is zero and the same formula holds.  Finally,
\begin{equation}
 1+\kappa G_r(\nu)>0
\tag{POS}\label{pt:eq:QAslope}
\end{equation}
whenever the $r$-fibre is nonconstant.
\end{lemma}

\begin{proof}
We give the algebra because both the common value of $\kappa$ and the strict
inequality~\eqref{pt:eq:QAslope} are used later.  Initially retain the notation
$F_q$.  For nondegenerate $p,r$, put
\[
 L_{p,r}(\mu,\nu):=F_{p\otimes r}(\mu\otimes\nu),
 \qquad x:=F_p(\mu),\quad y:=F_r(\nu).
\]
The map $L_{p,r}$ is separately affine.  By
Lemma~\ref{pt:lem:background}, every first-coordinate slice represents the
$F_p$-order and every second-coordinate slice represents the $F_r$-order.
Affine-utility uniqueness, first on a slice and then as the other coordinate
varies affinely, therefore gives unique constants
$A_{p,r},B_{p,r}>0$ and $C_{p,r}\in\mathbb R$ such that
\begin{equation}
 L_{p,r}(\mu,\nu)
 =A_{p,r}x+B_{p,r}y+C_{p,r}xy.
\label{pt:eq:pairbilinear}
\end{equation}
Here is the coefficient argument.  Fixing $\nu$ first gives
$L_{p,r}=\alpha(\nu)x+\gamma(\nu)$ with $\alpha(\nu)>0$.  At
$\mu=\delta_p$, the function
$\gamma(\nu)=L_{p,r}(\delta_p,\nu)$ is an affine representative of the
$F_r$-order and vanishes at $\delta_r$; hence
$\gamma(\nu)=B_{p,r}y$ for some $B_{p,r}>0$.  Choose
$\bar\mu$ with $\bar x:=F_p(\bar\mu)\ne0$.  The other slice has the form
\[
 L_{p,r}(\bar\mu,\nu)=A_{p,r}\bar x+D_{p,r}y,
\]
where $A_{p,r}>0$ is its value slope at $\nu=\delta_r$.  Comparing this with
$\alpha(\nu)\bar x+B_{p,r}y$ shows that
$\alpha(\nu)=A_{p,r}+C_{p,r}y$, where
$C_{p,r}:=(D_{p,r}-B_{p,r})/\bar x$.  This proves
\eqref{pt:eq:pairbilinear}; the symmetric calculation shows that the two
slice slopes obey
\begin{equation}
 A_{p,r}+C_{p,r}y>0,
 \qquad B_{p,r}+C_{p,r}x>0
\label{pt:eq:pair-slopes}
\end{equation}
throughout the two value ranges.

Exact relabelling makes the evaluations under the two associations of a
triple product equal.  Comparing the coefficients of $x,y,z$ gives
\begin{align}
 A_{p\otimes r,s}A_{p,r}&=A_{p,r\otimes s},\label{pt:eq:Acoc1}\\
 A_{p\otimes r,s}B_{p,r}&=B_{p,r\otimes s}A_{r,s},\label{pt:eq:Acoc2}\\
 B_{p\otimes r,s}&=B_{p,r\otimes s}B_{r,s}.
\label{pt:eq:Acoc3}
\end{align}
The coordinate swap similarly gives
\[
 B_{r,p}=A_{p,r},\qquad A_{r,p}=B_{p,r},\qquad C_{r,p}=C_{p,r}.
\tag{SW}\label{pt:eq:coefficient-swap}
\]
This coefficient comparison is legitimate: the laws
$(1-t)\delta_p+t\chi_p$ make $x$ range over a nondegenerate interval, and
the same holds for $y$ and $z$.
Let $\rho(p,r):=B_{p,r}/A_{p,r}$.  Dividing
\eqref{pt:eq:Acoc2} by~\eqref{pt:eq:Acoc1} yields
\[
 \rho(p,r)=\rho(p,r\otimes s)A_{r,s}.
\]
Fix a nondegenerate reference prior $q_0$ and set
$g(r):=\rho(q_0,r)$.  Then
\begin{equation}
 A_{r,s}=\frac{g(r)}{g(r\otimes s)}.
\label{pt:eq:gaugeA}
\end{equation}
The uniqueness of the coefficients in~\eqref{pt:eq:pairbilinear}, applied
after product relabellings, shows that $g$ is invariant under bijections.
Consequently the positive rescaling $G_p:=g(p)F_p$ preserves exact
relabelling.  Equation~\eqref{pt:eq:gaugeA} makes the new first coefficient
equal to one, and~\eqref{pt:eq:coefficient-swap} makes the second coefficient
equal to one as well.  We henceforth use $G$ for these gauged representatives.

Write the remaining coefficient as $\kappa_{p,r}$.  Associativity of
\[
 x\oplus_{p,r}y:=x+y+\kappa_{p,r}xy
\]
and comparison of the coefficients of $xy,xz,yz$ give
\[
 \kappa_{p,r}=\kappa_{p,r\otimes s},\qquad
 \kappa_{p\otimes r,s}=\kappa_{p,r\otimes s},\qquad
 \kappa_{p\otimes r,s}=\kappa_{r,s}.
\tag{KA}\label{pt:eq:kappa-assoc}
\]
(The $xyz$ coefficient gives the product of the corresponding two sides and
adds no restriction.)  Taking $s=q_0$ in~\eqref{pt:eq:kappa-assoc} shows that
$\kappa_{p,r}=\kappa_{r,q_0}$, so it is independent of $p$; the swap in
\eqref{pt:eq:coefficient-swap} makes it symmetric, hence independent of $r$.
Call the common value $\kappa$.  After the gauge,
\eqref{pt:eq:pair-slopes} is exactly
$1+\kappa G_r(\nu)>0$ and its symmetric counterpart.

For a singleton prior set $G\equiv0$; the canonical bijections
$A\cong A\times\{*\}$ give~\eqref{pt:eq:QA}.  For a boundary prior, define
$G_r$ by exact transport from its positive-probability support, as in
Lemma~\ref{pt:lem:supprestrict}.  This convention is relabelling invariant.
\end{proof}

\begin{corollary}[Exact relabelling after the product gauge]
\label{pt:cor:permutationinvariance}
For every bijection $\pi:A\to A'$, prior $q\in\Delta(A)$, and
$\mu\in\mathcal M_q$,
\[
 G_{q\pi}^{A'}(\pi_\#\mu)=G_q^A(\mu).
\]
\end{corollary}

\begin{proof}
The initial equality is Lemma~\ref{pt:lem:actionbase}; the proof of
Lemma~\ref{pt:lem:coherentnorm} shows that its gauge factor is invariant under
bijections.
\end{proof}

\begin{lemma}[Cross-prior block representation]
\label{pt:lem:blockbridge}
For arbitrary finite priors $q\in\Delta(A)$ and $r\in\Delta(B)$ and channels
$P,Q$ on the respective action sets,
\begin{equation}
 q^0\succeq_{P\sqcup Q}r^1
 \quad\Longleftrightarrow\quad
 G_q(\mu_{qP})\ge G_r(\mu_{r,Q}).
\tag{BC}\label{pt:eq:blockcomp}
\end{equation}
\end{lemma}

\begin{proof}
First suppose that $q,r$ have full support.  By~\eqref{pt:eq:QA} and
$G_q(\delta_q)=G_r(\delta_r)=0$,
\[
 G_{q\otimes r}(\mu_{qP}\otimes\delta_r)=G_q(\mu_{qP}),
 \qquad
 G_{q\otimes r}(\delta_q\otimes\mu_{r,Q})=G_r(\mu_{r,Q}).
\]
The two laws on the left lie in the same fibre and are implemented by
$P\otimes U_B$ and $U_A\otimes Q$, so the affine fibre representation
compares them by the displayed numbers.

It remains only to remove the uninformative product coordinates.  Projection
$A\times B\to A$ and Axiom~\textup{(A7)} give
$(q\otimes r,P\otimes U_B)\succeq(q,P)$.  Conversely the stochastic
embedding
\[
 T((a,b)\mid a')=\mathbf 1_{\{a=a'\}}r(b)
\]
sends $q$ to $q\otimes r$ and its Bayesian pushforward sends $P$ to
$P\otimes U_B$; Axiom~\textup{(A7)} gives the reverse comparison.  Thus the
two pairs are indifferent.  The same argument identifies
$(q\otimes r,U_A\otimes Q)$ with $(r,Q)$.  Block coherence and transitivity
now prove~\eqref{pt:eq:blockcomp}.  Arbitrary priors follow by restricting
both pairs and both representatives to their supports; singleton supports use
the zero convention.  The singleton record tags created by these projections
and embeddings are the canonical relabellings covered by the convention at
the start of the appendix.
\end{proof}

\subsubsection{Stage 4: sequential decomposition and coherent face scales}

Changing a continuation on one reached branch changes the posterior law by a
signed measure with zero total mass and zero barycentre.  The following space
is exactly the space of such feasible directions.
For a nonempty finite set $A$, let
\[
 W_A:=\left\{\eta:
 \begin{array}{l}
 \eta\text{ is a finitely supported signed measure on }\Delta(A),\\[-2pt]
 \eta(\Delta(A))=0,\quad \displaystyle\int p\,d\eta(p)=0
 \end{array}\right\}.
\]
If $\iota:B\hookrightarrow A$ is an injection, let
$i_\iota:W_B\to W_A$ push every posterior forward along $\iota$; for a
literal subset $B\subseteq A$, write $i_B$ for the inclusion case.  Let
$L_q^A$ be the linear part of the affine functional $G_q$ on the affine hull
of $\mathcal M_q$.

\begin{lemma}[Branch aggregation]
\label{pt:lem:branchagg}
Let $q$ have full support on $A$.  For every reached posterior $r$ with
$B:=\operatorname{supp}(r)$ and $|B|\ge2$, there is a unique
$\beta_A(q,r)>0$ satisfying
\begin{equation}
 L_q^A\circ i_B=\beta_A(q,r)L_{r_B}^B
 \quad\hbox{on }W_B.
\label{pt:eq:branch-linear}
\end{equation}
It depends only on $q,r$ and the face inclusion, not on the experiment that
reaches $r$.  If a first-stage channel $P_1:A\to\Delta(\{1,\ldots,k\})$ has branch masses
$m_i$ and posteriors $r_i$, then every list of continuation channels $Q_1,\ldots,Q_k$ satisfies
\begin{equation}
 G_q\bigl(\mu_{q,P_1\triangleright(Q_1,\ldots,Q_k)}\bigr)
 =G_q(\mu_{qP_1})+
 \sum_{i:m_i>0}m_i\beta_A(q,r_i)
 G_{r_i}(\mu_{r_i,Q_i}).
\tag{BA}\label{pt:eq:branchaggregation}
\end{equation}
For singleton $r_i$, the continuation value is zero and the corresponding
coefficient may be fixed by any positive convention.
\end{lemma}

\begin{proof}
We first record the common tangent space.  If $q$ has full support, then
\begin{equation}
 \operatorname{span}\{\mu-\nu:\mu,\nu\in\mathcal M_q\}=W_A.
\label{pt:eq:tangent-space}
\end{equation}
Only the reverse inclusion needs proof.  For $0\ne\eta\in W_A$, write its
Jordan decomposition as $\eta=\eta^+-\eta^-$, where both positive parts have
the same mass $c>0$ and, after division by $c$, the same barycentre $b$.
Choose $\varepsilon>0$ small enough that
$\tau=(q-\varepsilon b)/(1-\varepsilon)$ belongs to $\Delta(A)$.  Then
\[
 \mu^\pm:=\frac{\varepsilon}{c}\eta^\pm
          +(1-\varepsilon)\delta_\tau
 \in\mathcal M_q,
 \qquad
 \eta=\frac c\varepsilon(\mu^+-\mu^-),
\]
which proves~\eqref{pt:eq:tangent-space}.  The same argument intrinsically on
a face gives the space $W_B$.

Fix a nondegenerate $r$ supported on $B$.  It is reachable from $q$: choose
\[
 0<\varepsilon<\min\left\{1,
       \min_{b\in B}\frac{q(b)}{r(b)}\right\}
\]
and use a two-record experiment whose first record has conditional probability
$\varepsilon r(a)/q(a)$.  That record has mass $\varepsilon$ and posterior
$r$.

Implement $\nu,\nu'\in\mathcal M_{r_B}$ by channels on $B$, extend their
off-support rows arbitrarily to $A$, and pad their record alphabets if
necessary.  Vary only the continuation on the branch reaching $r$ from
$\nu$ to $\nu'$.  The aggregate posterior law changes by
$\varepsilon i_B(\nu-\nu')$.  Corollary~\ref{pt:cor:one-branch-a8}, support transport, and strictness give all
three signs:
\begin{align*}
 L_{r_B}^B(\nu-\nu')>0&\Longrightarrow
     L_q^A i_B(\nu-\nu')>0,\\
 L_{r_B}^B(\nu-\nu')<0&\Longrightarrow
     L_q^A i_B(\nu-\nu')<0,\\
 L_{r_B}^B(\nu-\nu')=0&\Longrightarrow
     L_q^A i_B(\nu-\nu')=0.
\end{align*}
The second line is the first applied after swapping the continuations.  For
the third, indifference supplies both weak branch comparisons, and
Corollary~\ref{pt:cor:one-branch-a8} supplies both aggregate inequalities.  By
\eqref{pt:eq:tangent-space}, the sign agreement extends by positive scaling
to all of $W_B$.  Positive trace orientation makes
\[
 L_{r_B}^B(\chi_{r_B}-\delta_{r_B})
 =G_{r_B}(\chi_{r_B})>0,
\]
so $L_{r_B}^B\ne0$.  Two nonzero linear functionals with the same positive
half-space have the same kernel and are positive scalar multiples; this gives
the unique coefficient in~\eqref{pt:eq:branch-linear}.  Since both linear
functionals in that equation are fixed before a reaching experiment is
chosen, the coefficient is path-independent.

For the aggregation formula, put $B_i:=\operatorname{supp}(r_i)$ and
\[
 \nu_i:=\pi_{B_i}\mu_{r_i,Q_i}\in\mathcal M_{(r_i)_{B_i}}.
\]
The compound law differs from the first-stage law by
\[
 \mu_{q,P_1\triangleright(Q_1,\ldots,Q_k)}-\mu_{qP_1}
 =\sum_{i:m_i>0}m_i\,i_{B_i}
   (\nu_i-\delta_{(r_i)_{B_i}}).
\]
Apply $L_q^A$, use~\eqref{pt:eq:branch-linear}, and use
$G_{r_i}(\delta_{r_i})=0$.  Singleton summands vanish.  This is
\eqref{pt:eq:branchaggregation}.
\end{proof}

For an injection $\iota:B\hookrightarrow A$ and full-support
$r\in\Delta(B)$, define
\[
 \beta_\iota(q,r):=\beta_A(q,\iota_\#r).
\]
The image posterior $\iota_\#r$ is reachable by the two-record construction
in the proof, and its support inclusion is precisely $i_\iota$.  Thus
\[
 L_q^A i_\iota=\beta_\iota(q,r)L_r^B
 \quad\hbox{on }W_B.
\]

\begin{lemma}[Cocycle, face coherence, and the chain gauge]
\label{pt:lem:chain}
The coefficients of Lemma~\ref{pt:lem:branchagg} admit positive scales
$a_q^A$ which are invariant under relabelling and satisfy
\begin{equation}
 \beta_\iota(q,r)=\frac{a_q^A}{a_r^B}
\tag{FS}\label{pt:eq:facescale}
\end{equation}
whenever $\iota:B\hookrightarrow A$, both sets are nondegenerate, $q$ and
$r$ have full support on their respective sets, and $\beta_\iota(q,r)$ is the
coefficient for reaching the face $\iota(B)$.  Hence the representatives
\[
 V_q:=\frac{G_q}{a_q}
\]
satisfy, for every $P_1:A\to\Delta(\{1,\ldots,k\})$ with branch masses $m_i$ and posteriors $r_i$, and every
list of continuation channels $Q_1,\ldots,Q_k$, the exact chain rule
\begin{equation}
 V_q\bigl(\mu_{q,P_1\triangleright(Q_1,\ldots,Q_k)}\bigr)
 =V_q(\mu_{qP_1})
  +\sum_{i:m_i>0}m_iV_{r_i}(\mu_{r_i,Q_i}).
\tag{CR}\label{pt:eq:chainrule}
\end{equation}
Singleton continuation terms are zero.
\end{lemma}

\begin{proof}
There are two normalization tasks: first solve the branch cocycle within each
fixed action-set dimension, then remove the residual scalar attached to
embedding an $m$-point face in an $n$-point simplex.

We first prove the cocycle before choosing scales.  For nested injections
$C\xhookrightarrow{\jmath}B\xhookrightarrow{\iota}A$ and full-support
priors $q,r,s$ on $A,B,C$, respectively, the defining identities on $W_C$
are
\[
 L_q^A i_{\iota\circ\jmath}
 =\beta_\iota(q,r)L_r^B i_\jmath
 =\beta_\iota(q,r)\beta_\jmath(r,s)L_s^C.
\]
The direct coefficient from $q$ to $s$ is unique because $L_s^C\ne0$ by
Axiom~\textup{(A4)}.  Therefore
\begin{equation}
 \beta_{\iota\circ\jmath}(q,s)
 =\beta_\iota(q,r)\beta_\jmath(r,s).
\label{pt:eq:branch-cocycle}
\end{equation}
This single identity covers full-support paths, full-to-face paths, and
nested support faces.

For a fixed nondegenerate $A$, apply~\eqref{pt:eq:branch-cocycle} to identity
embeddings.  Fix a full-support basepoint $q_A$ and put
\[
 a_q^A:=\beta_{\mathrm{id}_A}(q,q_A).
\]
Then
\[
 \beta_{\mathrm{id}_A}(q,r)=\frac{a_q^A}{a_r^A}.
\tag{WS}\label{pt:eq:within-scale}
\]
Choosing $q_A$ to be uniform and using exact relabelling of $G$ makes these
initial scales invariant under bijections.  They remain free up to one
positive multiplier for each cardinality.

For an injection $\iota:B\hookrightarrow A$, define its residual defect by
\[
 K(\iota):=\beta_\iota(q,r)\frac{a_r^B}{a_q^A}.
\]
Equation~\eqref{pt:eq:within-scale} shows that this number is independent of
$q$ and $r$: changing either prior merely inserts and cancels the
corresponding within-set coefficient.  Equation~\eqref{pt:eq:branch-cocycle}
gives
\[
 K(\iota\circ\jmath)=K(\iota)K(\jmath).
\]
Exact relabelling implies that $K(\iota)$ depends only on
$n=|A|$ and $m=|B|$; write it $K_{n,m}$.  Thus
\[
 K_{n,\ell}=K_{n,m}K_{m,\ell},\qquad K_{n,n}=1.
\]
Set $t_n:=K_{n,2}$ and $t_2:=1$.  Then
$K_{n,m}=t_n/t_m$.  Replacing every $a_q^A$ on an $n$-point set by
$t_n a_q^A$ changes the defect to
$K_{n,m}t_m/t_n=1$, while preserving~\eqref{pt:eq:within-scale} and
relabel invariance.  This proves~\eqref{pt:eq:facescale}.

For a boundary prior, $a_r$ means the scale on its support; the preceding
construction says precisely that this agrees with the scale obtained by
reaching that support as a face.  Divide~\eqref{pt:eq:branchaggregation} by
$a_q$ and use~\eqref{pt:eq:facescale}.  Every nondegenerate branch coefficient
becomes
\[
 \frac1{a_q}\,\beta_A(q,r_i)G_{r_i}
 =\frac1{a_{r_i}}G_{r_i}=V_{r_i}.
\]
Singleton values vanish, so their scales and coefficients may be fixed
arbitrarily.  This proves~\eqref{pt:eq:chainrule}.
\end{proof}

\subsubsection{Stage 5: compatibility and entropy identification}

The product and sequential evaluations of full revelation now meet.  Their
compatibility first removes the provisional product interaction and then
leaves a strongly additive uncertainty functional.

\begin{lemma}[Interaction collapse and universal chain scale]
\label{pt:lem:scalecoherence}
The coefficient $\kappa$ in~\eqref{pt:eq:QA} is zero.  Moreover there is a
single $C>0$ such that $a_q=C$ for every nondegenerate finite prior, read on
its support.  The same value may be assigned to singleton priors.
Consequently the final representatives $V_q=G_q/C$ satisfy both
\begin{align}
 V_{p\otimes r}(\mu\otimes\nu)
   &=V_p(\mu)+V_r(\nu),\tag{PA}\label{pt:eq:PArecovered}\\
 V_q\bigl(\mu_{q,P_1\triangleright(Q_1,\ldots,Q_k)}\bigr)
   &=V_q(\mu_{qP_1})+\sum_{i:m_i>0}m_i
     V_{r_i}(\mu_{r_i,Q_i}).\tag{CR$'$}\label{pt:eq:final-chain}
\end{align}
for every $P_1:A\to\Delta(\{1,\ldots,k\})$ with branch masses $m_i$ and posteriors $r_i$, and every list of
continuation channels $Q_1,\ldots,Q_k$.
The cross-prior bridge of Lemma~\ref{pt:lem:blockbridge} is valid with $V$
in place of $G$.
\end{lemma}

\begin{proof}
In this proof only, put
\[
 H_G(q):=G_q(\chi_q),\qquad Z(q):=1+\kappa H_G(q).
\]
For nondegenerate $q$, Axiom~\textup{(A4)} and zero normalisation give
$H_G(q)>0$, while the strict slice-slope
\eqref{pt:eq:QAslope} gives
\begin{equation}
 Z(q)>0.
\label{pt:eq:Zpositive}
\end{equation}
For a singleton, $H_G=0$ and $Z=1$.

\emph{Product revelation versus sequential revelation.}
For nondegenerate full-support $p,r$, product quasi-additivity gives
\begin{equation}
 H_G(p\otimes r)
 =H_G(p)+H_G(r)+\kappa H_G(p)H_G(r).
\tag{QH}\label{pt:eq:QH}
\end{equation}
Compute the same full revelation by first revealing the $A$-coordinate and
then the $B$-coordinate.  The first-stage law is
$\chi_p\otimes\delta_r$ and has value $H_G(p)$ by~\eqref{pt:eq:QA}.
After action $a$ is revealed, the posterior is $e_a\otimes r$ on the face
$\{a\}\times B$; support restriction and relabelling identify the
continuation value with $H_G(r)$.  Every such branch has coefficient
$a_{p\otimes r}/a_r$ by~\eqref{pt:eq:facescale}.  Since the branch weights
sum to one, Lemma~\ref{pt:lem:branchagg} gives
\begin{equation}
 H_G(p\otimes r)=H_G(p)+\frac{a_{p\otimes r}}{a_r}H_G(r).
\tag{SA}\label{pt:eq:seqA}
\end{equation}
Comparing~\eqref{pt:eq:QH} and~\eqref{pt:eq:seqA}, and using
$H_G(r)>0$, yields
\[
 \frac{a_{p\otimes r}}{a_r}=Z(p).
\tag{R1}\label{pt:eq:ratio1}
\]
Revealing $B$ first gives symmetrically
\[
 \frac{a_{p\otimes r}}{a_p}=Z(r).
\tag{R2}\label{pt:eq:ratio2}
\]
Thus $a_rZ(p)=a_pZ(r)$ for every pair of nondegenerate priors, even on
different finite sets.  By~\eqref{pt:eq:Zpositive}, there is a universal
$C>0$ such that
\begin{equation}
 a_q=CZ(q).
\tag{AS}\label{pt:eq:ascaleZ}
\end{equation}
Assign $a=C$ on singleton supports.

If $\kappa=0$, then $Z\equiv1$ and the conclusion already follows.  It
remains to rule out $\kappa\ne0$.

\emph{The grouping weight.}
Let $q$ have full support and let
$\mathcal P=\{A_1,\ldots,A_m\}$ partition its action set into nonempty
cells.  Put
\[
 p_i=q(A_i),\qquad q^i=q(\,\cdot\mid A_i),\qquad
 p=(p_1,\ldots,p_m).
\]
If $C_{\mathcal P}$ reports only the cell, undetectable-distinctions
neutrality gives
\[
 (q,C_{\mathcal P})\sim(p,\Id_{\{1,\ldots,m\}}).
\]
The cross-prior bridge therefore identifies the coarse value as
\begin{equation}
 G_q(\mu_{q,C_{\mathcal P}})=H_G(p).
\tag{BG}\label{pt:eq:blockvalue}
\end{equation}
Append full revelation within each cell.  The final law is $\chi_q$, and the
branch formula together with~\eqref{pt:eq:ascaleZ} gives
\begin{equation}
 H_G(q)=H_G(p)+Z(q)\sum_i p_i\frac{H_G(q^i)}{Z(q^i)}.
\tag{GR}\label{pt:eq:groupingH}
\end{equation}
Here singleton cells contribute zero.

Set $w(q):=1/Z(q)>0$.  Because $\kappa\ne0$ in the present paragraph,
$\kappa H_G(x)=Z(x)-1$.  Substitution into
\eqref{pt:eq:groupingH} gives
\[
 Z(q)-1=Z(p)-1+
 Z(q)\sum_i p_i\left(1-\frac1{Z(q^i)}\right).
\]
Cancelling and inverting yields the positive grouping-weight equation
\begin{equation}
 w(q)=w(p)\sum_i p_iw(q^i).
\tag{W}\label{pt:eq:weight}
\end{equation}
All distributions in what follows are read on their positive-probability
supports.  For $u\otimes v$, partition by the first coordinate: the coarse
distribution is $u$ and every conditional distribution is $v$.  Thus
\eqref{pt:eq:weight} gives
\begin{equation}
 w(u\otimes v)=w(u)w(v).
\tag{M}\label{pt:eq:wmult}
\end{equation}

\emph{Two groupings force the weight to be constant.}
Take arbitrary finite probability vectors $u,v$ and write
$x=w(u)$ and $y=w(v)$.  Let $T$ be the law of a triple
$(\ell,z_1,z_2)$: $\ell$ is uniform on $\{0,1\}$, and conditional on
$\ell=0$ the pair $(z_1,z_2)$ has law $u\otimes u$, while conditional on
$\ell=1$ it has law $v\otimes v$.  In tagged notation,
\[
 T:=\tfrac12(u\otimes u)^0\sqcup\tfrac12(v\otimes v)^1.
\]
Grouping first by $\ell$ and using~\eqref{pt:eq:wmult} gives
\begin{equation}
 w(T)=w(\tfrac12,\tfrac12)\frac{x^2+y^2}{2}.
\tag{E1}\label{pt:eq:twoeval1}
\end{equation}
Instead group first by $(\ell,z_1)$.  That marginal is
$S=\tfrac12u^0\sqcup\tfrac12v^1$, so
\[
 w(S)=w(\tfrac12,\tfrac12)\frac{x+y}{2}.
\]
The conditional law of $z_2$ is $u$ when $\ell=0$ and $v$ when $\ell=1$.
A second use of~\eqref{pt:eq:weight} therefore gives
\begin{equation}
 w(T)=w(S)\frac{x+y}{2}
 =w(\tfrac12,\tfrac12)\left(\frac{x+y}{2}\right)^2.
\tag{E2}\label{pt:eq:twoeval2}
\end{equation}
The common leading factor is positive.  Comparing
\eqref{pt:eq:twoeval1} and~\eqref{pt:eq:twoeval2} yields $(x-y)^2=0$.
Thus $w$ is constant on all finite probability vectors.  Its singleton value
is one, so $w\equiv1$.

It follows that $Z\equiv1$, hence $\kappa H_G(q)=0$ for every $q$.
Axiom~\textup{(A4)} supplies a nondegenerate $q$ with $H_G(q)>0$, so
$\kappa=0$, contradicting the temporary alternative.  Equation
\eqref{pt:eq:ascaleZ} now gives $a_q=C$ universally.
Dividing the product and branch formulas by $C$ proves
\eqref{pt:eq:PArecovered} and~\eqref{pt:eq:final-chain}.  Since the division
is by one common positive number, it also preserves the cross-prior bridge.
\end{proof}

\paragraph{Entropy identification.}

\begin{lemma}[Entropy reduction and Faddeev identification]
\label{pt:lem:faddeevsketch}
Define the value of full revelation by
\[
 \mathcal H(q):=V_q(\chi_q),
\]
so $\mathcal H(q)=H_G(q)/C$, with boundary priors read on their support.
Then every finite channel
$P:A\to\Delta(X)$ satisfies
\begin{equation}
 V_q(\mu_{qP})
 =\mathcal H(q)-\sum_{x:m_{qP}(x)>0}
    m_{qP}(x)\mathcal H(r_x^{\,q,P}).
\tag{ER}\label{pt:eq:entropy-reduction}
\end{equation}
Moreover,
\[
 \mathcal H(q)=\alpha\Sh(q)
\]
for a single $\alpha>0$.  Consequently
\[
 V_q(\mu_{qP})=\alpha I_{qP}(A;X).
\tag{MI}\label{pt:eq:value-MI}
\]
\end{lemma}

\begin{proof}
\emph{Entropy reduction.}
Assume first that $q$ has full support.  Append full revelation after $P$.
The final posterior law is $\chi_q$, while the continuation after record $x$
has full-revelation law $\chi_{r_x}$ on the support of $r_x$.  The exact
chain rule~\eqref{pt:eq:final-chain} therefore gives
\[
 \mathcal H(q)=V_q(\mu_{qP})
 +\sum_{x:m_{qP}(x)>0}m_{qP}(x)\mathcal H(r_x),
\]
which is~\eqref{pt:eq:entropy-reduction}.  A boundary prior is reduced to its
support; null action rows and null records affect neither side.

\emph{Strong additivity.}
Let $\mathcal P=\{A_1,\ldots,A_m\}$ be a partition of the support of $q$,
and put
\[
 p_i=q(A_i),\qquad q^i=q(\,\cdot\mid A_i),\qquad
 p=(p_1,\ldots,p_m).
\]
The channel $C_{\mathcal P}$ which reports the cell is neutral to the
corresponding full revelation on the cell index:
\[
 (q,C_{\mathcal P})\sim(p,\Id_{\{1,\ldots,m\}})
\]
by Lemma~\ref{pt:lem:neutrality}.  The common-scale cross-prior bridge from
Lemma~\ref{pt:lem:scalecoherence} therefore gives
\[
 V_q(\mu_{q,C_{\mathcal P}})=\mathcal H(p).
\]
Applying~\eqref{pt:eq:entropy-reduction} to $C_{\mathcal P}$ yields
\begin{equation}
 \mathcal H(q)=\mathcal H(p)+\sum_i p_i\mathcal H(q^i).
\tag{FA}\label{pt:eq:faddeev-recursion}
\end{equation}
Thus $\mathcal H$ is strongly additive.  Relabelling invariance of $V$ makes
$\mathcal H$ symmetric; support restriction makes it expansive under adding
zero-probability states; and $\mathcal H(\delta_a)=0$.  Record processing
makes every experiment weakly preferable to its no-information garbling, so
$V_q(\mu)\ge0$ on every fibre and hence $\mathcal H(q)\ge0$.
Axiom~\textup{(A4)} gives strict
positivity at every nondegenerate full-support prior.
Although \eqref{pt:eq:faddeev-recursion} was derived by partitioning the
positive support, it also holds with zero outer weights: delete the null cells
by support restriction, apply the displayed recursion, and reinsert them by
expansibility.

\emph{Continuity: the binary case.}
Write $h(t):=\mathcal H(t,1-t)$.  We derive, rather than assume, its
continuity from Axiom~\textup{(A2)}.  Fix $c\ge0$.  Choose a full-support
binary prior $\theta$ with $\mathcal H(\theta)>0$, as supplied by
Axiom~\textup{(A4)}.  Choose $n\ge1$ and $\lambda\in[0,1]$ so that
\[
 c=\lambda n\mathcal H(\theta).
\]
For $\theta_n:=\theta^{\otimes n}$, product additivity
\eqref{pt:eq:PArecovered} gives
$\mathcal H(\theta_n)=n\mathcal H(\theta)$.  Let $E_{n,\lambda}$ reveal the
$n$-bit action fully with probability $\lambda$ and report a common erasure
symbol otherwise.  Its posterior law is
$\lambda\chi_{\theta_n}+(1-\lambda)\delta_{\theta_n}$, so affinity gives
\[
 V_{\theta_n}(\mu_{\theta_n,E_{n,\lambda}})=c.
\]

Now fix the block environment
\[
 K_c:=\Id_{\{0,1\}}\sqcup E_{n,\lambda}.
\]
Let $q_t^0$ put the binary prior $(t,1-t)$ on its first block and let
$\theta_n^1$ put $\theta_n$ on its second block.  The cross-prior bridge says
\begin{align*}
 h(t)\ge c&\quad\Longleftrightarrow\quad q_t^0\succeq_{K_c}\theta_n^1,\\
 h(t)\le c&\quad\Longleftrightarrow\quad \theta_n^1\succeq_{K_c}q_t^0.
\end{align*}
The environment is fixed and $t\mapsto q_t^0$ is continuous, so
Axiom~\textup{(A2)} makes both level sets closed.  When $c<0$, the superlevel
set is all of $[0,1]$ and the sublevel set is empty because $h\ge0$.
Thus every upper and lower level set of $h$ is closed, and $h$ is continuous.

\emph{Continuity on every finite simplex.}
Proceed by induction on the number of coordinates.  For
$q=(q_1,(1-q_1)\bar q)$ with $q_1<1$, recursion
\eqref{pt:eq:faddeev-recursion} for the partition
$\{\{1\},\{2,\ldots,n\}\}$ gives
\begin{equation}
 \mathcal H(q)=h(q_1)+(1-q_1)\mathcal H(\bar q).
\label{pt:eq:continuity-induction}
\end{equation}
The right side is continuous away from $q_1=1$ by binary continuity and the
induction hypothesis.  At $q_1=0$, support restriction and $h(0)=0$ give
the same formula on the face.  At $q_1=1$, the induction hypothesis makes
$\mathcal H$ bounded on the compact $(n-1)$-simplex, so the second term in
\eqref{pt:eq:continuity-induction} tends to zero and the first tends to
$h(1)=0$.  Relabelling handles the other faces.  Hence $\mathcal H$ is
continuous on each closed finite simplex.

We have now verified all hypotheses of the strong-additivity form of
Faddeev's theorem: continuity on each finite simplex, invariance under
bijections, expansibility, vanishing on point masses, and
\eqref{pt:eq:faddeev-recursion}.  The finite entropy characterisation
\cite[Theorem~6]{BaezFritzLeinster2011} therefore gives
\[
 \mathcal H(q)=\alpha\Sh(q)
\]
for some $\alpha\ge0$.  Axiom~\textup{(A4)} forces $\alpha>0$.
Substitution in~\eqref{pt:eq:entropy-reduction} gives
\[
 V_q(\mu_{qP})
 =\alpha\left(\Sh(q)-\sum_x m_{qP}(x)
                    \Sh(r_x^{\,q,P})\right)
 =\alpha I_{qP}(A;X),
\]
where the sum is over positive-mass records.  This proves
\eqref{pt:eq:value-MI}.
\end{proof}

\begin{proof}[Completion of Proposition~\ref{pt:prop:main}]
After Lemma~\ref{pt:lem:scalecoherence}, the cross-prior bridge is represented
by $V$.  Lemma~\ref{pt:lem:faddeevsketch} identifies this value as
$\alpha I$ with one $\alpha>0$.  Hence, for arbitrary pairs $(q,P)$ and
$(r,Q)$,
\[
 q^0\succeq_{P\sqcup Q}r^1
 \quad\Longleftrightarrow\quad
 I_{qP}\ge I_{r,Q},
\tag{BMI}\label{pt:eq:final-block-MI}
\]
including boundary priors by support restriction.  For one fixed channel,
the duplication clause of Axiom~\textup{(A5)} identifies
$q\succeq_Pq'$ with the left side of~\eqref{pt:eq:final-block-MI} for
$Q=P$, proving the first assertion.  For the block clause,
Axiom~\textup{(A5)} deletes all irrelevant blocks and
\eqref{pt:eq:final-block-MI} compares the two remaining pairs.  A lottery
supported on block $i$ has a constant block tag, so
\[
 I_{q_i^i,\,\bigsqcup_{k\in J}P_k}
 =I_{q_iP_i},
\]
and similarly for block $j$.  This is exactly the stated common-scale block
comparison.
\end{proof}

\clearpage
\input{appendix_b_main_theorem_v5}
\clearpage
\section{Auxiliary results}\label{app:auxiliary-results}

This appendix collects proofs of the separability, uniqueness, sign, independence, and
choice-implication results stated in the main text.  None is used to prove
Theorem~\ref{thm:main}; separating them keeps Appendices A and B focused on
the representation argument.

\subsection{Separability, uniqueness, and sign}

For $q\in\D(A)$ and $K:A\to\D(O\times R)$, write $\pi^{qK}_{or}(a):=q(a)K(o,r\mid a)/p_{qK}(o,r)$ when
$p_{qK}(o,r)>0$ for the posterior induced by the visible pair, and
$\mu_{qK}:=\sum_{(o,r):p_{qK}(o,r)>0}p_{qK}(o,r)\,\delta_{\pi^{qK}_{or}}$ for its distribution.

\begin{corollary}[No payoff--trace complementarity]\label{cor:matching}
Suppose Axioms \textup{(A1)--(A8)} hold.  Fix $A$ and $q\in\D(A)$.  If channels $K$ and $L$ satisfy
$\E_{qK}[u(O)]=\E_{q,L}[u(O)]$ and $\mu_{qK}=\mu_{qL}$, then $(q,K)\sim(q,L)$.
\end{corollary}

\begin{proof}
The result follows from Theorem~\ref{thm:main} and
$\I_{qK}(A;O,R)=\Sh(q)-\int\Sh(\pi)d\mu_{qK}(\pi)$.
\end{proof}

\begin{proof}[Proof of Proposition~\ref{prop:uniqueness}]
Choose $o_*\in O$ with $u(o_*)=\underline u$.  First note that the range of $V$ is
$[\underline u,\infty)$.  The lower bound is immediate.  For $q\in\D(A)$ and $t\in[0,1]$, let a constant-$o_*$
record channel reveal the realised action with
probability $t$ and otherwise report $*$:
\[
 E_t(o_*,a'\mid a)=t\1_{a'=a},
 \qquad
 E_t(o_*,*\mid a)=1-t.
\]
Then $V(q,E_t)=\underline u+\lambda t\Sh(q)$.  At the uniform lottery on $n$ actions these values fill
$[\underline u,\underline u+\lambda\log n]$; the intervals are unbounded.

For part \textup{(i)}, let $W$ be a common-scale representation and define
$\varphi(v):=W(q,K)$ for any pair with $V(q,K)=v$.  This is well defined: equality of the two $V$-values gives
indifference in their two-block environment by Theorem~\ref{thm:main}, and the common-scale property gives equality
of the two $W$-values.  A strict inequality between $V$-values similarly gives a strict inequality between
$W$-values.  Thus $\varphi$ is strictly increasing and $W=\varphi\circ V$.  The converse is immediate from
Theorem~\ref{thm:main}.

For part \textup{(ii)}, the expected-utility and mutual-information chain rules give, for every
$P:A\to\D(\{1,\ldots,k\})$ and every
$K_1,\ldots,K_k,L_1,\ldots,L_k:A\to\D(O\times R)$, writing
$m(i):=\sum_aq(a)P(i\mid a)$,
\[
 V\bigl(q,P\triangleright(K_1,\ldots,K_k)\bigr)
 -V\bigl(q,P\triangleright(L_1,\ldots,L_k)\bigr)
 =\sum_{i:m(i)>0}m(i)
 \bigl[V(q_i,K_i)-V(q_i,L_i)\bigr].
\]
Hence every $\alpha V+\beta$, with $\alpha>0$, is common-scale and branch-additive.

Conversely, let the branch-additive common-scale representation be $W=\varphi\circ V$ and define
\[
 \psi(x):=\varphi(\underline u+x)-\varphi(\underline u),
 \qquad x\ge0.
\]
Fix $n\ge2$ and let $q$ be uniform on $n$ actions.  For $x\in[0,\lambda\log n]$, choose a constant-$o_*$ erasure
channel whose value is $\underline u+x$.  In the first compound, use this channel after the first branch
of a public coin with probabilities $\theta$ and $1-\theta$, and use the uninformative constant-$o_*$ channel after
the second.  In the comparison compound, use the uninformative channel after both branches.  The two compound values
are $\underline u+\theta x$ and $\underline u$, so \eqref{eq:branch-additive-index} gives
\[
 \psi(\theta x)=\theta\psi(x)
 \qquad(0\le\theta\le1).
\]
Taking $x=\lambda\log n$ shows that $\psi$ is linear on $[0,\lambda\log n]$.  These intervals are nested and
unbounded, so $\psi(x)=\alpha x$ on $[0,\infty)$ for one $\alpha>0$.  Therefore
$\varphi(v)=\alpha v+\beta$ on the range of $V$.

Finally, any second trace-tempered index $\widetilde V$ is common-scale, because expected utility and mutual
information are invariant under block embedding, and is branch-additive by the same two chain rules.  Applying
part \textup{(ii)} and then restricting to deterministic payoff outcomes and to constant-payoff record channels gives
$\widetilde u=\alpha u+\beta$ and $\widetilde\lambda=\alpha\lambda$.
\end{proof}

\begin{proof}[Proof of Corollary~\ref{cor:signduality}]
Part \textup{(i)} is Theorem~\ref{thm:main}.  For part \textup{(ii)}, reverse
the primitive relation in every environment.  The reversed family satisfies
Axioms~\textup{(A1)}, \textup{(A2)}, \textup{(A5)}, and \textup{(A8)}; Axiom~\textup{(A3)} holds with
its two witnessing outcomes interchanged.  It also satisfies the forward
versions of Axioms~\textup{(A4)}, \textup{(A6)}, and
\textup{(A7)}.  Theorem~\ref{thm:main} therefore represents the reversed
family by
\(
  \E[\widehat u(O)]+\widehat\lambda I(A;O,R)
\)
with \(\widehat\lambda>0\).  Negating that index represents the original
family, so it has the asserted form with
\(u=-\widehat u\) and \(\lambda=-\widehat\lambda<0\).

For part \textup{(iii)}, collapse the action alphabet of any pair $(q,K)$ to
a singleton.  Axiom~$(\mathrm{A}7^0)$ makes the pair indifferent to the
resulting one-action channel, whose output law is $p_{qK}$.  Axiom
$(\mathrm{A}6^0)$ then erases its record without changing value.  Thus every
pair is indifferent to the one-action channel that delivers precisely its
payoff marginal.  Axiom~\textup{(A5)} makes these reductions coherent across
blocks.  On payoff lotteries, Axiom~\textup{(A8)} applied to an
action-independent public coin gives mixture independence, Axiom
\textup{(A2)} gives continuity, and Axiom~\textup{(A3)} gives
non-triviality.  The mixture-space theorem therefore yields a nonconstant
expected-utility index.  Conversely, expected utility is unaffected by record
or action processing and satisfies the three indifference clauses.  The
positive and negative formulae satisfy their respective clauses by data
processing and the two chain rules.
\end{proof}

\subsection{Independence of the axioms}

\begin{proof}[Proof of Proposition~\ref{prop:independence}]
Fix a nonconstant $u$ and write $U(q,K):=\E_{qK}[u(O)]$, $Z:=(O,R)$, and $I:=\I(A;Z)$.  All scores below are
extended to block and compound channels by the same formula.

\emph{Failure of Axiom \textup{(A1)}.}
Use the Pareto relation on the two coordinates $(U,I)$:
\[
 (q,K)\succeq(p,L)
 \quad\Longleftrightarrow\quad
 U(q,K)\ge U(p,L)\ \text{and}\ I(q,K)\ge I(p,L).
\]
It is transitive but incomplete whenever one alternative has more payoff and the other more information.  Its graph is
closed.  Both data-processing axioms are coordinatewise monotonic.  Under a common-branch replacement, both the
payoff and information differences are multiplied by the same positive branch probability, so Axiom~\textup{(A8)}
holds in both directions.  Block coherence holds.  For Axiom~\textup{(A3)}, choose the two outcomes so that
$u(o^+)>u(o^-)$; for Axiom~\textup{(A4)}, the two test lotteries have equal payoff value and information values
$\log2$ and zero.  Thus both relevance axioms hold.

\emph{Failure of Axiom \textup{(A2)}.}
Order alternatives lexicographically by $(U,I)$, with payoff first.  This is a weak order and satisfies all structural,
processing, and continuation requirements.  Axiom~\textup{(A3)} holds through the first coordinate, and
Axiom~\textup{(A4)} through the second when the first is tied.  It is not continuous even within one fixed channel.  Let a
three-action channel reveal the action fully and let only action $2$ pay the high outcome.  Put
\[
 p=(1/4,1/2,1/4),\qquad q=(0.49,0.50,0.01),\qquad
 q_n=(0.49-\varepsilon_n,0.50+\varepsilon_n,0.01),
\]
where $\varepsilon_n\downarrow0$.  The lotteries $p$ and $q$ have the same expected utility but
$\Sh(q)<\Sh(p)$, so $p\succ_Kq$.  Every $q_n$ has strictly higher expected utility than $p$, so
$q_n\succ_Kp$.  Thus the graph in Axiom~\textup{(A2)} is not closed.

\emph{Failure of Axiom \textup{(A3)}.}
Use $V:=I$.  All other axioms hold.  In particular, the test in
Axiom~\textup{(A4)} has information values $\log2$ and zero.  But the two pure
lotteries in Axiom~\textup{(A3)} both have zero information, whatever their payoffs.

\emph{Failure of Axiom \textup{(A4)}.}
Use $V:=U$.  All other axioms hold.  The outcomes in Axiom~\textup{(A3)} can be
chosen to have different utility, but the two lotteries in Axiom~\textup{(A4)}
always have the same payoff and hence are indifferent.

\emph{Failure of Axiom \textup{(A5)}.}
Treat outer block labels as syntactic action tags.  For an action alphabet presented with at least three outer
blocks, put $w(a^0)=1$ and $w(a^i)=0$ for $i\ne0$; put $w(a)=0$ on every action alphabet with fewer than three
outer blocks.  For $\varepsilon>0$, use within each environment
\[
 V_\varepsilon(q,K):=U(q,K)+I(q,K)+\varepsilon\sum_aq(a)w(a).
\]
This defines a continuous weak order in every fixed environment.  Every comparison in Axioms~\textup{(A6)},
\textup{(A7)}, and \textup{(A8)} is made in a two-block environment, where $w=0$; those axioms therefore reduce
to their valid $U+I$ comparisons.  The fixed relevance tests can be chosen on untagged action sets and also reduce
to $U+I$.  But irrelevant-block coherence fails.  Take three identical one-action, constant-payoff, silent
channels.  Their first two block-supported lotteries are indifferent in the canonical two-block comparison, while
in the three-block environment the lottery supported on block $0$ is strictly preferred to the lottery supported
on block $1$ by $\varepsilon$.

\emph{Failure of Axiom \textup{(A6)}.}
For $\varepsilon>0$, use
\begin{equation}
 V(q,K):=U(q,K)+I(A;Z)-\varepsilon\Sh(Z\mid A).
 \label{eq:no-record-dp}
\end{equation}
This score is continuous and branch-additive, since both mutual information and conditional entropy obey the chain
rule.  Under action processing $A\to B$, the score difference between the original and processed alternatives is
\[
 \bigl[I(A;Z)-I(B;Z)\bigr]
 +\varepsilon\bigl[\Sh(Z\mid B)-\Sh(Z\mid A)\bigr]
 =(1+\varepsilon)I(A;Z\mid B)\ge0,
\]
so \eqref{eq:no-record-dp} satisfies Axiom~\textup{(A7)}.  Axiom~\textup{(A3)} holds through $U$.  In the
Axiom~\textup{(A4)} test, the first lottery has $I=\log2$ and $\Sh(Z\mid A)=0$, while the second has
$I=0$ and $\Sh(Z\mid A)=\log2$; hence the first is strictly preferred.  Record data
processing fails: with constant payoff and a fair record independent of $A$, the original value is $-\varepsilon\log2$;
deleting that noise gives value zero.

\emph{Failure of Axiom \textup{(A7)}.}
Let $G(q):=1-\sum_aq(a)^2$ be Gini entropy and define its posterior reduction
\[
 J_G(q,K):=G(q)-\int G(\pi)\,d\mu_{qK}(\pi),
 \qquad V:=U+J_G.
\]
Concavity of $G$ gives record data processing, and the potential-difference form gives the exact branch chain rule.
Axiom~\textup{(A3)} holds through $U$.  In the Axiom~\textup{(A4)} test, $J_G$ is $1/2$ under the first
lottery and zero under the second, so trace relevance holds.  Action processing can nevertheless increase $J_G$.  Take
$q=(1/8,1/8,3/4)$, a binary record with
\[
 \Pr(0\mid a_1)=\Pr(0\mid a_2)=0,
 \qquad \Pr(0\mid a_3)=1/4,
\]
and merge $a_1,a_2$ into one processed action.  Direct calculation gives $J_G=9/416$ before the merge and
$J_G=3/104$ afterwards.  Hence Axiom~\textup{(A7)} fails.

\emph{Failure of Axiom \textup{(A8)}.}
Use $V:=U+\min\{I,1\}$, with information measured in bits.  It is continuous, respects both data-processing axioms,
and satisfies block coherence.  Axiom~\textup{(A3)} holds through $U$, while the test in Axiom~\textup{(A4)}
has information values one and zero.  Let $q$ be uniform on four actions.  A first-stage record reveals which of two pairs
contains the action.  Hold the continuation following one record fixed and silent.  Following the other record,
keep payoffs constant and compare a silent continuation with one that reveals the action within its pair.  The latter continuation is strictly
preferred at the branch posterior, but the two compounds have respectively one and three-halves bits before capping
and are therefore indifferent.  The forward implication in Axiom~\textup{(A8)} fails.
\end{proof}

\subsection{Proofs of the choice implications}\label{app:choice-proofs}

Several results stated only in summary form in the text are given here in full.  Proofs appear in an order that
front-loads shared machinery.  Throughout this subsection, write
$Z:=O\times R$ and $K_a:=K(\cdot\mid a)\in\D(Z)$ for the row of action $a$; the distribution induced by $q$ is
$p_{qK}$, as in Section~\ref{sec:theorem}.

\begin{proof}[Proof of Proposition~\ref{prop:signature}]
(i) Under either lottery, each visible pair $(o_0,r_j)$, $j=1,2$, has probability $\tfrac12$: under $q'$ because
each $a_j$ produces $(o_0,r_j)$ for sure, under $q''$ because each of $a_3,a_4$ produces either pair with equal
probability.  Hence $p_{q'K}=p_{q''K}$.  On $\supp(q')$ the rows are degenerate and distinct, so the visible pair
reveals the action and $\I_{q'K}=\Sh(q')=\log2$; on $\supp(q'')$ the rows are identical, so the visible pair is
independent of the action and $\I_{q''K}=0$.  The payoff terms agree because the payoff marginals do, giving the
displayed gap.
(ii) All rows of $K^\circ$ are equal, so $\I_{qK^\circ}=0$ for every $q$; since the payoff is constantly $o_0$,
the representation gives indifference.
\end{proof}

\begin{proof}[Proof of Proposition~\ref{prop:optimal-choice}]
Expected utility is linear in $q$, while mutual information is concave in the input distribution of a fixed channel
\cite{CoverThomas2006}.  Continuity on the compact simplex gives existence.  Using
\[
 I_{qK}(A;Z)=\sum_aq(a)\KL\!\left(K_a\,\middle\|\,p_{qK}\right),
\]
its directional derivative in the mass on action $a$ equals the displayed divergence up to a term common to every
action.  The result is the Kuhn--Tucker condition for a concave programme.  At constant $u$ it reduces to the classical
capacity condition.
\end{proof}

\begin{proposition}[Optimal marginal and multiplicity]\label{prop:optimal-marginal}
Let $\lambda>0$ and write $C^*(K):=\max_{q\in\D(A)}V(q,K)$.  There is a unique $p^*\in\D(Z)$ such that
$p_{q^*K}=p^*$ for every optimal lottery $q^*$.  With
\[
 \mathcal A^*
 :=\bigl\{a\in A:
 \bar u(a)+\lambda
 \KL\!\left(K_a\,\middle\|\,p^*\right)=C^*(K)
 \bigr\},
\]
the set of optimal lotteries is exactly
\[
 \bigl\{q\in\D(A):
 p_{qK}=p^*,\ \supp(q)\subseteq\mathcal A^*
 \bigr\}.
\]
The optimum is unique whenever the rows $\{K_a:a\in\mathcal A^*\}$ are affinely independent.  If actions are
partitioned into classes with identical rows, $V(q,K)$ depends on $q$ within each class only through the class's
total mass.  The optimal lotteries are exactly the lifts of the optimal lotteries of the collapsed channel, and
each class's mass may be distributed arbitrarily among its members.
\end{proposition}

\begin{proof}[Proof of Proposition~\ref{prop:optimal-marginal}]
The entropy identity gives
\[
 V(q,K)
 =\lambda\Sh(p_{qK})
  +\sum_aq(a)\bigl[\bar u(a)-\lambda\Sh(K_a)\bigr].
\]
If $q_0$ and $q_1$ are optimal with $p_{q_0K}\ne p_{q_1K}$, strict concavity of Shannon entropy and linearity of
$q\mapsto p_{qK}$ and of the remaining sum give $V(tq_0+(1-t)q_1,K)>C^*(K)$ for $t\in(0,1)$, a contradiction.
Hence every optimum induces the same $p^*$.  By Proposition~\ref{prop:optimal-choice}, an optimum $q^*$ satisfies
the displayed equalities at some constant $c$; averaging them under $q^*$ gives $c=V(q^*,K)=C^*(K)$, so
$\supp(q^*)\subseteq\mathcal A^*$.  Conversely, if $p_{qK}=p^*$ and $\supp(q)\subseteq\mathcal A^*$, then
\[
 V(q,K)
 =\sum_aq(a)\bigl[\bar u(a)+\lambda\KL(K_a\|p^*)\bigr]
 =C^*(K).
\]
Affine independence makes the representation $p^*=\sum_aq(a)K_a$ with $\supp(q)\subseteq\mathcal A^*$ unique.
Identical rows share $\bar u(a)$ and $\Sh(K_a)$, so both $p_{qK}$ and the displayed expression for $V$ are
unchanged by redistribution within a class; collapsing the classes and lifting aggregate masses proves the last
claim.
\end{proof}

\begin{proposition}[Support and pure choice]\label{cor:support-pure}
Let $\lambda>0$ and let $Z_K:=\{z\in Z:K_a(z)>0\text{ for some }a\in A\}$.
\begin{enumerate}[label=\textup{(\roman*)},leftmargin=*]
\item Every optimal lottery $q^*$ satisfies $p_{q^*K}(z)>0$ for every $z\in Z_K$.  Consequently, if every action
$a$ has a private consequence---some $z_a$ with $K_a(z_a)>0$ and $K_b(z_a)=0$ for all $b\ne a$---then every
optimal lottery has full support.
\item The pure lottery $\delta_a$ is optimal if and only if
\[
 \bar u(b)+\lambda\KL\!\left(K_b\,\middle\|\,K_a\right)
 \le \bar u(a)
 \qquad\text{for every }b\in A,
\]
and uniquely optimal if these inequalities are strict for $b\ne a$.  In particular, if some rival reaches a
consequence outside the support of $K_a$, then $\delta_a$ is not optimal, however large its payoff advantage.
\end{enumerate}
\end{proposition}

\begin{proof}[Proof of Proposition~\ref{cor:support-pure}]
Throughout, $p^*$ denotes the common optimal marginal of Proposition~\ref{prop:optimal-marginal}.
(i) Suppose $p^*(z)=0$ for some $z\in Z_K$ and choose $a$ with $K_a(z)>0$.  Then
$\KL(K_a\|p^*)=+\infty$, which is incompatible with the finite on-support equality if some optimum uses $a$ and
with the off-support inequality of Proposition~\ref{prop:optimal-choice} if none does.  Under the
private-consequence condition, an optimum with $q^*(a)=0$ would force $p^*(z_a)=0$.

(ii) For $q=\delta_a$ the induced marginal is $K_a$ and the on-support constant is $\bar u(a)$; the displayed
system is the off-support condition of Proposition~\ref{prop:optimal-choice}, necessary and sufficient by
concavity.  Under strict inequalities, $\mathcal A^*=\{a\}$ in Proposition~\ref{prop:optimal-marginal}, so every
optimum is $\delta_a$.
\end{proof}

\begin{proposition}[Optimal mixing]\label{prop:binary}
Let $\lambda>0$ and $A=\{a,b\}$ with $K_a\ne K_b$ and
$\Delta:=\bar u(a)-\bar u(b)\ge0$.  For $x\in[0,1]$ let $q_x$ put mass $x$ on $b$, so that
$p_{q_xK}=(1-x)K_a+xK_b=:p_x$, and for $x\in(0,1)$ define
\[
 g(x):=\KL\!\left(K_b\,\middle\|\,p_x\right)-\KL\!\left(K_a\,\middle\|\,p_x\right).
\]
\begin{enumerate}[label=\textup{(\roman*)},leftmargin=*]
\item $g$ is continuous and strictly decreasing, with $g(0^+)=\KL(K_b\,\|\,K_a)$ and
$g(1^-)=-\KL(K_a\,\|\,K_b)$ as extended-real limits.  The pure lottery $\delta_a$ is optimal if and only if
$\Delta\ge\lambda\,\KL(K_b\,\|\,K_a)$.  Below this threshold the optimum is unique and interior,
$q^*=q_{x^*}$ with $x^*$ the unique solution of $\Delta=\lambda g(x)$.
\item The optimal mass $x^*(\Delta,\lambda)$ on the inferior action is nonincreasing in $\Delta$ and
nondecreasing in $\lambda$, strictly so on the mixing region; at $\Delta=0$ it is the capacity-achieving input
of the two-row channel, independently of $\lambda$.
\end{enumerate}
\end{proposition}

\begin{proof}
Write $f(x):=V(q_x,K)=\bar u(a)-x\Delta+\lambda I(x)$ with
$I(x)=\Sh(p_x)-(1-x)\Sh(K_a)-x\Sh(K_b)$.  Since $K_a\ne K_b$,
$x\mapsto p_x$ is affine and injective, so strict concavity of $\Sh$ \cite{CoverThomas2006} makes $f$ strictly
concave on $[0,1]$, and a unique maximiser exists.  On $(0,1)$ every $z\in\supp K_a\cup\supp K_b$ has $p_x(z)>0$
and, using $\sum_z(K_b(z)-K_a(z))=0$,
\[
 I'(x)=\Sh(K_a)-\Sh(K_b)-\sum_z\bigl(K_b(z)-K_a(z)\bigr)\log p_x(z)=g(x),
\]
so $f'(x)=-\Delta+\lambda g(x)$, with $g$ continuous and strictly decreasing as the derivative of a strictly
concave function.  As $x\downarrow0$, $p_x(z)\to K_a(z)$; if $\supp K_b\subseteq\supp K_a$ every summand converges
and $g(0^+)=\KL(K_b\|K_a)$, while if some $\bar z$ has $K_b(\bar z)>0=K_a(\bar z)$ the summand
$-K_b(\bar z)\log\bigl(xK_b(\bar z)\bigr)$ diverges to $+\infty$, matching the extended-sense divergence.
Exchanging the roles of $a$ and $b$ gives the limit at $1$.  By concavity and the mean value theorem,
$\delta_a$ is optimal if and only if $\lim_{x\downarrow0}f'(x)\le0$, i.e.\ $\Delta\ge\lambda\,\KL(K_b\|K_a)$;
this is the threshold stated in the text, and Proposition~\ref{cor:support-pure}\textup{(ii)} with two actions.
If $\Delta<\lambda\,\KL(K_b\|K_a)$, the maximiser is not $x=0$ by the threshold and not $x=1$
because $f'(1^-)=-\Delta-\lambda\KL(K_a\|K_b)<0$; it is therefore the unique root of $\Delta=\lambda g(x)$, which
exists because $g$ decreases strictly between its limits.  This proves \textup{(i)}.  The comparative statics of
\textup{(ii)} follow by inverting the
strictly decreasing $g$ at $\Delta/\lambda$: raising $\Delta$ raises the target, raising $\lambda$ lowers it; at
$\Delta=0$ the condition $g(x)=0$ is the capacity condition of Proposition~\ref{prop:optimal-choice} at constant
payoff.
\end{proof}

\begin{remark}[Overlapping traces and IIA]\label{rem:overlap-iia-example}
Let three actions have equal expected utility and rows
\[
 K_1=(1,0,0),\qquad K_2=(0,1,0),\qquad K_3=(1/2,0,1/2).
\]
With only the first two actions, the optimal lottery is $(1/2,1/2)$.  Adding the third gives
\[
 q^*=(1/3,4/9,2/9),\qquad p_{q^*K}=(4/9,4/9,1/9).
\]
Each active row has divergence $\log(9/4)$ from $p_{q^*K}$, so
Proposition~\ref{prop:optimal-choice} verifies optimality.  The ratio of the first two choice probabilities falls
from $1$ to $3/4$.  Thus overlapping rows can generate the red-bus/blue-bus pattern even when expected utilities
are equal.
\end{remark}

\begin{proposition}[Trace demand]\label{prop:trace-demand}
Fix $K$ and write $e(q):=\E_{qK}[u(O)]$, $i(q):=\I_{qK}(A;O,R)$,
\[
 V^*(\lambda):=\max_{q\in\D(A)}\{e(q)+\lambda i(q)\},
 \qquad
 Q(\lambda):=\arg\max_{q\in\D(A)}\{e(q)+\lambda i(q)\}.
\]
\begin{enumerate}[label=\textup{(\roman*)},leftmargin=*]
\item $V^*$ is convex on $(0,\infty)$, $i$ is affine on the convex compact set $Q(\lambda)$, and
\[
 (V^*)'_-(\lambda)=\min_{q\in Q(\lambda)}i(q),
 \qquad
 (V^*)'_+(\lambda)=\max_{q\in Q(\lambda)}i(q);
\]
at every point of differentiability, $(V^*)'(\lambda)=i(q)$ for every $q\in Q(\lambda)$.
\item Let $0<\lambda_1<\lambda_2$ and $q_j\in Q(\lambda_j)$, writing $e_j:=e(q_j)$, $i_j:=i(q_j)$.  Then
$i_1\le i_2$, $e_1\ge e_2$, and
\[
 \lambda_1(i_2-i_1)\le e_1-e_2\le\lambda_2(i_2-i_1).
\]
Hence either $(e_1,i_1)=(e_2,i_2)$, or $i_1<i_2$ and $e_1>e_2$ and
\[
 \lambda_1\le\frac{e_1-e_2}{i_2-i_1}\le\lambda_2.
\]
\end{enumerate}
\end{proposition}

\begin{proof}[Proof of Proposition~\ref{prop:trace-demand}]
(i) $V^*$ is a pointwise maximum of functions affine in $\lambda$, hence convex and, since
$0\le i\le\log|A|$, finite.  On the convex compact set $Q(\lambda)$ the objective is constant and $e$ is linear, so
$i=(V^*-e)/\lambda$ is affine there.  For $q\in Q(\lambda)$ and $\mu>0$,
$V^*(\mu)\ge V^*(\lambda)+(\mu-\lambda)i(q)$, so each such $i(q)$ is a subgradient.  For the right derivative,
choose $q_h\in Q(\lambda+h)$; then
\[
 \max_{q\in Q(\lambda)}i(q)
 \le\frac{V^*(\lambda+h)-V^*(\lambda)}{h}
 \le i(q_h),
\]
the first inequality by the subgradient bound, the second by optimality of $Q(\lambda)$ at $\lambda$.  By
compactness and continuity, accumulation points of $(q_h)_{h\downarrow0}$ lie in $Q(\lambda)$, so the right-hand
side accumulates in $i(Q(\lambda))$ and the right derivative equals the maximum.  The left derivative is symmetric,
and a convex function on an interval is differentiable off a countable set.

(ii) Optimality at the two weights gives
$e_1+\lambda_1i_1\ge e_2+\lambda_1i_2$ and $e_2+\lambda_2i_2\ge e_1+\lambda_2i_1$.  Adding yields
$(\lambda_2-\lambda_1)(i_2-i_1)\ge0$, hence $i_2\ge i_1$; the two inequalities rearrange to the displayed bracket,
which gives $e_1\ge e_2$ and, when the pairs differ, the final bound.
\end{proof}

\begin{corollary}[Extreme trace weights]\label{cor:lambda-limits}
Let $e^0:=\max_qe(q)$, $i^0:=\max\{i(q):e(q)=e^0\}$, $C:=\max_qi(q)$, and $e^C:=\max\{e(q):i(q)=C\}$.  If
$\lambda_n\downarrow0$ and $q_n\in Q(\lambda_n)$, every accumulation point of $(q_n)$ maximises $i$ among the
maximisers of $e$; if $\lambda_n\uparrow\infty$, every accumulation point maximises $e$ among the maximisers of
$i$.  Moreover
\[
 V^*(\lambda)=e^0+\lambda i^0+o(\lambda)
 \quad\text{as }\lambda\downarrow0,
 \qquad
 V^*(\lambda)=\lambda C+e^C+o(1)
 \quad\text{as }\lambda\uparrow\infty.
\]
\end{corollary}

\begin{proof}[Proof of Corollary~\ref{cor:lambda-limits}]
For $\lambda_n\downarrow0$, comparing $q_n$ with an $e$-maximiser attaining $i^0$ gives
\[
 0\le e^0-e(q_n)\le\lambda_n\bigl(i(q_n)-i^0\bigr).
\]
Boundedness of $i$ forces $e(q_n)\to e^0$, and the same inequality gives $i(q_n)\ge i^0$; continuity delivers the
accumulation claim.  Dividing by $\lambda_n$,
\[
 \frac{V^*(\lambda_n)-e^0}{\lambda_n}=i(q_n)-\frac{e^0-e(q_n)}{\lambda_n},
\]
where the last fraction lies in $[0,\,i(q_n)-i^0]$ and vanishes, proving the first expansion.  For
$\lambda_n\uparrow\infty$, comparing $q_n$ with an $i$-maximiser attaining $e^C$ gives
$0\le\lambda_n(C-i(q_n))\le e(q_n)-e^C$, and the symmetric argument completes the proof.
\end{proof}

\begin{proof}[Proof of Proposition~\ref{prop:garbling-price}]
The processor leaves the payoff coordinate unchanged, so the expected-utility terms coincide and the difference is
$\lambda[\I_{qK}(A;O,R)-\I_{qL}(A;O,R')]$.  Write $Z:=(O,R)$ and $Z':=(O,R')$.  By construction $A\to Z\to Z'$ is
a Markov chain, so expanding $\I(A;Z,Z')$ in both orders gives
\[
 \I(A;Z)-\I(A;Z')=\I(A;Z\mid Z')=\I(A;R\mid O,R'),
\]
the last equality because $O$ is part of $Z'$.  Conditional mutual information is nonnegative and vanishes exactly
when $A$ and $(O,R)$ are conditionally independent given $(O,R')$ \cite{CoverThomas2006}.  This is the zero-price
condition stated in the proposition.  Under a garbling, it is also equivalent to equality of the posterior laws,
$\mu_{qK}=\mu_{qL}$.  Indeed, the Markov property gives
$\Pr(A=\cdot\mid Z,Z')=\pi^{qK}_{Z}$ almost surely, while conditional independence gives
$\Pr(A=\cdot\mid Z,Z')=\pi^{qL}_{Z'}$; the two random posteriors then coincide almost surely.  Conversely, by the
identity in the proof of Corollary~\ref{cor:matching}, equality of the posterior laws gives equality of the
information terms, and the payoff terms already agree.
\end{proof}

\begin{corollary}[The recorded coin]\label{cor:recorded-coin}
Let $K,L:A\to\D(O\times R)$ and $t\in(0,1)$.  Let $tK\oplus(1-t)L$ be the public mixture that draws an
action-independent coin $Y$ with $\Pr(Y=0)=t$, applies $K$ if $Y=0$ and $L$ otherwise, and includes the coin in
the record; let $M_t:=tK+(1-t)L$ be the rowwise mixture, which hides it.  Then, for every $q\in\D(A)$,
\begin{enumerate}[label=\textup{(\roman*)},leftmargin=*]
\item $V\bigl(q,\,tK\oplus(1-t)L\bigr)=t\,V(q,K)+(1-t)\,V(q,L)$;
\item $V\bigl(q,\,tK\oplus(1-t)L\bigr)-V(q,M_t)=\lambda\,\I(A;Y\mid O,R)\ge0$, with equality if and only if $A$
and $Y$ are conditionally independent given $(O,R)$.
\end{enumerate}
\end{corollary}

\begin{proof}[Proof of Corollary~\ref{cor:recorded-coin}]
Part \textup{(i)} is the chain rule with an action-independent first stage:
$\I(A;Y)=0$, and each reached branch reproduces the joint law of the
corresponding constituent.  Part \textup{(ii)} follows from
Proposition~\ref{prop:garbling-price}: deleting the coin coordinate is a
deterministic payoff-preserving record processor carrying the public mixture
to $M_t$.
\end{proof}

\begin{definition}[Garbling and more-capable orders]\label{def:channel-orders}
Let $K:A\to\D(O\times R)$ and $L:A\to\D(O\times R')$ have the same payoff kernels.  Write $K\succeq_gL$ if
$L=KT$ for some record processor $T$, and write $K\succeq_cL$ if
$\I_{qK}(A;O,R)\ge\I_{qL}(A;O,R')$ for every $q\in\D(A)$; the latter is the \emph{more-capable} order on the
visible-consequence channels \cite{KornerMarton1977}.
\end{definition}

\begin{proposition}[Uniform dominance is more-capable, not garbling]\label{prop:blackwell}
Let $K$ and $L$ have identical payoff kernels.
\begin{enumerate}[label=\textup{(\roman*)},leftmargin=*]
\item $K\succeq_cL$ if and only if $V(q,K)\ge V(q,L)$ for every $q\in\D(A)$.
\item $K\succeq_gL\Rightarrow K\succeq_cL$, but the converse fails.  There are channels for which
$K\succeq_cL$ although $L\ne KT$ for every record processor $T$.
\end{enumerate}
\end{proposition}

\begin{proof}[Proof of Proposition~\ref{prop:blackwell}]
(i) Because the payoff kernels coincide,
$V(q,K)-V(q,L)=\lambda[\I_{qK}-\I_{qL}]$ with $\lambda>0$.  Uniform value dominance and
$K\succeq_cL$ are therefore the same statement.

(ii) If $L=KT$, Proposition~\ref{prop:garbling-price} gives $\I_{qK}\ge\I_{qL}$ at every $q$, and hence
$K\succeq_gL\Rightarrow K\succeq_cL$.  For failure of the converse, let $A=\{0,1\}$, fix $o^\ast\in O$, and let
both channels pay $o^\ast$ surely.  Let $K$ carry the binary
erasure record with erasure probability $\varepsilon=\tfrac12$ and $L$ the binary symmetric record with crossover
$p=\tfrac15$:
\[
 K(o^\ast,a\mid a)=1-\varepsilon,\quad K(o^\ast,\mathrm e\mid a)=\varepsilon,
 \qquad
 L(o^\ast,a\mid a)=1-p,\quad L(o^\ast,1-a\mid a)=p.
\]
The payoff kernels agree row by row, so it suffices to compare informations, and the comparison is invariant to the
base; compute in bits with $h$ the binary entropy.  For $q=(1-t,t)$, grouping records gives
\[
 \I_{qK}=(1-\varepsilon)\,h(t),
 \qquad
 \I_{qL}=h(p\ast t)-h(p),
 \qquad p\ast t:=p(1-t)+(1-p)t.
\]
By Mrs Gerber's Lemma \cite{WynerZiv1973}, $v\mapsto h\bigl(p\ast h^{-1}(v)\bigr)$ is convex on $[0,1]$, with
$h^{-1}$ valued in $[0,\tfrac12]$; since $h(p\ast t)$ is invariant under $t\mapsto1-t$, take $t\le\tfrac12$ and
$v=h(t)$.  The chord between $v=0$ and $v=1$ gives
\[
 h(p\ast t)\le(1-v)\,h(p)+v,
 \qquad\text{hence}\qquad
 \I_{qL}\le\bigl(1-h(p)\bigr)h(t)\le(1-\varepsilon)\,h(t)=\I_{qK},
\]
the last step because $\varepsilon=\tfrac12\le h(\tfrac15)$.  Thus $K\succeq_cL$.  Now suppose $L=KT$ for a
processor $T$, and write $t_r:=T(0\mid o^\ast,r)$ for $r\in\{0,1,\mathrm e\}$.  Matching the probability of the
processed record $0$ in the two rows,
\[
 (1-\varepsilon)t_0+\varepsilon t_{\mathrm e}=1-p,
 \qquad
 (1-\varepsilon)t_1+\varepsilon t_{\mathrm e}=p,
\]
and subtracting gives $t_0-t_1=(1-2p)/(1-\varepsilon)=\tfrac65>1$, impossible.  Compositions of record processors
are record processors, so no chain of garblings produces $L$ from $K$ either.
\end{proof}

The next definition and proposition formalise the side-bet elicitation of \S\ref{sec:measurement}.  Let $o^+$
and $o^-$ attain the maximum and minimum of $u$, write $\Delta u:=u(o^+)-u(o^-)>0$, and for $\beta\in[0,1]$ let
$w(\beta):=\beta u(o^+)+(1-\beta)u(o^-)$.

\begin{definition}[Side bet]\label{def:dilution}
For $K:A\to\D(O\times R)$, put $R^\ast:=R\sqcup\{\ast\}$ and regard $K$ as a channel on $R^\ast$ by zero
padding.  For $\varepsilon\in[0,1)$ and $\beta\in[0,1]$, let $P_\varepsilon:A\to\D(\{1,2\})$ be the public coin
$P_\varepsilon(2\mid a)=\varepsilon$, and let $B_\beta:A\to\D(O\times R^\ast)$ have the constant rows
$B_\beta(o^+,\ast\mid a)=\beta$ and $B_\beta(o^-,\ast\mid a)=1-\beta$.  Write
\[
 D_{\varepsilon,\beta}K:=P_\varepsilon\triangleright(K,B_\beta),
\]
as in \eqref{eq:compound}.
\end{definition}

\begin{proposition}[Willingness to pay for trace]\label{prop:wtp}
Fix $q\in\D(A)$, $p\in\D(B)$, and channels $K:A\to\D(O\times R)$ and $L:B\to\D(O\times S)$.
\begin{enumerate}[label=\textup{(\roman*)},leftmargin=*]
\item For every $\varepsilon\in[0,1)$ and $\beta\in[0,1]$,
\[
 V\bigl(q,\,D_{\varepsilon,\beta}K\bigr)=(1-\varepsilon)\,V(q,K)+\varepsilon\,w(\beta).
\]
\item If $\varepsilon\in(0,1)$ and $\beta,\beta'\in[0,1]$ satisfy
$\bigl(q,D_{\varepsilon,\beta}K\bigr)\sim\bigl(p,D_{\varepsilon,\beta'}L\bigr)$, then
\[
 \frac{\varepsilon}{1-\varepsilon}\bigl[w(\beta)-w(\beta')\bigr]
 =V(p,L)-V(q,K).
\]
When the material parts agree, $\E_{qK}[u(O)]=\E_{p,L}[u(O)]$, the left-hand side equals
$\lambda\bigl[\I_{pL}(B;O,S)-\I_{qK}(A;O,R)\bigr]$: information is priced linearly at rate $\lambda$.
\item Such a pair $(\beta,\beta')$ exists if and only if
$\lvert V(p,L)-V(q,K)\rvert\le\frac{\varepsilon}{1-\varepsilon}\Delta u$; a boundary response therefore brackets
the value difference instead of measuring it.
\end{enumerate}
\end{proposition}

\begin{proof}
(i) The material part of the compound is $(1-\varepsilon)\E_{qK}[u(O)]+\varepsilon w(\beta)$.  For the trace part,
the coin is action-independent, so both branch posteriors equal $q$ and the chain rule for the compound gives an
information term of $(1-\varepsilon)\I_{qK}(A;O,R)$, the side-bet branch contributing zero.
(ii) By the moreover clause of Theorem~\ref{thm:main}, the indifference holds if and only if
$(1-\varepsilon)V(q,K)+\varepsilon w(\beta)=(1-\varepsilon)V(p,L)+\varepsilon w(\beta')$; rearrange, and
substitute the representation when the material parts cancel.
(iii) $w(\beta)-w(\beta')$ is continuous on $[0,1]^2$ with range $[-\Delta u,\Delta u]$; apply the intermediate
value theorem to the equation in \textup{(ii)}.
\end{proof}

\begin{lemma}[The information--payoff plane]\label{lem:info-payoff-plane}
Write $\phi(q,K):=\bigl(\I_{qK}(A;O,R),\,E_{qK}\bigr)$, and let $\underline u:=\min_ou(o)$ and
$\overline u:=\max_ou(o)$.  The image of $\phi$ over all finite pairs $(q,K)$ is exactly
$[0,\infty)\times[\underline u,\overline u]$, and the indifference classes of the block-comparison relation are,
in these coordinates, the nonempty intersections of the lines $E=c-\lambda\I$ with that strip.
\end{lemma}

\begin{proof}
Inclusion is immediate.  Conversely, fix $(t,w)$ in the strip, choose $n$ with $\log n\ge t$, and on
$A=\{1,\dots,n\}$ use continuity of entropy along a segment from a point mass to the uniform lottery to find
$q_t$ with $\Sh(q_t)=t$.  Let $K(o^+,a\mid a)=\beta$ and $K(o^-,a\mid a)=1-\beta$ with $\beta$ chosen so that
$w(\beta)=w$.  The payoff lottery is action-independent with mean $w$, and the record reveals the action, so
$\phi(q_t,K)=(t,w)$.  The description of the indifference classes then reads the representation through $\phi$,
using the moreover clause of Theorem~\ref{thm:main}.
\end{proof}

\begin{proof}[Proof of Proposition~\ref{prop:single-price}]
The identity rearranges $E_{qK}+\lambda\I_{qK}=E_{pL}+\lambda\I_{pL}$.  For uniqueness, let
$\lambda'\ne\lambda$, put $t:=\tfrac12\min\{1,\,(\overline u-\underline u)/\lambda\}>0$, and use
Lemma~\ref{lem:info-payoff-plane} to realise pairs with $\phi$-coordinates $(t,\overline u-\lambda t)$ and
$(0,\overline u)$; they are indifferent, while $E+\lambda'\I$ differs across them by $(\lambda-\lambda')t\ne0$.
\end{proof}

To compare agents, let both families be defined over the same outcome set $O$, on the common domain of all finite
pairs, with mutual information in one fixed base.  A
channel $B$ is \emph{silent} if all its rows are equal; then the joint law under any $p$ is a product, so
$\I_{pB}=0$ and $V^i(p,B)$ is the expectation of $u_i$ under the $O$-marginal of the common row.  Since
one-action channels realise every payoff lottery in $\D(O)$, silent pairs realise exactly the payoff lotteries.

\begin{definition}[More trace-seeking]
Family 2 is \emph{more trace-seeking than} family 1 if
\[
 (q,K)\succeq^1(p,B)\quad\Longrightarrow\quad(q,K)\succeq^2(p,B)
\]
for all lotteries $q,p$, all channels $K$, and all silent channels $B$.
\end{definition}

\begin{proposition}[More trace-seeking]\label{prop:comparative}
Let $\{\succeq^1_K\}$ and $\{\succeq^2_K\}$ satisfy Axioms~\textup{(A1)--(A8)}.  Then family 2 is more
trace-seeking than family 1 if and only if the representations can be chosen with a common utility and
$\lambda_2\ge\lambda_1$.
\end{proposition}

\begin{proof}[Proof of Proposition~\ref{prop:comparative}]
Throughout, block comparisons are read through the moreover clause of Theorem~\ref{thm:main}.

Suppose the representations are $(u,\lambda_1)$ and $(u,\lambda_2)$ with $\lambda_2\ge\lambda_1$.  For any
$(q,K)$ and silent $(p,B)$,
\[
 V^2(q,K)-V^2(p,B)
 =\bigl[V^1(q,K)-V^1(p,B)\bigr]+(\lambda_2-\lambda_1)\I_{qK}
 \ge V^1(q,K)-V^1(p,B),
\]
since $\I_{qK}\ge0$ and the silent pair's value is material under both.  A weak (or strict) agent-1 preference
therefore transfers to agent 2, and reading the display from right to left gives the mirrored clause
$(p,B)\succeq^2(q,K)\Rightarrow(p,B)\succeq^1(q,K)$: the weak one-sided definition, its strict variant, and the
two-sided variant coincide on this class.

Conversely, suppose $\{\succeq^2_K\}$ is more trace-seeking than $\{\succeq^1_K\}$, with representations
$(u_1,\lambda_1)$ and $(u_2,\lambda_2)$.

\emph{Step 1: common utility.}  Let $D:=\{\mu-\nu:\mu,\nu\in\D(O)\}$ and $f_i(d):=\sum_ou_i(o)d(o)$; each
$f_i\ne0$ on $D$ because $u_i$ is nonconstant.  Silent pairs realise every payoff lottery, and the definition's
left argument ranges over all pairs, so silent-against-silent instances are in its scope:
$f_1(d)\ge0\Rightarrow f_2(d)\ge0$ on $D$.  If $f_1(d)=0$, applying this to $d$ and $-d$ gives $f_2(d)=0$, so
$\ker f_1\subseteq\ker f_2$ and $f_2=\alpha f_1$ for some $\alpha\in\mathbb R$.  If $\alpha=0$ then $u_2$ is
constant, contradicting Axiom~\textup{(A3)} for family 2.  (The case is not idle: an index with constant $u_2$,
which violates only that axiom, is indifferent across all silent pairs and is more trace-seeking than every
family whatsoever.)  If $\alpha<0$, pick $o^+,o^-$ with $u_1(o^+)>u_1(o^-)$: the silent instance
$\delta_{o^+}\succ^1\delta_{o^-}$ forces $\delta_{o^+}\succeq^2\delta_{o^-}$, while $\alpha<0$ gives the strict
reverse.  Hence $\alpha>0$ and $u_2=\alpha u_1+\beta$.  Since $((u_2-\beta)/\alpha,\lambda_2/\alpha)$ again
represents family 2, choose $u_2=u_1=:u$, replacing $\lambda_2$ by $\lambda_2/\alpha>0$.  The choice is
innocuous: any two common-utility normalisations rescale both weights by one positive factor, because
$u$ nonconstant forces the factors to agree, so the order of the weights does not depend on it.

\emph{Step 2: weight ordering.}  Write $\overline u:=\max_ou(o)$ and $\underline u:=\min_ou(o)$;
$\overline u>\underline u$ by Axiom~\textup{(A3)}.  Suppose $\lambda_2<\lambda_1$ and fix
$\lambda'\in(\lambda_2,\lambda_1)$.  Choose $t\in\bigl(0,(\overline u-\underline u)/\lambda'\bigr)$ and
$w:=\underline u+\lambda't\in(\underline u,\overline u)$.  By Lemma~\ref{lem:info-payoff-plane} there is a pair
$(q,K)$ with $\I_{qK}=t$ and $E_{qK}=\underline u$; its payoff lottery is action-independent, so after Step~1
both families assign it the same coordinates.  Choose a silent $(p,B)$ whose common-row $O$-marginal pays
$o^+$ with probability $(w-\underline u)/(\overline u-\underline u)$ and $o^-$ otherwise, so
$V^1(p,B)=V^2(p,B)=w$.  Then $V^1(q,K)=\underline u+\lambda_1t>w$, so $(q,K)\succ^1(p,B)$, and the definition
gives $(q,K)\succeq^2(p,B)$; but $V^2(q,K)=\underline u+\lambda_2t<w$, a contradiction.  Hence
$\lambda_2\ge\lambda_1$.
\end{proof}

\begin{remark}
The proof uses $\lambda_i>0$ only in the choice $E_{qK}=\underline u$; taking $E_{qK}=\overline u$ when
$\lambda'<0$ runs both directions verbatim.  The equivalence therefore holds across all three orientations of
Corollary~\ref{cor:signduality}: more trace-seeking and less trace-averse are one comparative.
\end{remark}

\begin{remark}[Eliciting $\lambda$]\label{rem:elicitation}
Normalise $u(o^+)=1$ and $u(o^-)=0$.  All comparisons below are block comparisons, hence within the primitive
domain, and no information quantity is estimated: the information gap is an accounting identity of the design.
\emph{Step 1.}  For each $o\in O$, elicit $\beta_o$ with $(\delta_\ast,K_o)\sim(\delta_\ast,B_{\beta_o})$, where
$K_o$ is the degenerate one-action channel used in Lemma~\ref{lem:relevance-bridge}; one-action channels carry zero information, so
$u(o)=\beta_o$.
\emph{Step 2.}  Fix $o\in O$, $n\ge2$, the uniform $q$ on $\{1,\dots,n\}$, a side-bet probability $\varepsilon\in(0,1)$, and
the constant-payoff channels $K^{\mathrm{id}}_o$ and $K^0_o$ used in Lemma~\ref{lem:relevance-bridge}.  Elicit $\beta$ with
\[
 \bigl(q,\,D_{\varepsilon,\beta}K^{0}_o\bigr)\sim\bigl(q,\,D_{\varepsilon,0}K^{\mathrm{id}}_o\bigr);
 \qquad\text{then}\qquad
 \lambda=\frac{\varepsilon}{1-\varepsilon}\cdot\frac{\beta}{\log n},
\]
by Proposition~\ref{prop:wtp}\textup{(i)} applied to $V(q,K^{\mathrm{id}}_o)=u(o)+\lambda\log n$ and
$V(q,K^{0}_o)=u(o)$.  Such a $\beta$ exists if and only if $\lambda\log n\le\varepsilon/(1-\varepsilon)$; a
subject still strictly preferring the right block at $\beta=1$ reveals
$\lambda>\varepsilon/\bigl((1-\varepsilon)\log n\bigr)$, and $\varepsilon$ is raised.
\emph{Controls.}  Orientation: Theorem~\ref{thm:main} implies
$(q,K^{\mathrm{id}}_o)\succ(q,K^{0}_o)$ at every payoff, although Axiom~\textup{(A4)} assumes the comparison at only one.  Garbling:
under the total record garbling of $K^{\mathrm{id}}_o$, both blocks carry zero information and indifference must
obtain; a surviving strict preference indicates that an uncontrolled feature of the display still traces the
action.  Transport: repeating Step 2 at another $n$, prior, or payoff level must return the same $\lambda$, by
Proposition~\ref{prop:single-price}; an estimate that declines in $\log n$ is the signature of the
capped-information criterion.
\end{remark}


\clearpage
\begin{thebibliography}{99}
\raggedright

\bibitem{AczelDaroczy1975}
J\'anos Acz\'el and Zolt\'an Dar\'oczy.
\newblock {\em On Measures of Information and Their Characterizations}.
\newblock Academic Press, New York, 1975.

\bibitem{AgranovOrtoleva2017}
Marina Agranov and Pietro Ortoleva.
\newblock Stochastic choice and preferences for randomization.
\newblock {\em Journal of Political Economy}, 125(1):40--68, 2017.

\bibitem{AndreoniBernheim2009}
James Andreoni and B. Douglas Bernheim.
\newblock Social image and the 50--50 norm: A theoretical and experimental analysis of audience effects.
\newblock {\em Econometrica}, 77(5):1607--1636, 2009.

\bibitem{AzrieliLehrer2008}
Yaron Azrieli and Ehud Lehrer.
\newblock The value of a stochastic information structure.
\newblock {\em Games and Economic Behavior}, 63(2):679--693, 2008.

\bibitem{BaezFritz2014}
John C. Baez and Tobias Fritz.
\newblock A {B}ayesian characterization of relative entropy.
\newblock {\em Theory and Applications of Categories}, 29(16):421--456, 2014.

\bibitem{BartlingFehrHerz2014}
Bj\"orn Bartling, Ernst Fehr, and Holger Herz.
\newblock The intrinsic value of decision rights.
\newblock {\em Econometrica}, 82(6):2005--2039, 2014.

\bibitem{BaezFritzLeinster2011}
John C. Baez, Tobias Fritz, and Tom Leinster.
\newblock A characterization of entropy in terms of information loss.
\newblock {\em Entropy}, 13(11):1945--1957, 2011.

\bibitem{BenabouTirole2002}
Roland B\'enabou and Jean Tirole.
\newblock Self-confidence and personal motivation.
\newblock {\em Quarterly Journal of Economics}, 117(3):871--915, 2002.

\bibitem{BenabouTirole2006}
Roland B\'enabou and Jean Tirole.
\newblock Incentives and prosocial behavior.
\newblock {\em American Economic Review}, 96(5):1652--1678, 2006.

\bibitem{Bernheim1994}
B. Douglas Bernheim.
\newblock A theory of conformity.
\newblock {\em Journal of Political Economy}, 102(5):841--877, 1994.

\bibitem{Blackwell1953}
David Blackwell.
\newblock Equivalent comparisons of experiments.
\newblock {\em Annals of Mathematical Statistics}, 24(2):265--272, 1953.

% ORPHANED while evidence paragraphs are parked (see end of Section 5):
%% \bibitem{BuffatPraxmarerSutter2023}
%% Justin Buffat, Matthias Praxmarer, and Matthias Sutter.
%% \newblock The intrinsic value of decision rights: A replication and an extension to team decision making.
%% \newblock {\em Journal of Economic Behavior \& Organization}, 209:560--571, 2023.

\bibitem{CabralesGossnerSerrano2013}
Antonio Cabrales, Olivier Gossner, and Roberto Serrano.
\newblock Entropy and the value of information for investors.
\newblock {\em American Economic Review}, 103(1):360--377, 2013.

\bibitem{CaplinDeanLeahy2022}
Andrew Caplin, Mark Dean, and John Leahy.
\newblock Rationally inattentive behavior: Characterizing and generalizing {S}hannon entropy.
\newblock {\em Journal of Political Economy}, 130(6):1676--1715, 2022.

\bibitem{CaplinLeahy2001}
Andrew Caplin and John Leahy.
\newblock Psychological expected utility theory and anticipatory feelings.
\newblock {\em Quarterly Journal of Economics}, 116(1):55--79, 2001.

\bibitem{CerreiaVioglioDillenbergerOrtolevaRiella2019}
Simone Cerreia-Vioglio, David Dillenberger, Pietro Ortoleva, and Gil Riella.
\newblock Deliberately stochastic.
\newblock {\em American Economic Review}, 109(7):2425--2445, 2019.

\bibitem{CoverThomas2006}
Thomas M. Cover and Joy A. Thomas.
\newblock {\em Elements of Information Theory}.
\newblock Wiley-Interscience, Hoboken, second edition, 2006.

\bibitem{Csiszar2008}
Imre Csisz\'ar.
\newblock Axiomatic characterizations of information measures.
\newblock {\em Entropy}, 10(3):261--273, 2008.

\bibitem{DekelLipmanRustichini2001}
Eddie Dekel, Barton L. Lipman, and Aldo Rustichini.
\newblock Representing preferences with a unique subjective state space.
\newblock {\em Econometrica}, 69(4):891--934, 2001.

\bibitem{Denti2022}
Tommaso Denti.
\newblock Posterior separable cost of information.
\newblock {\em American Economic Review}, 112(10):3215--3259, 2022.

\bibitem{Dostoevsky1864}
Fyodor Dostoevsky.
\newblock {\em Notes from the Underground}, Part I.
\newblock 1864.  Translated by Constance Garnett, 1918.

\bibitem{ElyFrankelKamenica2015}
Jeffrey Ely, Alexander Frankel, and Emir Kamenica.
\newblock Suspense and surprise.
\newblock {\em Journal of Political Economy}, 123(1):215--260, 2015.

\bibitem{Faddeev1956}
D. K. Faddeev.
\newblock On the concept of entropy of a finite probabilistic scheme.
\newblock {\em Uspekhi Matematicheskikh Nauk}, 11(1):227--231, 1956.
\newblock In Russian.

\bibitem{FrankelKamenica2019}
Alexander Frankel and Emir Kamenica.
\newblock Quantifying information and uncertainty.
\newblock {\em American Economic Review}, 109(10):3650--3680, 2019.

\bibitem{GhirardatoMarinacci2002}
Paolo Ghirardato and Massimo Marinacci.
\newblock Ambiguity made precise: A comparative foundation.
\newblock {\em Journal of Economic Theory}, 102(2):251--289, 2002.

\bibitem{GilboaLehrer1991}
Itzhak Gilboa and Ehud Lehrer.
\newblock The value of information---an axiomatic approach.
\newblock {\em Journal of Mathematical Economics}, 20(5):443--459, 1991.

\bibitem{GlazerKonrad1996}
Amihai Glazer and Kai A. Konrad.
\newblock A signaling explanation for charity.
\newblock {\em American Economic Review}, 86(4):1019--1028, 1996.

\bibitem{GrantKajiiPolak1998}
Simon Grant, Atsushi Kajii, and Ben Polak.
\newblock Intrinsic preference for information.
\newblock {\em Journal of Economic Theory}, 83(2):233--259, 1998.

\bibitem{GulPesendorfer2001}
Faruk Gul and Wolfgang Pesendorfer.
\newblock Temptation and self-control.
\newblock {\em Econometrica}, 69(6):1403--1435, 2001.

\bibitem{HersteinMilnor1953}
I. N. Herstein and John Milnor.
\newblock An axiomatic approach to measurable utility.
\newblock {\em Econometrica}, 21(2):291--297, 1953.

% ORPHANED while Hobson is uncited (cut from Section 5, 2026-08-12):
% \bibitem{Hobson1969}
% Arthur Hobson.
% \newblock A new theorem of information theory.
% \newblock {\em Journal of Statistical Physics}, 1(3):383--391, 1969.

\bibitem{Khinchin1957}
Aleksandr I. Khinchin.
\newblock {\em Mathematical Foundations of Information Theory}.
\newblock Dover, New York, 1957.

\bibitem{KlyubinPolaniNehaniv2005}
Alexander S. Klyubin, Daniel Polani, and Chrystopher L. Nehaniv.
\newblock Empowerment: A universal agent-centric measure of control.
\newblock In {\em Proceedings of the 2005 IEEE Congress on Evolutionary Computation}, volume~1, pages 128--135. IEEE,
  2005.

\bibitem{KornerMarton1977}
J\'anos K\"orner and Katalin Marton.
\newblock Comparison of two noisy channels.
\newblock In Imre Csisz\'ar and Peter Elias, editors, {\em Topics in Information Theory}, Colloquia Mathematica
  Societatis J\'anos Bolyai, volume~16, pages 411--423. North-Holland, Amsterdam, 1977.

\bibitem{Kreps1979}
David M. Kreps.
\newblock A representation theorem for preference for flexibility.
\newblock {\em Econometrica}, 47(3):565--577, 1979.

\bibitem{KrepsPorteus1978}
David M. Kreps and Evan L. Porteus.
\newblock Temporal resolution of uncertainty and dynamic choice theory.
\newblock {\em Econometrica}, 46(1):185--200, 1978.

% ORPHANED while evidence paragraphs are parked (see end of Section 5):
%% \bibitem{Liljeholm2022}
%% Mimi Liljeholm.
%% \newblock Flexible control as surrogate reward or dynamic reward maximization.
%% \newblock {\em Cognition}, 229:105262, 2022.

\bibitem{MatejkaMcKay2015}
Filip Mat\v{e}jka and Alisdair McKay.
\newblock Rational inattention to discrete choices: A new foundation for the multinomial logit model.
\newblock {\em American Economic Review}, 105(1):272--298, 2015.

\bibitem{MistryLiljeholm2016}
Prachi Mistry and Mimi Liljeholm.
\newblock Instrumental divergence and the value of control.
\newblock {\em Scientific Reports}, 6:36295, 2016.

\bibitem{NortonLiljeholm2020}
Kaitlyn G. Norton and Mimi Liljeholm.
\newblock The rostrolateral prefrontal cortex mediates a preference for high-agency environments.
\newblock {\em Journal of Neuroscience}, 40(22):4401--4409, 2020.

% ORPHANED while evidence paragraphs are parked (see end of Section 5):
%% \bibitem{OwensGrossmanFackler2014}
%% David Owens, Zachary Grossman, and Ryan Fackler.
%% \newblock The control premium: A preference for payoff autonomy.
%% \newblock {\em American Economic Journal: Microeconomics}, 6(4):138--161, 2014.

\bibitem{PattanaikXu1990}
Prasanta K. Pattanaik and Yongsheng Xu.
\newblock On ranking opportunity sets in terms of freedom of choice.
\newblock {\em Recherches \'Economiques de Louvain}, 56(3--4):383--390, 1990.

\bibitem{PomattoStrackTamuz2023}
Luciano Pomatto, Philipp Strack, and Omer Tamuz.
\newblock The cost of information: The case of constant marginal costs.
\newblock {\em American Economic Review}, 113(5):1360--1393, 2023.

\bibitem{Renyi1961}
Alfr\'ed R\'enyi.
\newblock On measures of entropy and information.
\newblock In {\em Proceedings of the Fourth Berkeley Symposium on Mathematical Statistics and Probability}, volume~1,
  pages 547--561. University of California Press, 1961.

\bibitem{SalgeGlackinPolani2014}
Christoph Salge, Cornelius Glackin, and Daniel Polani.
\newblock Empowerment---an introduction.
\newblock In Mikhail Prokopenko, editor, {\em Guided Self-Organization: Inception}, pages 67--114. Springer, Berlin,
  2014.

\bibitem{Savage1954}
Leonard J. Savage.
\newblock {\em The Foundations of Statistics}.
\newblock Wiley, New York, 1954.

\bibitem{Shannon1948}
Claude E. Shannon.
\newblock A mathematical theory of communication.
\newblock {\em Bell System Technical Journal}, 27:379--423, 623--656, 1948.

\bibitem{Sims2003}
Christopher A. Sims.
\newblock Implications of rational inattention.
\newblock {\em Journal of Monetary Economics}, 50(3):665--690, 2003.

\bibitem{Spence1973}
Michael Spence.
\newblock Job market signaling.
\newblock {\em Quarterly Journal of Economics}, 87(3):355--374, 1973.

\bibitem{WynerZiv1973}
Aaron D. Wyner and Jacob Ziv.
\newblock A theorem on the entropy of certain binary sequences and applications: Part {I}.
\newblock {\em IEEE Transactions on Information Theory}, 19(6):769--772, 1973.

\end{thebibliography}



\end{document}
